\section{Introduction}
\label{topologies-section-introduction}
In this document we explain what the different topologies on the
category of schemes are. Some references are \cite{SGA1} and \cite{Ner}.
Before doing so we would like to point out that there are many
different choices of sites (as defined in
Sites, Definition \ref{sites-definition-site}) which give rise to
the same notion of sheaf on the underlying category. Hence
our choices may be slightly different from those in the references
but ultimately lead to the same cohomology groups, etc.

\section{The general procedure}
\label{topologies-section-procedure}
In this section we explain a general procedure for producing the
sites we will be working with. Suppose we want to study sheaves
over schemes with respect to some topology $\tau$. In order to
get a site, as in Sites, Definition \ref{sites-definition-site},
of schemes with that topology we have to do some work. Namely,
we cannot simply say ``consider all schemes with the Zariski topology''
since that would give a ``big'' category. Instead, in each section of
this chapter we will proceed as follows:
\begin{enumerate}
\item We define a class $\text{Cov}_\tau$ of coverings of schemes
satisfying the axioms of Sites, Definition \ref{sites-definition-site}.
It will always be the case that a Zariski open covering of
a scheme is a covering for $\tau$.
\item We single out a notion of standard
$\tau$-covering within the category of affine schemes.
\item We define what is an ``absolute'' big $\tau$-site $\Sch_\tau$.
These are the sites one gets by appropriately choosing a set of schemes
and a set of coverings.
\item For any object $S$ of $\Sch_\tau$
we define the big $\tau$-site $(\Sch/S)_\tau$ and for suitable
$\tau$ the small\footnote{The words big and
small here do not relate to bigness/smallness of the corresponding
categories.} $\tau$-site $S_\tau$.
\item In addition there is a site $(\textit{Aff}/S)_\tau$ using the
notion of standard $\tau$-covering of affines\footnote{In the case
of the ph topology we deviate very slightly from this approach, see
Definition \ref{topologies-definition-big-small-ph} and the surrounding discussion.}
whose category of sheaves
is equivalent to the category of sheaves on $(\Sch/S)_\tau$.
\end{enumerate}
The above is a little clumsy in that we do not end up with a canonical
choice for the big $\tau$-site of a scheme, or even the small
$\tau$-site of a scheme. If you are willing to ignore set theoretic
difficulties, then you can work with classes and end up with
canonical big and small sites...

\section{The Zariski topology}
\label{topologies-section-zariski}
\begin{definition}[Zariski covering]
\label{topologies-definition-zariski-covering}
\source{stacks:020O}
\href{https://stacks.math.columbia.edu/tag/020O}{\texttt{Stacks tag 020O}}
\group{topologies}

Let $T$ be a scheme. A {\it Zariski covering of $T$} is a family
of morphisms $\{f_i : T_i \to T\}_{i \in I}$ of schemes
such that each $f_i$ is an open immersion and such
that $T = \bigcup f_i(T_i)$.
\end{definition}

This defines a (proper) class of coverings.
Next, we show that this notion satisfies the conditions of
Sites, Definition \ref{sites-definition-site}.

\begin{lemma}[Zariski]
\label{topologies-lemma-zariski}
\source{stacks:020P}
\href{https://stacks.math.columbia.edu/tag/020P}{\texttt{Stacks tag 020P}}
\group{topologies}

Let $T$ be a scheme.
\begin{enumerate}
\item If $T' \to T$ is an isomorphism then $\{T' \to T\}$
is a Zariski covering of $T$.
\item If $\{T_i \to T\}_{i\in I}$ is a Zariski covering and for each
$i$ we have a Zariski covering $\{T_{ij} \to T_i\}_{j\in J_i}$, then
$\{T_{ij} \to T\}_{i \in I, j\in J_i}$ is a Zariski covering.
\item If $\{T_i \to T\}_{i\in I}$ is a Zariski covering
and $T' \to T$ is a morphism of schemes then
$\{T' \times_T T_i \to T'\}_{i\in I}$ is a Zariski covering.
\end{enumerate}
\end{lemma}

\begin{proof}
Omitted.
\end{proof}

\begin{lemma}[Zariski affine]
\label{topologies-lemma-zariski-affine}
\source{stacks:020Q}
\href{https://stacks.math.columbia.edu/tag/020Q}{\texttt{Stacks tag 020Q}}
\group{topologies}

Let $T$ be an affine scheme. Let $\{T_i \to T\}_{i \in I}$ be a
Zariski covering of $T$. Then there exists a Zariski covering
$\{U_j \to T\}_{j = 1, \ldots, m}$ which is a refinement
of $\{T_i \to T\}_{i \in I}$ such that each $U_j$ is a standard
open of $T$, see
Schemes, Definition \ref{schemes-definition-standard-covering}.
Moreover, we may choose each $U_j$ to be an open of one of the $T_i$.
\end{lemma}

\begin{proof}
Follows as $T$ is quasi-compact and standard opens form a basis
for its topology. This is also proved in
Schemes, Lemma \ref{schemes-lemma-standard-open}.
\end{proof}

Thus we define the corresponding standard coverings of affines as follows.

\begin{definition}[Standard Zariski]
\label{topologies-definition-standard-Zariski}
\source{stacks:020R}
\href{https://stacks.math.columbia.edu/tag/020R}{\texttt{Stacks tag 020R}}
\group{topologies}

Compare Schemes, Definition \ref{schemes-definition-standard-covering}.
Let $T$ be an affine scheme. A {\it standard Zariski covering}
of $T$ is a Zariski covering $\{U_j \to T\}_{j = 1, \ldots, m}$
with each $U_j \to T$ inducing an isomorphism with a standard affine open
of $T$.
\end{definition}

\begin{definition}[Big zariski site]
\label{topologies-definition-big-zariski-site}
\source{stacks:020S}
\href{https://stacks.math.columbia.edu/tag/020S}{\texttt{Stacks tag 020S}}
\group{topologies}

A {\it big Zariski site} is any site $\Sch_{Zar}$ as in
Sites, Definition \ref{sites-definition-site} constructed as follows:
\begin{enumerate}
\item Choose any set of schemes $S_0$, and any set of Zariski coverings
$\text{Cov}_0$ among these schemes.
\item As underlying category of $\Sch_{Zar}$
take any category $\Sch_\alpha$ constructed as in
Sets, Lemma \ref{sets-lemma-construct-category} starting with the set $S_0$.
\item As coverings of $\Sch_{Zar}$ choose any set of coverings as in
Sets, Lemma \ref{sets-lemma-coverings-site} starting with the
category $\Sch_\alpha$ and the class of Zariski coverings,
and the set $\text{Cov}_0$ chosen above.
\end{enumerate}
\end{definition}

It is shown in Sites, Lemma \ref{sites-lemma-choice-set-coverings-immaterial}
that, after having chosen the category $\Sch_\alpha$, the
category of sheaves on $\Sch_\alpha$ does not depend on the
choice of coverings chosen in (3) above. In other words, the topos
$\Sh(\Sch_{Zar})$ only depends on the choice of
the category $\Sch_\alpha$. It is shown in
Sets, Lemma \ref{sets-lemma-what-is-in-it} that these categories
are closed under many constructions of algebraic geometry, e.g.,
fibre products and taking open and closed subschemes. We can also show
that the exact choice of $\Sch_\alpha$ does not matter
too much, see Section \ref{topologies-section-change-alpha}.

\medskip\noindent
Another approach would be to assume the existence of a
strongly inaccessible cardinal and to define $\Sch_{Zar}$
to be the category of schemes contained in a chosen universe with
set of coverings the Zariski coverings contained in that same
universe.

\medskip\noindent
Before we continue with the introduction of the big Zariski site of
a scheme $S$, let us point out that the topology on a big Zariski site
$\Sch_{Zar}$ is in some sense induced from the Zariski topology
on the category of all schemes.

\begin{lemma}[Zariski induced]
\label{topologies-lemma-zariski-induced}
\source{stacks:03WV}
\href{https://stacks.math.columbia.edu/tag/03WV}{\texttt{Stacks tag 03WV}}
\group{topologies}

\uses{topologies-definition-big-zariski-site}
Let $\Sch_{Zar}$ be a big Zariski site as in
Definition \ref{topologies-definition-big-zariski-site}.
Let $T \in \Ob(\Sch_{Zar})$.
Let $\{T_i \to T\}_{i \in I}$ be an arbitrary Zariski covering of $T$.
There exists a covering $\{U_j \to T\}_{j \in J}$ of $T$ in the site
$\Sch_{Zar}$ which is tautologically equivalent (see
Sites, Definition \ref{sites-definition-combinatorial-tautological})
to $\{T_i \to T\}_{i \in I}$.
\end{lemma}

\begin{proof}
Since each $T_i \to T$ is an open immersion, we see by
Sets, Lemma \ref{sets-lemma-what-is-in-it}
that each $T_i$ is isomorphic to an object $V_i$ of $\Sch_{Zar}$.
The covering $\{V_i \to T\}_{i \in I}$ is tautologically equivalent
to $\{T_i \to T\}_{i \in I}$ (using the identity map on $I$ both ways).
Moreover, $\{V_i \to T\}_{i \in I}$ is combinatorially equivalent to a
covering $\{U_j \to T\}_{j \in J}$ of $T$ in the site $\Sch_{Zar}$ by
Sets, Lemma \ref{sets-lemma-coverings-site}.
\end{proof}

\begin{definition}[Big small Zariski]
\label{topologies-definition-big-small-Zariski}
\source{stacks:020T}
\href{https://stacks.math.columbia.edu/tag/020T}{\texttt{Stacks tag 020T}}
\group{topologies}

Let $S$ be a scheme. Let $\Sch_{Zar}$ be a big Zariski
site containing $S$.
\begin{enumerate}
\item The {\it big Zariski site of $S$}, denoted
$(\Sch/S)_{Zar}$, is the site $\Sch_{Zar}/S$
introduced in Sites, Section \ref{sites-section-localize}.
\item The {\it small Zariski site of $S$}, which we denote
$S_{Zar}$, is the full subcategory of $(\Sch/S)_{Zar}$
whose objects are those $U/S$ such that $U \to S$ is an open immersion.
A covering of $S_{Zar}$ is any covering $\{U_i \to U\}$ of
$(\Sch/S)_{Zar}$ with $U \in \Ob(S_{Zar})$.
\item The {\it big affine Zariski site of $S$}, denoted
$(\textit{Aff}/S)_{Zar}$, is the full subcategory of
$(\Sch/S)_{Zar}$ consisting of objects $U/S$ such that $U$ is an
affine scheme. A covering of $(\textit{Aff}/S)_{Zar}$ is any covering
$\{U_i \to U\}$ of $(\Sch/S)_{Zar}$ with $U \in \Ob((\textit{Aff}/S)_{Zar})$
which is a standard Zariski covering.
\item The {\it small affine Zariski site of $S$}, denoted
$S_{affine, Zar}$, is the full subcategory of $S_{Zar}$
whose objects are those $U/S$ such that $U$ is an affine scheme.
A covering of $S_{affine, Zar}$ is any covering $\{U_i \to U\}$ of
$S_{Zar}$ with $U \in \Ob(S_{affine, Zar})$
which is a standard Zariski covering.
\end{enumerate}
\end{definition}

It is not completely clear that the small Zariski site,
the big affine Zariski site, and the small affine Zariski site
are sites. We check this now.

\begin{lemma}[Verify site Zariski]
\label{topologies-lemma-verify-site-Zariski}
\source{stacks:020U}
\href{https://stacks.math.columbia.edu/tag/020U}{\texttt{Stacks tag 020U}}
\group{topologies}

Let $S$ be a scheme. Let $\Sch_{Zar}$ be a big Zariski
site containing $S$. The structures $S_{Zar}$, $(\textit{Aff}/S)_{Zar}$,
and $S_{affine, Zar}$ defined above are sites.
\end{lemma}

\begin{proof}
Let us show that $S_{Zar}$ is a site. It is a category with a
given set of families of morphisms with fixed target. Thus we
have to show properties (1), (2) and (3) of
Sites, Definition \ref{sites-definition-site}.
Since $(\Sch/S)_{Zar}$ is a site, it suffices to prove
that given any covering $\{U_i \to U\}$ of $(\Sch/S)_{Zar}$
with $U \in \Ob(S_{Zar})$ we also have $U_i \in \Ob(S_{Zar})$.
This follows from the definitions
as the composition of open immersions is an open immersion.

\medskip\noindent
Let us show that $(\textit{Aff}/S)_{Zar}$ is a site.
Reasoning as above, it suffices to show that the collection
of standard Zariski coverings of affines satisfies properties
(1), (2) and (3) of
Sites, Definition \ref{sites-definition-site}.
Let $R$ be a ring. Let $f_1, \ldots, f_n \in R$ generate the unit ideal.
For each $i \in \{1, \ldots, n\}$ let $g_{i1}, \ldots, g_{in_i} \in R_{f_i}$
be elements generating the unit ideal of $R_{f_i}$. Write
$g_{ij} = f_{ij}/f_i^{e_{ij}}$ which is possible. After replacing
$f_{ij}$ by $f_i f_{ij}$ if necessary, we have that
$D(f_{ij}) \subset D(f_i) \cong \Spec(R_{f_i})$ is
equal to $D(g_{ij}) \subset \Spec(R_{f_i})$. Hence we see that
the family of morphisms $\{D(g_{ij}) \to \Spec(R)\}$
is a standard Zariski covering. From these considerations
it follows that (2) holds for standard Zariski coverings.
We omit the verification of (1) and (3).

\medskip\noindent
We omit the proof that $S_{affine, Zar}$ is a site.
\end{proof}

\begin{lemma}[Fibre products Zariski]
\label{topologies-lemma-fibre-products-Zariski}
\source{stacks:020V}
\href{https://stacks.math.columbia.edu/tag/020V}{\texttt{Stacks tag 020V}}
\group{topologies}

Let $S$ be a scheme. Let $\Sch_{Zar}$ be a big Zariski
site containing $S$. The underlying categories of the sites
$\Sch_{Zar}$, $(\Sch/S)_{Zar}$,
$S_{Zar}$, $(\textit{Aff}/S)_{Zar}$, and $S_{affine, Zar}$ have fibre products.
In each case the obvious functor into the category $\Sch$ of
all schemes commutes with taking fibre products. The categories
$(\Sch/S)_{Zar}$, and $S_{Zar}$ both have a final object,
namely $S/S$.
\end{lemma}

\begin{proof}
For $\Sch_{Zar}$ it is true by construction, see
Sets, Lemma \ref{sets-lemma-what-is-in-it}.
Suppose we have $U \to S$, $V \to U$, $W \to U$ morphisms
of schemes with $U, V, W \in \Ob(\Sch_{Zar})$.
The fibre product $V \times_U W$ in $\Sch_{Zar}$
is a fibre product in $\Sch$ and
is the fibre product of $V/S$ with $W/S$ over $U/S$ in
the category of all schemes over $S$, and hence also a
fibre product in $(\Sch/S)_{Zar}$.
This proves the result for $(\Sch/S)_{Zar}$.
If $U \to S$, $V \to U$ and $W \to U$ are open immersions then so is
$V \times_U W \to S$ and hence we get the result for $S_{Zar}$.
If $U, V, W$ are affine, so is $V \times_U W$ and hence the
result for $(\textit{Aff}/S)_{Zar}$ and $S_{affine, Zar}$.
\end{proof}

Next, we check that the big, resp.\ small affine site defines the same
topos as the big, resp.\ small site.

\begin{lemma}[Affine big site Zariski]
\label{topologies-lemma-affine-big-site-Zariski}
\source{stacks:020W}
\href{https://stacks.math.columbia.edu/tag/020W}{\texttt{Stacks tag 020W}}
\group{topologies}

Let $S$ be a scheme. Let $\Sch_{Zar}$ be a big Zariski
site containing $S$.
The functor $(\textit{Aff}/S)_{Zar} \to (\Sch/S)_{Zar}$
is a special cocontinuous functor. Hence it induces an equivalence
of topoi from $\Sh((\textit{Aff}/S)_{Zar})$ to
$\Sh((\Sch/S)_{Zar})$.
\end{lemma}

\begin{proof}
\uses{topologies-lemma-zariski-affine}

The notion of a special cocontinuous functor is introduced in
Sites, Definition \ref{sites-definition-special-cocontinuous-functor}.
Thus we have to verify assumptions (1) -- (5) of
Sites, Lemma \ref{sites-lemma-equivalence}.
Denote the inclusion functor
$u : (\textit{Aff}/S)_{Zar} \to (\Sch/S)_{Zar}$.
Being cocontinuous just means that any Zariski covering of
$T/S$, $T$ affine, can be refined by a standard Zariski covering of $T$.
This is the content of
Lemma \ref{topologies-lemma-zariski-affine}.
Hence (1) holds. We see $u$ is continuous simply because a standard
Zariski covering is a Zariski covering. Hence (2) holds.
Parts (3) and (4) follow immediately from the fact that $u$ is
fully faithful. And finally condition (5) follows from the
fact that every scheme has an affine open covering.
\end{proof}

\begin{lemma}[Alternative zariski]
\label{topologies-lemma-alternative-zariski}
\source{stacks:0F1B}
\href{https://stacks.math.columbia.edu/tag/0F1B}{\texttt{Stacks tag 0F1B}}
\group{topologies}

Let $S$ be a scheme.  Let $\Sch_{Zar}$ be a big Zariski
site containing $S$. The functor $S_{affine, Zar} \to S_{Zar}$
is a special cocontinuous functor. Hence it induces an equivalence
of topoi from $\Sh(S_{affine, Zar})$ to $\Sh(S_{Zar})$.
\end{lemma}

\begin{proof}
\uses{topologies-lemma-affine-big-site-Zariski}

Omitted. Hint: compare with the proof of
Lemma \ref{topologies-lemma-affine-big-site-Zariski}.
\end{proof}

Let us check that the notion of a sheaf on the small Zariski site
corresponds to notion of a sheaf on $S$.

\begin{lemma}[Zariski usual]
\label{topologies-lemma-Zariski-usual}
\source{stacks:020X}
\href{https://stacks.math.columbia.edu/tag/020X}{\texttt{Stacks tag 020X}}
\group{topologies}

The category of sheaves on $S_{Zar}$ is equivalent to the
category of sheaves on the underlying topological space of $S$.
\end{lemma}

\begin{proof}
We will use repeatedly that for any object
$U/S$ of $S_{Zar}$ the morphism $U \to S$ is an isomorphism
onto an open subscheme.
Let $\mathcal{F}$ be a sheaf on $S$. Then we define a sheaf
on $S_{Zar}$ by the rule $\mathcal{F}'(U/S) = \mathcal{F}(\Im(U \to S))$.
For the converse, we choose for every open subscheme $U \subset S$ an object
$U'/S \in \Ob(S_{Zar})$ with $\Im(U' \to S) = U$
(here you have to use Sets, Lemma \ref{sets-lemma-what-is-in-it}).
Given a sheaf $\mathcal{G}$ on $S_{Zar}$ we define a sheaf on $S$ by setting
$\mathcal{G}'(U) = \mathcal{G}(U'/S)$. To see that $\mathcal{G}'$ is
a sheaf we use that for any open covering $U = \bigcup_{i \in I} U_i$
the covering $\{U_i \to U\}_{i \in I}$
is combinatorially equivalent to a covering $\{U_j' \to U'\}_{j \in J}$
in $S_{Zar}$ by Sets, Lemma \ref{sets-lemma-coverings-site},
and we use Sites, Lemma \ref{sites-lemma-tautological-same-sheaf}.
Details omitted.
\end{proof}

From now on we will not make any distinction between a sheaf on
$S_{Zar}$ or a sheaf on $S$. We will always use the procedures
of the proof of the lemma to go between the two notions.
Next, we establish some relationships between the topoi
associated to these sites.

\begin{lemma}[Put in T]
\label{topologies-lemma-put-in-T}
\source{stacks:020Y}
\href{https://stacks.math.columbia.edu/tag/020Y}{\texttt{Stacks tag 020Y}}
\group{topologies}

Let $\Sch_{Zar}$ be a big Zariski site.
Let $f : T \to S$ be a morphism in $\Sch_{Zar}$.
The functor $T_{Zar} \to (\Sch/S)_{Zar}$
is cocontinuous and induces a morphism of topoi
$$
i_f :
\Sh(T_{Zar})
\longrightarrow
\Sh((\Sch/S)_{Zar})
$$
For a sheaf $\mathcal{G}$ on $(\Sch/S)_{Zar}$
we have the formula $(i_f^{-1}\mathcal{G})(U/T) = \mathcal{G}(U/S)$.
The functor $i_f^{-1}$ also has a left adjoint $i_{f, !}$ which commutes
with fibre products and equalizers.
\end{lemma}

\begin{proof}
\uses{topologies-lemma-fibre-products-Zariski}

Denote the functor $u : T_{Zar} \to (\Sch/S)_{Zar}$.
In other words, given and open immersion $j : U \to T$ corresponding
to an object of $T_{Zar}$ we set $u(U \to T) = (f \circ j : U \to S)$.
This functor commutes with fibre products, see
Lemma \ref{topologies-lemma-fibre-products-Zariski}.
Moreover, $T_{Zar}$ has equalizers (as any two morphisms with the same
source and target are the same) and $u$ commutes with them.
It is clearly cocontinuous.
It is also continuous as $u$ transforms coverings to coverings and
commutes with fibre products. Hence the lemma follows from
Sites, Lemmas \ref{sites-lemma-when-shriek}
and \ref{sites-lemma-preserve-equalizers}.
\end{proof}

\begin{lemma}[At the bottom]
\label{topologies-lemma-at-the-bottom}
\source{stacks:020Z}
\href{https://stacks.math.columbia.edu/tag/020Z}{\texttt{Stacks tag 020Z}}
\group{topologies}

\uses{topologies-lemma-put-in-T}
Let $S$ be a scheme. Let $\Sch_{Zar}$ be a big Zariski
site containing $S$.
The inclusion functor $S_{Zar} \to (\Sch/S)_{Zar}$
satisfies the hypotheses of Sites, Lemma \ref{sites-lemma-bigger-site}
and hence induces a morphism of sites
$$
\pi_S : (\Sch/S)_{Zar} \longrightarrow S_{Zar}
$$
and a morphism of topoi
$$
i_S : \Sh(S_{Zar}) \longrightarrow \Sh((\Sch/S)_{Zar})
$$
such that $\pi_S \circ i_S = \text{id}$. Moreover, $i_S = i_{\text{id}_S}$
with $i_{\text{id}_S}$ as in Lemma \ref{topologies-lemma-put-in-T}. In particular the
functor $i_S^{-1} = \pi_{S, *}$ is described by the rule
$i_S^{-1}(\mathcal{G})(U/S) = \mathcal{G}(U/S)$.
\end{lemma}

\begin{proof}
\uses{topologies-lemma-put-in-T}

In this case the functor $u : S_{Zar} \to (\Sch/S)_{Zar}$,
in addition to the properties seen in the proof of
Lemma \ref{topologies-lemma-put-in-T} above, also is fully faithful
and transforms the final object into the final object.
The lemma follows.
\end{proof}

\begin{definition}[Restriction small zariski]
\label{topologies-definition-restriction-small-zariski}
\source{stacks:04BS}
\href{https://stacks.math.columbia.edu/tag/04BS}{\texttt{Stacks tag 04BS}}
\group{topologies}

\uses{topologies-lemma-at-the-bottom}
In the situation of
Lemma \ref{topologies-lemma-at-the-bottom}
the functor $i_S^{-1} = \pi_{S, *}$ is often
called the {\it restriction to the small Zariski site}, and for a sheaf
$\mathcal{F}$ on the big Zariski site we denote $\mathcal{F}|_{S_{Zar}}$
this restriction.
\end{definition}

With this notation in place we have for a sheaf $\mathcal{F}$ on the
big site and a sheaf $\mathcal{G}$ on the small site that
\begin{align*}
\Mor_{\Sh(S_{Zar})}(\mathcal{F}|_{S_{Zar}}, \mathcal{G})
& =
\Mor_{\Sh((\Sch/S)_{Zar})}(\mathcal{F},
i_{S, *}\mathcal{G}) \\
\Mor_{\Sh(S_{Zar})}(\mathcal{G}, \mathcal{F}|_{S_{Zar}})
& =
\Mor_{\Sh((\Sch/S)_{Zar})}(\pi_S^{-1}\mathcal{G},
\mathcal{F})
\end{align*}
Moreover, we have $(i_{S, *}\mathcal{G})|_{S_{Zar}} = \mathcal{G}$
and we have $(\pi_S^{-1}\mathcal{G})|_{S_{Zar}} = \mathcal{G}$.

\begin{lemma}[Morphism big]
\label{topologies-lemma-morphism-big}
\source{stacks:0210}
\href{https://stacks.math.columbia.edu/tag/0210}{\texttt{Stacks tag 0210}}
\group{topologies}

Let $\Sch_{Zar}$ be a big Zariski site.
Let $f : T \to S$ be a morphism in $\Sch_{Zar}$.
The functor
$$
u : (\Sch/T)_{Zar} \longrightarrow (\Sch/S)_{Zar},
\quad
V/T \longmapsto V/S
$$
is cocontinuous, and has a continuous right adjoint
$$
v : (\Sch/S)_{Zar} \longrightarrow (\Sch/T)_{Zar},
\quad
(U \to S) \longmapsto (U \times_S T \to T).
$$
They induce the same morphism of topoi
$$
f_{big} :
\Sh((\Sch/T)_{Zar})
\longrightarrow
\Sh((\Sch/S)_{Zar})
$$
We have $f_{big}^{-1}(\mathcal{G})(U/T) = \mathcal{G}(U/S)$.
We have $f_{big, *}(\mathcal{F})(U/S) = \mathcal{F}(U \times_S T/T)$.
Also, $f_{big}^{-1}$ has a left adjoint $f_{big!}$ which commutes with
fibre products and equalizers.
\end{lemma}

\begin{proof}
\uses{topologies-lemma-put-in-T}

The functor $u$ is cocontinuous, continuous, and commutes with fibre products
and equalizers (details omitted; compare with proof of
Lemma \ref{topologies-lemma-put-in-T}).
Hence
Sites, Lemmas \ref{sites-lemma-when-shriek} and
\ref{sites-lemma-preserve-equalizers}
apply and we deduce the formula
for $f_{big}^{-1}$ and the existence of $f_{big!}$. Moreover,
the functor $v$ is a right adjoint because given $U/T$ and $V/S$
we have $\Mor_S(u(U), V) = \Mor_T(U, V \times_S T)$
as desired. Thus we may apply
Sites, Lemmas \ref{sites-lemma-have-functor-other-way} and
\ref{sites-lemma-have-functor-other-way-morphism}
to get the formula for $f_{big, *}$.
\end{proof}

\begin{lemma}[Morphism big small]
\label{topologies-lemma-morphism-big-small}
\source{stacks:0211}
\href{https://stacks.math.columbia.edu/tag/0211}{\texttt{Stacks tag 0211}}
\group{topologies}

\uses{topologies-lemma-put-in-T, topologies-lemma-at-the-bottom, topologies-lemma-Zariski-usual}
Let $\Sch_{Zar}$ be a big Zariski site.
Let $f : T \to S$ be a morphism in $\Sch_{Zar}$.
\begin{enumerate}
\item We have $i_f = f_{big} \circ i_T$ with $i_f$ as in
Lemma \ref{topologies-lemma-put-in-T} and $i_T$ as in
Lemma \ref{topologies-lemma-at-the-bottom}.
\item The functor $S_{Zar} \to T_{Zar}$,
$(U \to S) \mapsto (U \times_S T \to T)$ is continuous and induces
a morphism of sites
$$
f_{small} :
T_{Zar}
\longrightarrow
S_{Zar}
$$
The functors $f_{small}^{-1}$ and $f_{small, *}$ agree with
the usual notions $f^{-1}$ and $f_*$ if we identify sheaves
on $T_{Zar}$, resp.\ $S_{Zar}$ with sheaves on $T$, resp.\ $S$
via Lemma \ref{topologies-lemma-Zariski-usual}.
\item We have a commutative diagram of morphisms of sites
$$
\begin{array}{cc}T_{Zar} \underset{f_{small}}{\downarrow} & (\Sch/T)_{Zar} \overset{f_{big}}{\downarrow}\overset{\pi_T}{\longleftarrow} \\ S_{Zar} & (\Sch/S)_{Zar} \underset{\pi_S}{\longleftarrow}\end{array}
$$
so that $f_{small} \circ \pi_T = \pi_S \circ f_{big}$ as morphisms of topoi.
\item We have $f_{small} = \pi_S \circ f_{big} \circ i_T = \pi_S \circ i_f$.
\end{enumerate}
\end{lemma}

\begin{proof}
The equality $i_f = f_{big} \circ i_T$ follows from the
equality $i_f^{-1} = i_T^{-1} \circ f_{big}^{-1}$ which is
clear from the descriptions of these functors above.
Thus we see (1).

\medskip\noindent
Statement (2): See Sites, Example \ref{sites-example-continuous-map}.

\medskip\noindent
Part (3) follows because $\pi_S$ and $\pi_T$ are given by
the inclusion functors and $f_{small}$ and $f_{big}$ by the
base change functor $U \mapsto U \times_S T$.

\medskip\noindent
Statement (4) follows from (3) by precomposing with $i_T$.
\end{proof}

In the situation of the lemma, using the terminology of
Definition \ref{topologies-definition-restriction-small-zariski}
we have: for $\mathcal{F}$ a sheaf on the big Zariski site of $T$
$$
(f_{big, *}\mathcal{F})|_{S_{Zar}} =
f_{small, *}(\mathcal{F}|_{T_{Zar}}),
$$
This equality is clear from the commutativity of the diagram of
sites of the lemma, since restriction to the small Zariski site of
$T$, resp.\ $S$ is given by $\pi_{T, *}$, resp.\ $\pi_{S, *}$. A similar
formula involving pullbacks and restrictions is false.

\begin{lemma}[Composition]
\label{topologies-lemma-composition}
\source{stacks:0212}
\href{https://stacks.math.columbia.edu/tag/0212}{\texttt{Stacks tag 0212}}
\group{topologies}

Given schemes $X$, $Y$, $Z$ in $(\Sch/S)_{Zar}$
and morphisms $f : X \to Y$, $g : Y \to Z$ we have
$g_{big} \circ f_{big} = (g \circ f)_{big}$ and
$g_{small} \circ f_{small} = (g \circ f)_{small}$.
\end{lemma}

\begin{proof}
\uses{topologies-lemma-morphism-big, topologies-lemma-Zariski-usual}

This follows from the simple description of pushforward
and pullback for the functors on the big sites from
Lemma \ref{topologies-lemma-morphism-big}. For the functors
on the small sites this is
Sheaves, Lemma \ref{sheaves-lemma-pushforward-composition}
via the identification of Lemma \ref{topologies-lemma-Zariski-usual}.
\end{proof}

\begin{lemma}[Morphism big small cartesian diagram]
\label{topologies-lemma-morphism-big-small-cartesian-diagram}
\source{stacks:0DD9}
\href{https://stacks.math.columbia.edu/tag/0DD9}{\texttt{Stacks tag 0DD9}}
\group{topologies}

Let $\Sch_{Zar}$ be a big Zariski site. Consider a cartesian diagram
$$
\begin{array}{cc}T' \underset{g'}{\longrightarrow}\underset{f'}{\downarrow} & T \overset{f}{\downarrow} \\ S' \overset{g}{\longrightarrow} & S\end{array}
$$
in $\Sch_{Zar}$. Then
$i_g^{-1} \circ f_{big, *} = f'_{small, *} \circ (i_{g'})^{-1}$
and $g_{big}^{-1} \circ f_{big, *} = f'_{big, *} \circ (g'_{big})^{-1}$.
\end{lemma}

\begin{proof}
\uses{topologies-lemma-put-in-T, topologies-lemma-morphism-big-small}

Since the diagram is cartesian, we have for $U'/S'$
that $U' \times_{S'} T' = U' \times_S T$. Hence both
$i_g^{-1} \circ f_{big, *}$ and $f'_{small, *} \circ (i_{g'})^{-1}$
send a sheaf $\mathcal{F}$ on $(\Sch/T)_{Zar}$ to the sheaf
$U' \mapsto \mathcal{F}(U' \times_{S'} T')$ on $S'_{Zar}$
(use Lemmas \ref{topologies-lemma-put-in-T} and \ref{topologies-lemma-morphism-big-small}).
The second equality can be proved in the same manner or can be
deduced from the very general
Sites, Lemma \ref{sites-lemma-localize-morphism}.
\end{proof}

We can think about a sheaf on the big Zariski site of $S$ as a collection
of ``usual'' sheaves on all schemes over $S$.

\begin{lemma}[Characterize sheaf big]
\label{topologies-lemma-characterize-sheaf-big}
\source{stacks:0213}
\href{https://stacks.math.columbia.edu/tag/0213}{\texttt{Stacks tag 0213}}
\group{topologies}

Let $S$ be a scheme contained in a big Zariski site $\Sch_{Zar}$.
A sheaf $\mathcal{F}$ on the big Zariski site $(\Sch/S)_{Zar}$
is given by the following data:
\begin{enumerate}
\item for every $T/S \in \Ob((\Sch/S)_{Zar})$ a sheaf
$\mathcal{F}_T$ on $T$,
\item for every $f : T' \to T$ in
$(\Sch/S)_{Zar}$ a map
$c_f : f^{-1}\mathcal{F}_T \to \mathcal{F}_{T'}$.
\end{enumerate}
These data are subject to the following conditions:
\begin{enumerate}
\item[(a)] given any $f : T' \to T$ and $g : T'' \to T'$ in
$(\Sch/S)_{Zar}$ the composition $c_g \circ g^{-1}c_f$
is equal to $c_{f \circ g}$, and
\item[(b)] if $f : T' \to T$ in $(\Sch/S)_{Zar}$ is an
open immersion then $c_f$ is an isomorphism.
\end{enumerate}
\end{lemma}

\begin{proof}
\uses{topologies-lemma-Zariski-usual, topologies-lemma-put-in-T}

This lemma follows from a purely sheaf theoretic statement discussed
in Sites, Remark \ref{sites-remark-variant-glue-sheaves-absolute}.
We also give a direct proof in this case.

\medskip\noindent
Given a sheaf $\mathcal{F}$ on $\Sh((\Sch/S)_{Zar})$
we set $\mathcal{F}_T = i_p^{-1}\mathcal{F}$ where $p : T \to S$
is the structure morphism. Note that
$\mathcal{F}_T(U) = \mathcal{F}(U'/S)$ for any open $U \subset T$,
and $U' \to T$ an open immersion in $(\Sch/T)_{Zar}$
with image $U$, see Lemmas \ref{topologies-lemma-Zariski-usual} and \ref{topologies-lemma-put-in-T}.
Hence given $f : T' \to T$ over $S$ and $U, U' \to T$ we get a canonical
map $\mathcal{F}_T(U) = \mathcal{F}(U'/S) \to \mathcal{F}(U'\times_T T'/S)
= \mathcal{F}_{T'}(f^{-1}(U))$ where the middle is the restriction map
of $\mathcal{F}$ with respect to the morphism
$U' \times_T T' \to U'$ over $S$. The collection of these maps are
compatible with restrictions, and hence define an $f$-map $c_f$
from $\mathcal{F}_T$ to $\mathcal{F}_{T'}$, see
Sheaves, Definition \ref{sheaves-definition-f-map} and the discussion
surrounding it. It is clear that $c_{f \circ g}$ is the composition of
$c_f$ and $c_g$, since composition of restriction maps of $\mathcal{F}$
gives restriction maps.

\medskip\noindent
Conversely, given a system $(\mathcal{F}_T, c_f)$ as in the lemma
we may define a presheaf $\mathcal{F}$ on $\Sh((\Sch/S)_{Zar})$
by simply setting $\mathcal{F}(T/S) = \mathcal{F}_T(T)$. As restriction
mapping, given $f : T' \to T$ we set for $s \in \mathcal{F}(T)$
the pullback $f^*(s)$ equal to $c_f(s)$ (where we think of $c_f$ as
an $f$-map again). The condition on the $c_f$ guarantees that
pullbacks satisfy the required functoriality property.
We omit the verification that this is a sheaf.
It is clear that the constructions so defined are mutually inverse.
\end{proof}

\section{The \'etale topology}
\label{topologies-section-etale}
Let $S$ be a scheme. We would like to define the \'etale-topology on
the category of schemes over $S$. According to our general principle
we first introduce the notion of an \'etale covering.

\begin{definition}[Etale covering]
\label{topologies-definition-etale-covering}
\source{stacks:0215}
\href{https://stacks.math.columbia.edu/tag/0215}{\texttt{Stacks tag 0215}}
\group{topologies}

Let $T$ be a scheme. An {\it \'etale covering of $T$} is a family
of morphisms $\{f_i : T_i \to T\}_{i \in I}$ of schemes
such that each $f_i$ is \'etale and such that $T = \bigcup f_i(T_i)$.
\end{definition}

\begin{lemma}[Zariski etale]
\label{topologies-lemma-zariski-etale}
\source{stacks:0216}
\href{https://stacks.math.columbia.edu/tag/0216}{\texttt{Stacks tag 0216}}
\group{topologies}

Any Zariski covering is an \'etale covering.
\end{lemma}

\begin{proof}
This is clear from the definitions and the fact that an open immersion
is an \'etale morphism, see
Morphisms, Lemma \ref{morphisms-lemma-open-immersion-etale}.
\end{proof}

Next, we show that this notion satisfies the conditions of
Sites, Definition \ref{sites-definition-site}.

\begin{lemma}[Etale]
\label{topologies-lemma-etale}
\source{stacks:0217}
\href{https://stacks.math.columbia.edu/tag/0217}{\texttt{Stacks tag 0217}}
\group{topologies}

Let $T$ be a scheme.
\begin{enumerate}
\item If $T' \to T$ is an isomorphism then $\{T' \to T\}$
is an \'etale covering of $T$.
\item If $\{T_i \to T\}_{i\in I}$ is an \'etale covering and for each
$i$ we have an \'etale covering $\{T_{ij} \to T_i\}_{j\in J_i}$, then
$\{T_{ij} \to T\}_{i \in I, j\in J_i}$ is an \'etale covering.
\item If $\{T_i \to T\}_{i\in I}$ is an \'etale covering
and $T' \to T$ is a morphism of schemes then
$\{T' \times_T T_i \to T'\}_{i\in I}$ is an \'etale covering.
\end{enumerate}
\end{lemma}

\begin{proof}
Omitted.
\end{proof}

\begin{lemma}[Etale affine]
\label{topologies-lemma-etale-affine}
\source{stacks:0218}
\href{https://stacks.math.columbia.edu/tag/0218}{\texttt{Stacks tag 0218}}
\group{topologies}

Let $T$ be an affine scheme.
Let $\{T_i \to T\}_{i \in I}$ be an \'etale covering of $T$.
Then there exists an \'etale covering
$\{U_j \to T\}_{j = 1, \ldots, m}$ which is a refinement
of $\{T_i \to T\}_{i \in I}$ such that each $U_j$ is an affine
scheme. Moreover, we may choose each $U_j$ to be open affine
in one of the $T_i$.
\end{lemma}

\begin{proof}
Omitted.
\end{proof}

Thus we define the corresponding standard coverings of affines as follows.

\begin{definition}[Standard etale]
\label{topologies-definition-standard-etale}
\source{stacks:0219}
\href{https://stacks.math.columbia.edu/tag/0219}{\texttt{Stacks tag 0219}}
\group{topologies}

Let $T$ be an affine scheme. A {\it standard \'etale covering}
of $T$ is a family $\{f_j : U_j \to T\}_{j = 1, \ldots, m}$
with each $U_j$ is affine and \'etale over $T$ and
$T = \bigcup f_j(U_j)$.
\end{definition}

In the definition above we do {\bf not} assume the morphisms $f_j$ are
standard \'etale. The reason is that if we did then the standard \'etale
coverings would not define a site on $\textit{Aff}/S$, for example because of
Algebra, Lemma \ref{algebra-lemma-standard-etale} part (4).
On the other hand, an \'etale morphism of affines is automatically
standard smooth, see
Algebra, Lemma \ref{algebra-lemma-etale-standard-smooth}.
Hence a standard \'etale covering is a standard smooth
covering and a standard syntomic covering.

\begin{definition}[Big etale site]
\label{topologies-definition-big-etale-site}
\source{stacks:021A}
\href{https://stacks.math.columbia.edu/tag/021A}{\texttt{Stacks tag 021A}}
\group{topologies}

A {\it big \'etale site} is any site $\Sch_\etale$ as in
Sites, Definition \ref{sites-definition-site} constructed as follows:
\begin{enumerate}
\item Choose any set of schemes $S_0$, and any set of \'etale coverings
$\text{Cov}_0$ among these schemes.
\item As underlying category take any category $\Sch_\alpha$
constructed as in Sets, Lemma \ref{sets-lemma-construct-category}
starting with the set $S_0$.
\item Choose any set of coverings as in
Sets, Lemma \ref{sets-lemma-coverings-site} starting with the
category $\Sch_\alpha$ and the class of \'etale coverings,
and the set $\text{Cov}_0$ chosen above.
\end{enumerate}
\end{definition}

See the remarks following Definition \ref{topologies-definition-big-zariski-site}
for motivation and explanation regarding the definition of big sites.

\medskip\noindent
Before we continue with the introduction of the big \'etale site of
a scheme $S$, let us point out that the topology on a big \'etale site
$\Sch_\etale$ is in some sense induced from the \'etale
topology on the category of all schemes.

\begin{lemma}[Etale induced]
\label{topologies-lemma-etale-induced}
\source{stacks:03WW}
\href{https://stacks.math.columbia.edu/tag/03WW}{\texttt{Stacks tag 03WW}}
\group{topologies}

\uses{topologies-definition-big-etale-site}
Let $\Sch_\etale$ be a big \'etale site as in
Definition \ref{topologies-definition-big-etale-site}.
Let $T \in \Ob(\Sch_\etale)$.
Let $\{T_i \to T\}_{i \in I}$ be an arbitrary \'etale covering of $T$.
\begin{enumerate}
\item There exists a covering $\{U_j \to T\}_{j \in J}$ of $T$ in the site
$\Sch_\etale$ which refines $\{T_i \to T\}_{i \in I}$.
\item If $\{T_i \to T\}_{i \in I}$ is a standard \'etale covering, then
it is tautologically equivalent to a covering in $\Sch_\etale$.
\item If $\{T_i \to T\}_{i \in I}$ is a Zariski covering, then
it is tautologically equivalent to a covering in $\Sch_\etale$.
\end{enumerate}
\end{lemma}

\begin{proof}
\uses{topologies-lemma-etale}

For each $i$ choose an affine open covering $T_i = \bigcup_{j \in J_i} T_{ij}$
such that each $T_{ij}$ maps into an affine open subscheme of $T$. By
Lemma \ref{topologies-lemma-etale}
the refinement $\{T_{ij} \to T\}_{i \in I, j \in J_i}$ is an \'etale covering
of $T$ as well. Hence we may assume each $T_i$ is affine, and maps into
an affine open $W_i$ of $T$. Applying
Sets, Lemma \ref{sets-lemma-what-is-in-it}
we see that $W_i$ is isomorphic to an object of $\Sch_\etale$.
But then $T_i$ as a finite type scheme over $W_i$
is isomorphic to an object $V_i$ of $\Sch_\etale$ by a second
application of
Sets, Lemma \ref{sets-lemma-what-is-in-it}.
The covering $\{V_i \to T\}_{i \in I}$ refines $\{T_i \to T\}_{i \in I}$
(because they are isomorphic).
Moreover, $\{V_i \to T\}_{i \in I}$ is combinatorially equivalent to a
covering $\{U_j \to T\}_{j \in J}$ of $T$ in the site
$\Sch_\etale$ by
Sets, Lemma \ref{sets-lemma-what-is-in-it}.
The covering $\{U_j \to T\}_{j \in J}$ is a refinement as in (1).
In the situation of (2), (3) each of the
schemes $T_i$ is isomorphic to an object of $\Sch_\etale$ by
Sets, Lemma \ref{sets-lemma-what-is-in-it},
and another application of
Sets, Lemma \ref{sets-lemma-coverings-site}
gives what we want.
\end{proof}

\begin{definition}[Big small etale]
\label{topologies-definition-big-small-etale}
\source{stacks:021B}
\href{https://stacks.math.columbia.edu/tag/021B}{\texttt{Stacks tag 021B}}
\group{topologies}

Let $S$ be a scheme. Let $\Sch_\etale$ be a big \'etale
site containing $S$.
\begin{enumerate}
\item The {\it big \'etale site of $S$}, denoted
$(\Sch/S)_\etale$, is the site
$\Sch_\etale/S$ introduced in
Sites, Section \ref{sites-section-localize}.
\item The {\it small \'etale site of $S$}, which we denote
$S_\etale$, is the full subcategory of
$(\Sch/S)_\etale$
whose objects are those $U/S$ such that $U \to S$ is \'etale.
A covering of $S_\etale$ is any covering $\{U_i \to U\}$ of
$(\Sch/S)_\etale$ with $U \in \Ob(S_\etale)$.
\item The {\it big affine \'etale site of $S$}, denoted
$(\textit{Aff}/S)_\etale$, is the full subcategory of
$(\Sch/S)_\etale$ whose objects are those $U/S$
such that $U$ is an affine scheme.
A covering of $(\textit{Aff}/S)_\etale$ is any covering
$\{U_i \to U\}$ of $(\Sch/S)_\etale$ with $U \in \Ob((\textit{Aff}/S)_\etale)$
which is a standard \'etale covering.
\item The {\it small affine \'etale site of $S$}, denoted
$S_{affine, \etale}$, is the full subcategory of $S_\etale$
whose objects are those $U/S$ such that $U$ is an affine scheme.
A covering of $S_{affine, \etale}$ is any covering $\{U_i \to U\}$
of $S_\etale$ with $U \in \Ob(S_{affine, \etale})$ which is a standard
\'etale covering.
\end{enumerate}
\end{definition}

It is not completely clear that the big affine \'etale site,
the small \'etale site, and the small affine \'etale site are sites.
We check this now.

\begin{lemma}[Verify site etale]
\label{topologies-lemma-verify-site-etale}
\source{stacks:021C}
\href{https://stacks.math.columbia.edu/tag/021C}{\texttt{Stacks tag 021C}}
\group{topologies}

Let $S$ be a scheme. Let $\Sch_\etale$ be a big \'etale
site containing $S$.
The structures $S_\etale$, $(\textit{Aff}/S)_\etale$, and $S_{affine, \etale}$
are sites.
\end{lemma}

\begin{proof}
Let us show that $S_\etale$ is a site. It is a category with a
given set of families of morphisms with fixed target. Thus we
have to show properties (1), (2) and (3) of
Sites, Definition \ref{sites-definition-site}.
Since $(\Sch/S)_\etale$ is a site, it suffices to prove
that given any covering $\{U_i \to U\}$ of $(\Sch/S)_\etale$
with $U \in \Ob(S_\etale)$ we also have
$U_i \in \Ob(S_\etale)$.
This follows from the definitions as the composition of \'etale morphisms
is an \'etale morphism.

\medskip\noindent
Let us show that $(\textit{Aff}/S)_\etale$ is a site.
Reasoning as above, it suffices to show that the collection
of standard \'etale coverings of affines satisfies properties
(1), (2) and (3) of
Sites, Definition \ref{sites-definition-site}.
This is clear since for example, given a standard \'etale
covering $\{T_i \to T\}_{i\in I}$ and for each
$i$ we have a standard \'etale covering $\{T_{ij} \to T_i\}_{j\in J_i}$, then
$\{T_{ij} \to T\}_{i \in I, j\in J_i}$ is a standard \'etale covering
because $\bigcup_{i\in I} J_i$ is finite and each $T_{ij}$ is affine.

\medskip\noindent
We omit the proof that $S_{affine, \'etale}$ is a site.
\end{proof}

\begin{lemma}[Fibre products etale]
\label{topologies-lemma-fibre-products-etale}
\source{stacks:021D}
\href{https://stacks.math.columbia.edu/tag/021D}{\texttt{Stacks tag 021D}}
\group{topologies}

Let $S$ be a scheme. Let $\Sch_\etale$ be a big \'etale
site containing $S$. The underlying categories of the sites
$\Sch_\etale$, $(\Sch/S)_\etale$, $S_\etale$, $(\textit{Aff}/S)_\etale$,
and $S_{affine, \etale}$ have fibre products.
In each case the obvious functor into the category $\Sch$ of
all schemes commutes with taking fibre products. The categories
$(\Sch/S)_\etale$, and $S_\etale$ both have a
final object, namely $S/S$.
\end{lemma}

\begin{proof}
For $\Sch_\etale$ it is true by construction, see
Sets, Lemma \ref{sets-lemma-what-is-in-it}.
Suppose we have $U \to S$, $V \to U$, $W \to U$ morphisms
of schemes with $U, V, W \in \Ob(\Sch_\etale)$.
The fibre product $V \times_U W$ in $\Sch_\etale$
is a fibre product in $\Sch$ and
is the fibre product of $V/S$ with $W/S$ over $U/S$ in
the category of all schemes over $S$, and hence also a
fibre product in $(\Sch/S)_\etale$.
This proves the result for $(\Sch/S)_\etale$.
If $U \to S$, $V \to U$ and $W \to U$ are \'etale then so is
$V \times_U W \to S$ and hence we get the result for $S_\etale$.
If $U, V, W$ are affine, so is $V \times_U W$ and hence the
result for $(\textit{Aff}/S)_\etale$ and $S_{affine, \etale}$.
\end{proof}

Next, we check that the big, resp.\ small affine site defines the same
topos as the big, resp.\ small site.

\begin{lemma}[Affine big site etale]
\label{topologies-lemma-affine-big-site-etale}
\source{stacks:021E}
\href{https://stacks.math.columbia.edu/tag/021E}{\texttt{Stacks tag 021E}}
\group{topologies}

Let $S$ be a scheme. Let $\Sch_\etale$ be a big \'etale
site containing $S$.
The functor
$(\textit{Aff}/S)_\etale \to (\Sch/S)_\etale$
is special cocontinuous and induces an equivalence of topoi from
$\Sh((\textit{Aff}/S)_\etale)$ to
$\Sh((\Sch/S)_\etale)$.
\end{lemma}

\begin{proof}
\uses{topologies-lemma-etale-affine}

The notion of a special cocontinuous functor is introduced in
Sites, Definition \ref{sites-definition-special-cocontinuous-functor}.
Thus we have to verify assumptions (1) -- (5) of
Sites, Lemma \ref{sites-lemma-equivalence}.
Denote the inclusion functor
$u : (\textit{Aff}/S)_\etale \to (\Sch/S)_\etale$.
Being cocontinuous just means that any \'etale covering of
$T/S$, $T$ affine, can be refined by a standard \'etale covering of $T$.
This is the content of
Lemma \ref{topologies-lemma-etale-affine}.
Hence (1) holds. We see $u$ is continuous simply because a standard
\'etale covering is a \'etale covering. Hence (2) holds.
Parts (3) and (4) follow immediately from the fact that $u$ is
fully faithful. And finally condition (5) follows from the
fact that every scheme has an affine open covering.
\end{proof}

\begin{lemma}[Alternative]
\label{topologies-lemma-alternative}
\source{stacks:04HR}
\href{https://stacks.math.columbia.edu/tag/04HR}{\texttt{Stacks tag 04HR}}
\group{topologies}

Let $S$ be a scheme. Let $\Sch_\etale$ be a big \'etale
site containing $S$. The functor $S_{affine, \etale} \to S_\etale$
is special cocontinuous and induces an equivalence of topoi from
$\Sh(S_{affine, \etale})$ to $\Sh(S_\etale)$.
\end{lemma}

\begin{proof}
\uses{topologies-lemma-affine-big-site-etale}

Omitted. Hint: compare with the proof of
Lemma \ref{topologies-lemma-affine-big-site-etale}.
\end{proof}

Next, we establish some relationships between the topoi
associated to these sites.

\begin{lemma}[Put in T etale]
\label{topologies-lemma-put-in-T-etale}
\source{stacks:021F}
\href{https://stacks.math.columbia.edu/tag/021F}{\texttt{Stacks tag 021F}}
\group{topologies}

Let $\Sch_\etale$ be a big \'etale site.
Let $f : T \to S$ be a morphism in $\Sch_\etale$.
The functor $T_\etale \to (\Sch/S)_\etale$
is cocontinuous and induces a morphism of topoi
$$
i_f :
\Sh(T_\etale)
\longrightarrow
\Sh((\Sch/S)_\etale)
$$
For a sheaf $\mathcal{G}$ on $(\Sch/S)_\etale$
we have the formula $(i_f^{-1}\mathcal{G})(U/T) = \mathcal{G}(U/S)$.
The functor $i_f^{-1}$ also has a left adjoint $i_{f, !}$ which commutes
with fibre products and equalizers.
\end{lemma}

\begin{proof}
\uses{topologies-lemma-fibre-products-etale}

Denote the functor $u : T_\etale \to (\Sch/S)_\etale$.
In other words, given an \'etale morphism $j : U \to T$ corresponding
to an object of $T_\etale$ we set $u(U \to T) = (f \circ j : U \to S)$.
This functor commutes with fibre products, see
Lemma \ref{topologies-lemma-fibre-products-etale}.
Let $a, b : U \to V$ be two morphisms in $T_\etale$.
In this case the equalizer of $a$ and $b$ (in the category of schemes) is
$$
V \times_{\Delta_{V/T}, V \times_T V, (a, b)} U
$$
which is a fibre product of schemes \'etale over $T$, hence \'etale
over $T$. Thus $T_\etale$ has equalizers and $u$ commutes with them.
It is clearly cocontinuous.
It is also continuous as $u$ transforms coverings to coverings and
commutes with fibre products. Hence the Lemma follows from
Sites, Lemmas \ref{sites-lemma-when-shriek}
and \ref{sites-lemma-preserve-equalizers}.
\end{proof}

\begin{lemma}[At the bottom etale]
\label{topologies-lemma-at-the-bottom-etale}
\source{stacks:021G}
\href{https://stacks.math.columbia.edu/tag/021G}{\texttt{Stacks tag 021G}}
\group{topologies}

\uses{topologies-lemma-put-in-T-etale}
Let $S$ be a scheme. Let $\Sch_\etale$ be a big \'etale
site containing $S$.
The inclusion functor $S_\etale \to (\Sch/S)_\etale$
satisfies the hypotheses of Sites, Lemma \ref{sites-lemma-bigger-site}
and hence induces a morphism of sites
$$
\pi_S : (\Sch/S)_\etale \longrightarrow S_\etale
$$
and a morphism of topoi
$$
i_S : \Sh(S_\etale) \longrightarrow \Sh((\Sch/S)_\etale)
$$
such that $\pi_S \circ i_S = \text{id}$. Moreover, $i_S = i_{\text{id}_S}$
with $i_{\text{id}_S}$ as in Lemma \ref{topologies-lemma-put-in-T-etale}.
In particular the functor $i_S^{-1} = \pi_{S, *}$ is described by the rule
$i_S^{-1}(\mathcal{G})(U/S) = \mathcal{G}(U/S)$.
\end{lemma}

\begin{proof}
\uses{topologies-lemma-put-in-T-etale}

In this case the functor
$u : S_\etale \to (\Sch/S)_\etale$,
in addition to the properties seen in the proof of
Lemma \ref{topologies-lemma-put-in-T-etale} above, also is fully faithful
and transforms the final object into the final object.
The lemma follows from Sites, Lemma \ref{sites-lemma-bigger-site}.
\end{proof}

\begin{definition}[Restriction small etale]
\label{topologies-definition-restriction-small-etale}
\source{stacks:04BT}
\href{https://stacks.math.columbia.edu/tag/04BT}{\texttt{Stacks tag 04BT}}
\group{topologies}

\uses{topologies-lemma-at-the-bottom-etale}
In the situation of
Lemma \ref{topologies-lemma-at-the-bottom-etale}
the functor $i_S^{-1} = \pi_{S, *}$ is often
called the {\it restriction to the small \'etale site}, and for a sheaf
$\mathcal{F}$ on the big \'etale site we denote
$\mathcal{F}|_{S_\etale}$ this restriction.
\end{definition}

With this notation in place we have for a sheaf $\mathcal{F}$ on the
big site and a sheaf $\mathcal{G}$ on the small site that
\begin{align*}
\Mor_{\Sh(S_\etale)}(
\mathcal{F}|_{S_\etale},
\mathcal{G})
& =
\Mor_{\Sh((\Sch/S)_\etale)}(
\mathcal{F},
i_{S, *}\mathcal{G}) \\
\Mor_{\Sh(S_\etale)}(
\mathcal{G},
\mathcal{F}|_{S_\etale})
& =
\Mor_{\Sh((\Sch/S)_\etale)}(
\pi_S^{-1}\mathcal{G},
\mathcal{F})
\end{align*}
Moreover, we have $(i_{S, *}\mathcal{G})|_{S_\etale} = \mathcal{G}$
and we have $(\pi_S^{-1}\mathcal{G})|_{S_\etale} = \mathcal{G}$.

\begin{lemma}[Morphism big etale]
\label{topologies-lemma-morphism-big-etale}
\source{stacks:021H}
\href{https://stacks.math.columbia.edu/tag/021H}{\texttt{Stacks tag 021H}}
\group{topologies}

Let $\Sch_\etale$ be a big \'etale site.
Let $f : T \to S$ be a morphism in $\Sch_\etale$.
The functor
$$
u :
(\Sch/T)_\etale
\longrightarrow
(\Sch/S)_\etale,
\quad
V/T \longmapsto V/S
$$
is cocontinuous, and has a continuous right adjoint
$$
v :
(\Sch/S)_\etale
\longrightarrow
(\Sch/T)_\etale,
\quad
(U \to S) \longmapsto (U \times_S T \to T).
$$
They induce the same morphism of topoi
$$
f_{big} :
\Sh((\Sch/T)_\etale)
\longrightarrow
\Sh((\Sch/S)_\etale)
$$
We have $f_{big}^{-1}(\mathcal{G})(U/T) = \mathcal{G}(U/S)$.
We have $f_{big, *}(\mathcal{F})(U/S) = \mathcal{F}(U \times_S T/T)$.
Also, $f_{big}^{-1}$ has a left adjoint $f_{big!}$ which commutes with
fibre products and equalizers.
\end{lemma}

\begin{proof}
\uses{topologies-lemma-put-in-T-etale}

The functor $u$ is cocontinuous, continuous and commutes with fibre products
and equalizers (details omitted; compare with the proof of
Lemma \ref{topologies-lemma-put-in-T-etale}).
Hence
Sites, Lemmas \ref{sites-lemma-when-shriek} and
\ref{sites-lemma-preserve-equalizers}
apply and we deduce the formula
for $f_{big}^{-1}$ and the existence of $f_{big!}$. Moreover,
the functor $v$ is a right adjoint because given $U/T$ and $V/S$
we have $\Mor_S(u(U), V) = \Mor_T(U, V \times_S T)$
as desired. Thus we may apply
Sites, Lemmas \ref{sites-lemma-have-functor-other-way} and
\ref{sites-lemma-have-functor-other-way-morphism} to get the
formula for $f_{big, *}$.
\end{proof}

\begin{lemma}[Morphism big small etale]
\label{topologies-lemma-morphism-big-small-etale}
\source{stacks:021I}
\href{https://stacks.math.columbia.edu/tag/021I}{\texttt{Stacks tag 021I}}
\group{topologies}

\uses{topologies-lemma-put-in-T-etale, topologies-lemma-at-the-bottom-etale}
Let $\Sch_\etale$ be a big \'etale site.
Let $f : T \to S$ be a morphism in $\Sch_\etale$.
\begin{enumerate}
\item We have $i_f = f_{big} \circ i_T$ with $i_f$ as in
Lemma \ref{topologies-lemma-put-in-T-etale} and $i_T$ as in
Lemma \ref{topologies-lemma-at-the-bottom-etale}.
\item The functor $S_\etale \to T_\etale$,
$(U \to S) \mapsto (U \times_S T \to T)$ is continuous and induces
a morphism of sites
$$
f_{small} : T_\etale \longrightarrow S_\etale
$$
We have $f_{small, *}(\mathcal{F})(U/S) = \mathcal{F}(U \times_S T/T)$.
\item We have a commutative diagram of morphisms of sites
$$
\begin{array}{cc}T_\etale \underset{f_{small}}{\downarrow} & (\Sch/T)_\etale \overset{f_{big}}{\downarrow}\overset{\pi_T}{\longleftarrow} \\ S_\etale & (\Sch/S)_\etale \underset{\pi_S}{\longleftarrow}\end{array}
$$
so that $f_{small} \circ \pi_T = \pi_S \circ f_{big}$ as morphisms of topoi.
\item We have $f_{small} = \pi_S \circ f_{big} \circ i_T = \pi_S \circ i_f$.
\end{enumerate}
\end{lemma}

\begin{proof}
\uses{topologies-lemma-etale, topologies-lemma-fibre-products-etale}

The equality $i_f = f_{big} \circ i_T$ follows from the
equality $i_f^{-1} = i_T^{-1} \circ f_{big}^{-1}$ which is
clear from the descriptions of these functors above.
Thus we see (1).

\medskip\noindent
The functor
$u :
S_\etale
\to
T_\etale$, $u(U \to S) = (U \times_S T \to T)$
transforms coverings into coverings and commutes with fibre products,
see Lemma \ref{topologies-lemma-etale} (3) and \ref{topologies-lemma-fibre-products-etale}.
Moreover, both $S_\etale$, $T_\etale$ have final objects,
namely $S/S$ and $T/T$ and $u(S/S) = T/T$. Hence by
Sites, Proposition \ref{sites-proposition-get-morphism}
the functor $u$ corresponds to a morphism of sites
$T_\etale \to S_\etale$. This in turn gives rise to the
morphism of topoi, see
Sites, Lemma \ref{sites-lemma-morphism-sites-topoi}. The description
of the pushforward is clear from these references.

\medskip\noindent
Part (3) follows because $\pi_S$ and $\pi_T$ are given by the
inclusion functors and $f_{small}$ and $f_{big}$ by the
base change functors $U \mapsto U \times_S T$.

\medskip\noindent
Statement (4) follows from (3) by precomposing with $i_T$.
\end{proof}

In the situation of the lemma, using the terminology of
Definition \ref{topologies-definition-restriction-small-etale}
we have: for $\mathcal{F}$ a sheaf on the big \'etale site of $T$
$$
(f_{big, *}\mathcal{F})|_{S_\etale} =
f_{small, *}(\mathcal{F}|_{T_\etale}),
$$
This equality is clear from the commutativity of the diagram of
sites of the lemma, since restriction to the small \'etale site of
$T$, resp.\ $S$ is given by $\pi_{T, *}$, resp.\ $\pi_{S, *}$. A similar
formula involving pullbacks and restrictions is false.

\begin{lemma}[Composition etale]
\label{topologies-lemma-composition-etale}
\source{stacks:021J}
\href{https://stacks.math.columbia.edu/tag/021J}{\texttt{Stacks tag 021J}}
\group{topologies}

Given schemes $X$, $Y$, $Y$ in $\Sch_\etale$
and morphisms $f : X \to Y$, $g : Y \to Z$ we have
$g_{big} \circ f_{big} = (g \circ f)_{big}$ and
$g_{small} \circ f_{small} = (g \circ f)_{small}$.
\end{lemma}

\begin{proof}
\uses{topologies-lemma-morphism-big-etale, topologies-lemma-morphism-big-small-etale}

This follows from the simple description of pushforward
and pullback for the functors on the big sites from
Lemma \ref{topologies-lemma-morphism-big-etale}. For the functors
on the small sites this follows from the description of
the pushforward functors in Lemma \ref{topologies-lemma-morphism-big-small-etale}.
\end{proof}

\begin{lemma}[Morphism big small cartesian diagram etale]
\label{topologies-lemma-morphism-big-small-cartesian-diagram-etale}
\source{stacks:0DDA}
\href{https://stacks.math.columbia.edu/tag/0DDA}{\texttt{Stacks tag 0DDA}}
\group{topologies}

Let $\Sch_\etale$ be a big \'etale site. Consider a cartesian diagram
$$
\begin{array}{cc}T' \underset{g'}{\longrightarrow}\underset{f'}{\downarrow} & T \overset{f}{\downarrow} \\ S' \overset{g}{\longrightarrow} & S\end{array}
$$
in $\Sch_\etale$. Then
$i_g^{-1} \circ f_{big, *} = f'_{small, *} \circ (i_{g'})^{-1}$
and $g_{big}^{-1} \circ f_{big, *} = f'_{big, *} \circ (g'_{big})^{-1}$.
\end{lemma}

\begin{proof}
\uses{topologies-lemma-put-in-T-etale, topologies-lemma-morphism-big-etale}

Since the diagram is cartesian, we have for $U'/S'$
that $U' \times_{S'} T' = U' \times_S T$. Hence both
$i_g^{-1} \circ f_{big, *}$ and $f'_{small, *} \circ (i_{g'})^{-1}$
send a sheaf $\mathcal{F}$ on $(\Sch/T)_\etale$ to the sheaf
$U' \mapsto \mathcal{F}(U' \times_{S'} T')$ on $S'_\etale$
(use Lemmas \ref{topologies-lemma-put-in-T-etale} and \ref{topologies-lemma-morphism-big-etale}).
The second equality can be proved in the same manner or can be
deduced from the very general
Sites, Lemma \ref{sites-lemma-localize-morphism}.
\end{proof}

We can think about a sheaf on the big \'etale site of $S$ as a collection
of ``usual'' sheaves on all schemes over $S$.

\begin{lemma}[Characterize sheaf big etale]
\label{topologies-lemma-characterize-sheaf-big-etale}
\source{stacks:021K}
\href{https://stacks.math.columbia.edu/tag/021K}{\texttt{Stacks tag 021K}}
\group{topologies}

Let $S$ be a scheme contained in a big \'etale site
$\Sch_\etale$.
A sheaf $\mathcal{F}$ on the big \'etale site
$(\Sch/S)_\etale$ is given by the following data:
\begin{enumerate}
\item for every $T/S \in \Ob((\Sch/S)_\etale)$ a sheaf
$\mathcal{F}_T$ on $T_\etale$,
\item for every $f : T' \to T$ in
$(\Sch/S)_\etale$ a map
$c_f : f_{small}^{-1}\mathcal{F}_T \to \mathcal{F}_{T'}$.
\end{enumerate}
These data are subject to the following conditions:
\begin{enumerate}
\item[(a)] given any $f : T' \to T$ and $g : T'' \to T'$ in
$(\Sch/S)_\etale$ the composition
$c_g \circ g_{small}^{-1}c_f$ is equal to $c_{f \circ g}$, and
\item[(b)] if $f : T' \to T$ in $(\Sch/S)_\etale$
is \'etale then $c_f$ is an isomorphism.
\end{enumerate}
\end{lemma}

\begin{proof}
\uses{topologies-lemma-put-in-T-etale}

This lemma follows from a purely sheaf theoretic statement discussed
in Sites, Remark \ref{sites-remark-variant-glue-sheaves-absolute}.
We also give a direct proof in this case.

\medskip\noindent
Given a sheaf $\mathcal{F}$ on $\Sh((\Sch/S)_\etale)$
we set $\mathcal{F}_T = i_p^{-1}\mathcal{F}$ where $p : T \to S$
is the structure morphism. Note that
$\mathcal{F}_T(U) = \mathcal{F}(U/S)$ for any $U \to T$
in $T_\etale$ see Lemma \ref{topologies-lemma-put-in-T-etale}.
Hence given $f : T' \to T$ over $S$ and $U \to T$ we get a canonical
map $\mathcal{F}_T(U) = \mathcal{F}(U/S) \to \mathcal{F}(U \times_T T'/S)
= \mathcal{F}_{T'}(U \times_T T')$ where the middle is the restriction map
of $\mathcal{F}$ with respect to the morphism
$U \times_T T' \to U$ over $S$. The collection of these maps are
compatible with restrictions, and hence define a map
$c'_f : \mathcal{F}_T \to f_{small, *}\mathcal{F}_{T'}$ where
$u : T_\etale \to T'_\etale$ is the base change functor
associated to $f$. By adjunction of $f_{small, *}$ (see
Sites, Section \ref{sites-section-continuous-functors}) with
$f_{small}^{-1}$ this is the same as a map
$c_f : f_{small}^{-1}\mathcal{F}_T \to \mathcal{F}_{T'}$.
It is clear that $c'_{f \circ g}$ is the composition of
$c'_f$ and $f_{small, *}c'_g$, since composition of restriction maps
of $\mathcal{F}$ gives restriction maps, and this gives the desired
relationship among $c_f$, $c_g$ and $c_{f \circ g}$.

\medskip\noindent
Conversely, given a system $(\mathcal{F}_T, c_f)$ as in the lemma
we may define a presheaf $\mathcal{F}$ on
$\Sh((\Sch/S)_\etale)$
by simply setting $\mathcal{F}(T/S) = \mathcal{F}_T(T)$. As restriction
mapping, given $f : T' \to T$ we set for $s \in \mathcal{F}(T)$
the pullback $f^*(s)$ equal to $c_f(s)$ where we think of $c_f$ as
a map $\mathcal{F}_T \to f_{small, *}\mathcal{F}_{T'}$ again.
The condition on the $c_f$ guarantees that
pullbacks satisfy the required functoriality property.
We omit the verification that this is a sheaf.
It is clear that the constructions so defined are mutually inverse.
\end{proof}

\section{The smooth topology}
\label{topologies-section-smooth}
In this section we define the smooth topology.
This is a bit pointless as it will turn out later (see
More on Morphisms, Section \ref{more-morphisms-section-etale-over-smooth})
that this topology defines the same topos as the
\'etale topology. But still it makes sense and it is used
occasionally.

\begin{definition}[Smooth covering]
\label{topologies-definition-smooth-covering}
\source{stacks:021Z}
\href{https://stacks.math.columbia.edu/tag/021Z}{\texttt{Stacks tag 021Z}}
\group{topologies}

Let $T$ be a scheme. A {\it smooth covering of $T$} is a family
of morphisms $\{f_i : T_i \to T\}_{i \in I}$ of schemes
such that each $f_i$ is smooth and such
that $T = \bigcup f_i(T_i)$.
\end{definition}

\begin{lemma}[Zariski etale smooth]
\label{topologies-lemma-zariski-etale-smooth}
\source{stacks:0220}
\href{https://stacks.math.columbia.edu/tag/0220}{\texttt{Stacks tag 0220}}
\group{topologies}

Any \'etale covering is a smooth covering, and a fortiori,
any Zariski covering is a smooth covering.
\end{lemma}

\begin{proof}
\uses{topologies-lemma-zariski-etale}

This is clear from the definitions, the fact that an \'etale morphism is
smooth see
Morphisms, Definition \ref{morphisms-definition-etale}
and Lemma \ref{topologies-lemma-zariski-etale}.
\end{proof}

Next, we show that this notion satisfies the conditions of
Sites, Definition \ref{sites-definition-site}.

\begin{lemma}[Smooth]
\label{topologies-lemma-smooth}
\source{stacks:0221}
\href{https://stacks.math.columbia.edu/tag/0221}{\texttt{Stacks tag 0221}}
\group{topologies}

Let $T$ be a scheme.
\begin{enumerate}
\item If $T' \to T$ is an isomorphism then $\{T' \to T\}$
is a smooth covering of $T$.
\item If $\{T_i \to T\}_{i\in I}$ is a smooth covering and for each
$i$ we have a smooth covering $\{T_{ij} \to T_i\}_{j\in J_i}$, then
$\{T_{ij} \to T\}_{i \in I, j\in J_i}$ is a smooth covering.
\item If $\{T_i \to T\}_{i\in I}$ is a smooth covering
and $T' \to T$ is a morphism of schemes then
$\{T' \times_T T_i \to T'\}_{i\in I}$ is a smooth covering.
\end{enumerate}
\end{lemma}

\begin{proof}
Omitted.
\end{proof}

\begin{lemma}[Smooth affine]
\label{topologies-lemma-smooth-affine}
\source{stacks:0222}
\href{https://stacks.math.columbia.edu/tag/0222}{\texttt{Stacks tag 0222}}
\group{topologies}

Let $T$ be an affine scheme.
Let $\{T_i \to T\}_{i \in I}$ be a smooth covering of $T$.
Then there exists a smooth covering
$\{U_j \to T\}_{j = 1, \ldots, m}$ which is a refinement
of $\{T_i \to T\}_{i \in I}$ such that each $U_j$ is an affine
scheme, and such that each morphism $U_j \to T$ is standard
smooth, see Morphisms, Definition \ref{morphisms-definition-smooth}.
Moreover, we may choose each $U_j$ to be open affine in one of the $T_i$.
\end{lemma}

\begin{proof}
Omitted, but see Algebra, Lemma \ref{algebra-lemma-smooth-syntomic}.
\end{proof}

Thus we define the corresponding standard coverings of affines as follows.

\begin{definition}[Standard smooth]
\label{topologies-definition-standard-smooth}
\source{stacks:0223}
\href{https://stacks.math.columbia.edu/tag/0223}{\texttt{Stacks tag 0223}}
\group{topologies}

Let $T$ be an affine scheme. A {\it standard smooth covering}
of $T$ is a family $\{f_j : U_j \to T\}_{j = 1, \ldots, m}$
with each $U_j$ is affine, $U_j \to T$ standard smooth
and $T = \bigcup f_j(U_j)$.
\end{definition}

\begin{definition}[Big smooth site]
\label{topologies-definition-big-smooth-site}
\source{stacks:03WY}
\href{https://stacks.math.columbia.edu/tag/03WY}{\texttt{Stacks tag 03WY}}
\group{topologies}

A {\it big smooth site} is any site $\Sch_{smooth}$ as in
Sites, Definition \ref{sites-definition-site} constructed as follows:
\begin{enumerate}
\item Choose any set of schemes $S_0$, and any set of smooth coverings
$\text{Cov}_0$ among these schemes.
\item As underlying category take any category $\Sch_\alpha$
constructed as in Sets, Lemma \ref{sets-lemma-construct-category}
starting with the set $S_0$.
\item Choose any set of coverings as in
Sets, Lemma \ref{sets-lemma-coverings-site} starting with the
category $\Sch_\alpha$ and the class of smooth coverings,
and the set $\text{Cov}_0$ chosen above.
\end{enumerate}
\end{definition}

See the remarks following Definition \ref{topologies-definition-big-zariski-site}
for motivation and explanation regarding the definition of big sites.

\medskip\noindent
Before we continue with the introduction of the big smooth site of
a scheme $S$, let us point out that the topology on a big smooth site
$\Sch_{smooth}$ is in some sense induced from the smooth topology
on the category of all schemes.

\begin{lemma}[Smooth induced]
\label{topologies-lemma-smooth-induced}
\source{stacks:03WZ}
\href{https://stacks.math.columbia.edu/tag/03WZ}{\texttt{Stacks tag 03WZ}}
\group{topologies}

\uses{topologies-definition-big-smooth-site}
Let $\Sch_{smooth}$ be a big smooth site as in
Definition \ref{topologies-definition-big-smooth-site}.
Let $T \in \Ob(\Sch_{smooth})$.
Let $\{T_i \to T\}_{i \in I}$ be an arbitrary smooth covering of $T$.
\begin{enumerate}
\item There exists a covering $\{U_j \to T\}_{j \in J}$ of $T$ in the site
$\Sch_{smooth}$ which refines $\{T_i \to T\}_{i \in I}$.
\item If $\{T_i \to T\}_{i \in I}$ is a standard smooth covering, then
it is tautologically equivalent to a covering of $\Sch_{smooth}$.
\item If $\{T_i \to T\}_{i \in I}$ is a Zariski covering, then
it is tautologically equivalent to a covering of $\Sch_{smooth}$.
\end{enumerate}
\end{lemma}

\begin{proof}
\uses{topologies-lemma-smooth}

For each $i$ choose an affine open covering $T_i = \bigcup_{j \in J_i} T_{ij}$
such that each $T_{ij}$ maps into an affine open subscheme of $T$. By
Lemma \ref{topologies-lemma-smooth}
the refinement $\{T_{ij} \to T\}_{i \in I, j \in J_i}$ is a smooth covering
of $T$ as well. Hence we may assume each $T_i$ is affine, and maps into
an affine open $W_i$ of $T$. Applying
Sets, Lemma \ref{sets-lemma-what-is-in-it}
we see that $W_i$ is isomorphic to an object of $\Sch_{smooth}$.
But then $T_i$ as a finite type scheme over $W_i$
is isomorphic to an object $V_i$ of $\Sch_{smooth}$ by a second
application of
Sets, Lemma \ref{sets-lemma-what-is-in-it}.
The covering $\{V_i \to T\}_{i \in I}$ refines $\{T_i \to T\}_{i \in I}$
(because they are isomorphic).
Moreover, $\{V_i \to T\}_{i \in I}$ is combinatorially equivalent to a
covering $\{U_j \to T\}_{j \in J}$ of $T$ in the site
$\Sch_{smooth}$ by
Sets, Lemma \ref{sets-lemma-what-is-in-it}.
The covering $\{U_j \to T\}_{j \in J}$ is a refinement as in (1).
In the situation of (2), (3) each of the
schemes $T_i$ is isomorphic to an object of $\Sch_{smooth}$ by
Sets, Lemma \ref{sets-lemma-what-is-in-it},
and another application of
Sets, Lemma \ref{sets-lemma-coverings-site}
gives what we want.
\end{proof}

\begin{definition}[Big small smooth]
\label{topologies-definition-big-small-smooth}
\source{stacks:03X0}
\href{https://stacks.math.columbia.edu/tag/03X0}{\texttt{Stacks tag 03X0}}
\group{topologies}

Let $S$ be a scheme. Let $\Sch_{smooth}$ be a big smooth
site containing $S$.
\begin{enumerate}
\item The {\it big smooth site of $S$}, denoted
$(\Sch/S)_{smooth}$, is the site $\Sch_{smooth}/S$
introduced in Sites, Section \ref{sites-section-localize}.
\item The {\it big affine smooth site of $S$}, denoted
$(\textit{Aff}/S)_{smooth}$, is the full subcategory of
$(\Sch/S)_{smooth}$ whose objects are affine $U/S$.
A covering of $(\textit{Aff}/S)_{smooth}$ is any covering
$\{U_i \to U\}$ of $(\Sch/S)_{smooth}$ which is a
standard smooth covering.
\end{enumerate}
\end{definition}

Next, we check that the big affine site defines the same
topos as the big site.

\begin{lemma}[Affine big site smooth]
\label{topologies-lemma-affine-big-site-smooth}
\source{stacks:06VC}
\href{https://stacks.math.columbia.edu/tag/06VC}{\texttt{Stacks tag 06VC}}
\group{topologies}

Let $S$ be a scheme. Let $\Sch_{smooth}$ be a big smooth
site containing $S$.
The functor
$(\textit{Aff}/S)_{smooth} \to (\Sch/S)_{smooth}$
is special cocontinuous and induces an equivalence of topoi from
$\Sh((\textit{Aff}/S)_{smooth})$ to
$\Sh((\Sch/S)_{smooth})$.
\end{lemma}

\begin{proof}
\uses{topologies-lemma-smooth-affine}

The notion of a special cocontinuous functor is introduced in
Sites, Definition \ref{sites-definition-special-cocontinuous-functor}.
Thus we have to verify assumptions (1) -- (5) of
Sites, Lemma \ref{sites-lemma-equivalence}.
Denote the inclusion functor
$u : (\textit{Aff}/S)_{smooth} \to (\Sch/S)_{smooth}$.
Being cocontinuous just means that any smooth covering of
$T/S$, $T$ affine, can be refined by a standard smooth covering of $T$.
This is the content of
Lemma \ref{topologies-lemma-smooth-affine}.
Hence (1) holds. We see $u$ is continuous simply because a standard
smooth covering is a smooth covering. Hence (2) holds.
Parts (3) and (4) follow immediately from the fact that $u$ is
fully faithful. And finally condition (5) follows from the
fact that every scheme has an affine open covering.
\end{proof}

To be continued...

\begin{lemma}[Morphism big smooth]
\label{topologies-lemma-morphism-big-smooth}
\source{stacks:04HC}
\href{https://stacks.math.columbia.edu/tag/04HC}{\texttt{Stacks tag 04HC}}
\group{topologies}

Let $\Sch_{smooth}$ be a big smooth site.
Let $f : T \to S$ be a morphism in $\Sch_{smooth}$.
The functor
$$
u : (\Sch/T)_{smooth} \longrightarrow (\Sch/S)_{smooth},
\quad
V/T \longmapsto V/S
$$
is cocontinuous, and has a continuous right adjoint
$$
v : (\Sch/S)_{smooth} \longrightarrow (\Sch/T)_{smooth},
\quad
(U \to S) \longmapsto (U \times_S T \to T).
$$
They induce the same morphism of topoi
$$
f_{big} :
\Sh((\Sch/T)_{smooth})
\longrightarrow
\Sh((\Sch/S)_{smooth})
$$
We have $f_{big}^{-1}(\mathcal{G})(U/T) = \mathcal{G}(U/S)$.
We have $f_{big, *}(\mathcal{F})(U/S) = \mathcal{F}(U \times_S T/T)$.
Also, $f_{big}^{-1}$ has a left adjoint $f_{big!}$ which commutes with
fibre products and equalizers.
\end{lemma}

\begin{proof}
The functor $u$ is cocontinuous, continuous, and commutes with fibre products
and equalizers. Hence
Sites, Lemmas \ref{sites-lemma-when-shriek} and
\ref{sites-lemma-preserve-equalizers}
apply and we deduce the formula
for $f_{big}^{-1}$ and the existence of $f_{big!}$. Moreover,
the functor $v$ is a right adjoint because given $U/T$ and $V/S$
we have $\Mor_S(u(U), V) = \Mor_T(U, V \times_S T)$
as desired. Thus we may apply
Sites, Lemmas \ref{sites-lemma-have-functor-other-way} and
\ref{sites-lemma-have-functor-other-way-morphism} to get the
formula for $f_{big, *}$.
\end{proof}

\section{The syntomic topology}
\label{topologies-section-syntomic}
In this section we define the syntomic topology.
This topology is quite interesting in that it often
has the same cohomology groups as the fppf topology
but is technically easier to deal with.

\begin{definition}[Syntomic covering]
\label{topologies-definition-syntomic-covering}
\source{stacks:0225}
\href{https://stacks.math.columbia.edu/tag/0225}{\texttt{Stacks tag 0225}}
\group{topologies}

Let $T$ be a scheme. An {\it syntomic covering of $T$} is a family
of morphisms $\{f_i : T_i \to T\}_{i \in I}$ of schemes
such that each $f_i$ is syntomic and such
that $T = \bigcup f_i(T_i)$.
\end{definition}

\begin{lemma}[Zariski etale smooth syntomic]
\label{topologies-lemma-zariski-etale-smooth-syntomic}
\source{stacks:0226}
\href{https://stacks.math.columbia.edu/tag/0226}{\texttt{Stacks tag 0226}}
\group{topologies}

Any smooth covering is a syntomic covering, and a fortiori,
any \'etale or Zariski covering is a syntomic covering.
\end{lemma}

\begin{proof}
\uses{topologies-lemma-zariski-etale-smooth}

This is clear from the definitions and the fact that a smooth
morphism is syntomic, see
Morphisms, Lemma \ref{morphisms-lemma-smooth-syntomic}
and Lemma \ref{topologies-lemma-zariski-etale-smooth}.
\end{proof}

Next, we show that this notion satisfies the conditions of
Sites, Definition \ref{sites-definition-site}.

\begin{lemma}[Syntomic]
\label{topologies-lemma-syntomic}
\source{stacks:0227}
\href{https://stacks.math.columbia.edu/tag/0227}{\texttt{Stacks tag 0227}}
\group{topologies}

Let $T$ be a scheme.
\begin{enumerate}
\item If $T' \to T$ is an isomorphism then $\{T' \to T\}$
is a syntomic covering of $T$.
\item If $\{T_i \to T\}_{i\in I}$ is a syntomic covering and for each
$i$ we have a syntomic covering $\{T_{ij} \to T_i\}_{j\in J_i}$, then
$\{T_{ij} \to T\}_{i \in I, j\in J_i}$ is a syntomic covering.
\item If $\{T_i \to T\}_{i\in I}$ is a syntomic covering
and $T' \to T$ is a morphism of schemes then
$\{T' \times_T T_i \to T'\}_{i\in I}$ is a syntomic covering.
\end{enumerate}
\end{lemma}

\begin{proof}
Omitted.
\end{proof}

\begin{lemma}[Syntomic affine]
\label{topologies-lemma-syntomic-affine}
\source{stacks:0228}
\href{https://stacks.math.columbia.edu/tag/0228}{\texttt{Stacks tag 0228}}
\group{topologies}

Let $T$ be an affine scheme.
Let $\{T_i \to T\}_{i \in I}$ be a syntomic covering of $T$.
Then there exists a syntomic covering
$\{U_j \to T\}_{j = 1, \ldots, m}$ which is a refinement
of $\{T_i \to T\}_{i \in I}$ such that each $U_j$ is an affine
scheme, and such that each morphism $U_j \to T$ is standard
syntomic, see Morphisms, Definition \ref{morphisms-definition-syntomic}.
Moreover, we may choose each $U_j$ to be open affine in one of the $T_i$.
\end{lemma}

\begin{proof}
Omitted, but see Algebra, Lemma \ref{algebra-lemma-syntomic}.
\end{proof}

Thus we define the corresponding standard coverings of affines as follows.

\begin{definition}[Standard syntomic]
\label{topologies-definition-standard-syntomic}
\source{stacks:0229}
\href{https://stacks.math.columbia.edu/tag/0229}{\texttt{Stacks tag 0229}}
\group{topologies}

Let $T$ be an affine scheme. A {\it standard syntomic covering} of $T$ is
a family $\{f_j : U_j \to T\}_{j = 1, \ldots, m}$ with each $U_j$ is
affine, $U_j \to T$ standard syntomic and $T = \bigcup f_j(U_j)$.
\end{definition}

\begin{definition}[Big syntomic site]
\label{topologies-definition-big-syntomic-site}
\source{stacks:03X1}
\href{https://stacks.math.columbia.edu/tag/03X1}{\texttt{Stacks tag 03X1}}
\group{topologies}

A {\it big syntomic site} is any site $\Sch_{syntomic}$ as in
Sites, Definition \ref{sites-definition-site} constructed as follows:
\begin{enumerate}
\item Choose any set of schemes $S_0$, and any set of syntomic coverings
$\text{Cov}_0$ among these schemes.
\item As underlying category take any category $\Sch_\alpha$
constructed as in Sets, Lemma \ref{sets-lemma-construct-category}
starting with the set $S_0$.
\item Choose any set of coverings as in
Sets, Lemma \ref{sets-lemma-coverings-site} starting with the
category $\Sch_\alpha$ and the class of syntomic coverings,
and the set $\text{Cov}_0$ chosen above.
\end{enumerate}
\end{definition}

See the remarks following Definition \ref{topologies-definition-big-zariski-site}
for motivation and explanation regarding the definition of big sites.

\medskip\noindent
Before we continue with the introduction of the big syntomic site of
a scheme $S$, let us point out that the topology on a big syntomic site
$\Sch_{syntomic}$ is in some sense induced from the syntomic topology
on the category of all schemes.

\begin{lemma}[Syntomic induced]
\label{topologies-lemma-syntomic-induced}
\source{stacks:03X2}
\href{https://stacks.math.columbia.edu/tag/03X2}{\texttt{Stacks tag 03X2}}
\group{topologies}

\uses{topologies-definition-big-syntomic-site}
Let $\Sch_{syntomic}$ be a big syntomic site as in
Definition \ref{topologies-definition-big-syntomic-site}.
Let $T \in \Ob(\Sch_{syntomic})$.
Let $\{T_i \to T\}_{i \in I}$ be an arbitrary syntomic covering of $T$.
\begin{enumerate}
\item There exists a covering $\{U_j \to T\}_{j \in J}$ of $T$ in the site
$\Sch_{syntomic}$ which refines $\{T_i \to T\}_{i \in I}$.
\item If $\{T_i \to T\}_{i \in I}$ is a standard syntomic covering, then
it is tautologically equivalent to a covering in $\Sch_{syntomic}$.
\item If $\{T_i \to T\}_{i \in I}$ is a Zariski covering, then
it is tautologically equivalent to a covering in $\Sch_{syntomic}$.
\end{enumerate}
\end{lemma}

\begin{proof}
\uses{topologies-lemma-syntomic}

For each $i$ choose an affine open covering $T_i = \bigcup_{j \in J_i} T_{ij}$
such that each $T_{ij}$ maps into an affine open subscheme of $T$. By
Lemma \ref{topologies-lemma-syntomic}
the refinement $\{T_{ij} \to T\}_{i \in I, j \in J_i}$ is a syntomic covering
of $T$ as well. Hence we may assume each $T_i$ is affine, and maps into
an affine open $W_i$ of $T$. Applying
Sets, Lemma \ref{sets-lemma-what-is-in-it}
we see that $W_i$ is isomorphic to an object of $\Sch_{syntomic}$.
But then $T_i$ as a finite type scheme over $W_i$
is isomorphic to an object $V_i$ of $\Sch_{syntomic}$ by a second
application of
Sets, Lemma \ref{sets-lemma-what-is-in-it}.
The covering $\{V_i \to T\}_{i \in I}$ refines $\{T_i \to T\}_{i \in I}$
(because they are isomorphic).
Moreover, $\{V_i \to T\}_{i \in I}$ is combinatorially equivalent to a
covering $\{U_j \to T\}_{j \in J}$ of $T$ in the site
$\Sch_{syntomic}$ by
Sets, Lemma \ref{sets-lemma-what-is-in-it}.
The covering $\{U_j \to T\}_{j \in J}$ is a covering as in (1).
In the situation of (2), (3) each of the
schemes $T_i$ is isomorphic to an object of $\Sch_{syntomic}$ by
Sets, Lemma \ref{sets-lemma-what-is-in-it},
and another application of
Sets, Lemma \ref{sets-lemma-coverings-site}
gives what we want.
\end{proof}

\begin{definition}[Big small syntomic]
\label{topologies-definition-big-small-syntomic}
\source{stacks:03X3}
\href{https://stacks.math.columbia.edu/tag/03X3}{\texttt{Stacks tag 03X3}}
\group{topologies}

Let $S$ be a scheme. Let $\Sch_{syntomic}$ be a big syntomic
site containing $S$.
\begin{enumerate}
\item The {\it big syntomic site of $S$}, denoted
$(\Sch/S)_{syntomic}$, is the site $\Sch_{syntomic}/S$
introduced in Sites, Section \ref{sites-section-localize}.
\item The {\it big affine syntomic site of $S$}, denoted
$(\textit{Aff}/S)_{syntomic}$, is the full subcategory of
$(\Sch/S)_{syntomic}$ whose objects are affine $U/S$.
A covering of $(\textit{Aff}/S)_{syntomic}$ is any covering
$\{U_i \to U\}$ of $(\Sch/S)_{syntomic}$ which is a
standard syntomic covering.
\end{enumerate}
\end{definition}

Next, we check that the big affine site defines the same
topos as the big site.

\begin{lemma}[Affine big site syntomic]
\label{topologies-lemma-affine-big-site-syntomic}
\source{stacks:06VD}
\href{https://stacks.math.columbia.edu/tag/06VD}{\texttt{Stacks tag 06VD}}
\group{topologies}

Let $S$ be a scheme. Let $\Sch_{syntomic}$ be a big syntomic
site containing $S$.
The functor
$(\textit{Aff}/S)_{syntomic} \to (\Sch/S)_{syntomic}$
is special cocontinuous and induces an equivalence of topoi from
$\Sh((\textit{Aff}/S)_{syntomic})$ to
$\Sh((\Sch/S)_{syntomic})$.
\end{lemma}

\begin{proof}
\uses{topologies-lemma-syntomic-affine}

The notion of a special cocontinuous functor is introduced in
Sites, Definition \ref{sites-definition-special-cocontinuous-functor}.
Thus we have to verify assumptions (1) -- (5) of
Sites, Lemma \ref{sites-lemma-equivalence}.
Denote the inclusion functor
$u : (\textit{Aff}/S)_{syntomic} \to (\Sch/S)_{syntomic}$.
Being cocontinuous just means that any syntomic covering of
$T/S$, $T$ affine, can be refined by a standard syntomic covering of $T$.
This is the content of
Lemma \ref{topologies-lemma-syntomic-affine}.
Hence (1) holds. We see $u$ is continuous simply because a standard
syntomic covering is a syntomic covering. Hence (2) holds.
Parts (3) and (4) follow immediately from the fact that $u$ is
fully faithful. And finally condition (5) follows from the
fact that every scheme has an affine open covering.
\end{proof}

To be continued...

\begin{lemma}[Morphism big syntomic]
\label{topologies-lemma-morphism-big-syntomic}
\source{stacks:04HD}
\href{https://stacks.math.columbia.edu/tag/04HD}{\texttt{Stacks tag 04HD}}
\group{topologies}

Let $\Sch_{syntomic}$ be a big syntomic site.
Let $f : T \to S$ be a morphism in $\Sch_{syntomic}$.
The functor
$$
u : (\Sch/T)_{syntomic} \longrightarrow (\Sch/S)_{syntomic},
\quad
V/T \longmapsto V/S
$$
is cocontinuous, and has a continuous right adjoint
$$
v : (\Sch/S)_{syntomic} \longrightarrow (\Sch/T)_{syntomic},
\quad
(U \to S) \longmapsto (U \times_S T \to T).
$$
They induce the same morphism of topoi
$$
f_{big} :
\Sh((\Sch/T)_{syntomic})
\longrightarrow
\Sh((\Sch/S)_{syntomic})
$$
We have $f_{big}^{-1}(\mathcal{G})(U/T) = \mathcal{G}(U/S)$.
We have $f_{big, *}(\mathcal{F})(U/S) = \mathcal{F}(U \times_S T/T)$.
Also, $f_{big}^{-1}$ has a left adjoint $f_{big!}$ which commutes with
fibre products and equalizers.
\end{lemma}

\begin{proof}
The functor $u$ is cocontinuous, continuous, and commutes with fibre products
and equalizers. Hence
Sites, Lemmas \ref{sites-lemma-when-shriek} and
\ref{sites-lemma-preserve-equalizers}
apply and we deduce the formula
for $f_{big}^{-1}$ and the existence of $f_{big!}$. Moreover,
the functor $v$ is a right adjoint because given $U/T$ and $V/S$
we have $\Mor_S(u(U), V) = \Mor_T(U, V \times_S T)$
as desired. Thus we may apply
Sites, Lemmas \ref{sites-lemma-have-functor-other-way} and
\ref{sites-lemma-have-functor-other-way-morphism} to get the
formula for $f_{big, *}$.
\end{proof}

\section{The fppf topology}
\label{topologies-section-fppf}
Let $S$ be a scheme. We would like to define the fppf-topology\footnote{
The letters fppf stand for ``fid\`element plat de pr\'esentation finie''.} on
the category of schemes over $S$. According to our general principle
we first introduce the notion of an fppf-covering.

\begin{definition}[Fppf covering]
\label{topologies-definition-fppf-covering}
\source{stacks:021M}
\href{https://stacks.math.columbia.edu/tag/021M}{\texttt{Stacks tag 021M}}
\group{topologies}

Let $T$ be a scheme. An {\it fppf covering of $T$} is a family
of morphisms $\{f_i : T_i \to T\}_{i \in I}$ of schemes
such that each $f_i$ is flat, locally of finite presentation and such
that $T = \bigcup f_i(T_i)$.
\end{definition}

\begin{lemma}[Zariski etale smooth syntomic fppf]
\label{topologies-lemma-zariski-etale-smooth-syntomic-fppf}
\source{stacks:021N}
\href{https://stacks.math.columbia.edu/tag/021N}{\texttt{Stacks tag 021N}}
\group{topologies}

Any syntomic covering is an fppf covering, and a fortiori,
any smooth, \'etale, or Zariski covering is an fppf covering.
\end{lemma}

\begin{proof}
\uses{topologies-lemma-zariski-etale-smooth-syntomic}

This is clear from the definitions, the fact that a syntomic morphism
is flat and locally of finite presentation, see
Morphisms, Lemmas
\ref{morphisms-lemma-syntomic-locally-finite-presentation} and
\ref{morphisms-lemma-syntomic-flat},
and
Lemma \ref{topologies-lemma-zariski-etale-smooth-syntomic}.
\end{proof}

Next, we show that this notion satisfies the conditions of
Sites, Definition \ref{sites-definition-site}.

\begin{lemma}[Fppf]
\label{topologies-lemma-fppf}
\source{stacks:021O}
\href{https://stacks.math.columbia.edu/tag/021O}{\texttt{Stacks tag 021O}}
\group{topologies}

Let $T$ be a scheme.
\begin{enumerate}
\item If $T' \to T$ is an isomorphism then $\{T' \to T\}$
is an fppf covering of $T$.
\item If $\{T_i \to T\}_{i\in I}$ is an fppf covering and for each
$i$ we have an fppf covering $\{T_{ij} \to T_i\}_{j\in J_i}$, then
$\{T_{ij} \to T\}_{i \in I, j\in J_i}$ is an fppf covering.
\item If $\{T_i \to T\}_{i\in I}$ is an fppf covering
and $T' \to T$ is a morphism of schemes then
$\{T' \times_T T_i \to T'\}_{i\in I}$ is an fppf covering.
\end{enumerate}
\end{lemma}

\begin{proof}
The first assertion is clear.
The second follows as the composition of flat morphisms is flat
(see Morphisms, Lemma \ref{morphisms-lemma-composition-flat})
and the composition of morphisms of finite presentation is
of finite presentation
(see Morphisms, Lemma \ref{morphisms-lemma-composition-finite-presentation}).
The third follows as the base change of a flat morphism is flat
(see Morphisms, Lemma \ref{morphisms-lemma-base-change-flat})
and the base change of a morphism of finite presentation is
of finite presentation
(see Morphisms, Lemma \ref{morphisms-lemma-base-change-finite-presentation}).
Moreover, the base change of a surjective family of morphisms is surjective
(proof omitted).
\end{proof}

\begin{lemma}[Fppf affine]
\label{topologies-lemma-fppf-affine}
\source{stacks:021P}
\href{https://stacks.math.columbia.edu/tag/021P}{\texttt{Stacks tag 021P}}
\group{topologies}

Let $T$ be an affine scheme.
Let $\{T_i \to T\}_{i \in I}$ be an fppf covering of $T$.
Then there exists an fppf covering
$\{U_j \to T\}_{j = 1, \ldots, m}$ which is a refinement
of $\{T_i \to T\}_{i \in I}$ such that each $U_j$ is an affine
scheme. Moreover, we may choose each $U_j$ to be open affine
in one of the $T_i$.
\end{lemma}

\begin{proof}
This follows directly from the definitions using that a
morphism which is flat and locally of finite presentation is open,
see Morphisms, Lemma \ref{morphisms-lemma-fppf-open}.
\end{proof}

Thus we define the corresponding standard coverings of affines as follows.

\begin{definition}[Standard fppf]
\label{topologies-definition-standard-fppf}
\source{stacks:021Q}
\href{https://stacks.math.columbia.edu/tag/021Q}{\texttt{Stacks tag 021Q}}
\group{topologies}

Let $T$ be an affine scheme. A {\it standard fppf covering}
of $T$ is a family $\{f_j : U_j \to T\}_{j = 1, \ldots, m}$
with each $U_j$ is affine, flat and of finite presentation over $T$
and $T = \bigcup f_j(U_j)$.
\end{definition}

\begin{definition}[Big fppf site]
\label{topologies-definition-big-fppf-site}
\source{stacks:021R}
\href{https://stacks.math.columbia.edu/tag/021R}{\texttt{Stacks tag 021R}}
\group{topologies}

A {\it big fppf site} is any site $\Sch_{fppf}$ as in
Sites, Definition \ref{sites-definition-site} constructed as follows:
\begin{enumerate}
\item Choose any set of schemes $S_0$, and any set of fppf coverings
$\text{Cov}_0$ among these schemes.
\item As underlying category take any category $\Sch_\alpha$
constructed as in Sets, Lemma \ref{sets-lemma-construct-category}
starting with the set $S_0$.
\item Choose any set of coverings as in
Sets, Lemma \ref{sets-lemma-coverings-site} starting with the
category $\Sch_\alpha$ and the class of fppf coverings,
and the set $\text{Cov}_0$ chosen above.
\end{enumerate}
\end{definition}

See the remarks following Definition \ref{topologies-definition-big-zariski-site}
for motivation and explanation regarding the definition of big sites.

\medskip\noindent
Before we continue with the introduction of the big fppf site of
a scheme $S$, let us point out that the topology on a big fppf site
$\Sch_{fppf}$ is in some sense induced from the fppf topology
on the category of all schemes.

\begin{lemma}[Fppf induced]
\label{topologies-lemma-fppf-induced}
\source{stacks:03WX}
\href{https://stacks.math.columbia.edu/tag/03WX}{\texttt{Stacks tag 03WX}}
\group{topologies}

\uses{topologies-definition-big-fppf-site}
Let $\Sch_{fppf}$ be a big fppf site as in
Definition \ref{topologies-definition-big-fppf-site}.
Let $T \in \Ob(\Sch_{fppf})$.
Let $\{T_i \to T\}_{i \in I}$ be an arbitrary fppf covering of $T$.
\begin{enumerate}
\item There exists a covering $\{U_j \to T\}_{j \in J}$ of $T$ in the site
$\Sch_{fppf}$ which refines $\{T_i \to T\}_{i \in I}$.
\item If $\{T_i \to T\}_{i \in I}$ is a standard fppf covering, then
it is tautologically equivalent to a covering of $\Sch_{fppf}$.
\item If $\{T_i \to T\}_{i \in I}$ is a Zariski covering, then
it is tautologically equivalent to a covering of $\Sch_{fppf}$.
\end{enumerate}
\end{lemma}

\begin{proof}
\uses{topologies-lemma-fppf}

For each $i$ choose an affine open covering $T_i = \bigcup_{j \in J_i} T_{ij}$
such that each $T_{ij}$ maps into an affine open subscheme of $T$. By
Lemma \ref{topologies-lemma-fppf}
the refinement $\{T_{ij} \to T\}_{i \in I, j \in J_i}$ is an fppf covering
of $T$ as well. Hence we may assume each $T_i$ is affine, and maps into
an affine open $W_i$ of $T$. Applying
Sets, Lemma \ref{sets-lemma-what-is-in-it}
we see that $W_i$ is isomorphic to an object of $\Sch_{fppf}$.
But then $T_i$ as a finite type scheme over $W_i$
is isomorphic to an object $V_i$ of $\Sch_{fppf}$ by a second
application of
Sets, Lemma \ref{sets-lemma-what-is-in-it}.
The covering $\{V_i \to T\}_{i \in I}$ refines $\{T_i \to T\}_{i \in I}$
(because they are isomorphic).
Moreover, $\{V_i \to T\}_{i \in I}$ is combinatorially equivalent to a
covering $\{U_j \to T\}_{j \in J}$ of $T$ in the site
$\Sch_{fppf}$ by
Sets, Lemma \ref{sets-lemma-what-is-in-it}.
The covering $\{U_j \to T\}_{j \in J}$ is a refinement as in (1).
In the situation of (2), (3) each of the
schemes $T_i$ is isomorphic to an object of $\Sch_{fppf}$ by
Sets, Lemma \ref{sets-lemma-what-is-in-it},
and another application of
Sets, Lemma \ref{sets-lemma-coverings-site}
gives what we want.
\end{proof}

\begin{definition}[Big small fppf]
\label{topologies-definition-big-small-fppf}
\source{stacks:021S}
\href{https://stacks.math.columbia.edu/tag/021S}{\texttt{Stacks tag 021S}}
\group{topologies}

Let $S$ be a scheme. Let $\Sch_{fppf}$ be a big fppf
site containing $S$.
\begin{enumerate}
\item The {\it big fppf site of $S$}, denoted
$(\Sch/S)_{fppf}$, is the site $\Sch_{fppf}/S$
introduced in Sites, Section \ref{sites-section-localize}.
\item The {\it big affine fppf site of $S$}, denoted
$(\textit{Aff}/S)_{fppf}$, is the full subcategory of
$(\Sch/S)_{fppf}$ whose objects are affine $U/S$.
A covering of $(\textit{Aff}/S)_{fppf}$ is any covering
$\{U_i \to U\}$ of $(\Sch/S)_{fppf}$ which is a
standard fppf covering.
\end{enumerate}
\end{definition}

It is not completely clear that
the big affine fppf site is a site. We check this now.

\begin{lemma}[Verify site fppf]
\label{topologies-lemma-verify-site-fppf}
\source{stacks:021T}
\href{https://stacks.math.columbia.edu/tag/021T}{\texttt{Stacks tag 021T}}
\group{topologies}

Let $S$ be a scheme. Let $\Sch_{fppf}$ be a big fppf
site containing $S$. Then $(\textit{Aff}/S)_{fppf}$ is a site.
\end{lemma}

\begin{proof}
\uses{topologies-lemma-verify-site-etale}

Let us show that $(\textit{Aff}/S)_{fppf}$ is a site.
Reasoning as in the proof of Lemma \ref{topologies-lemma-verify-site-etale}
it suffices to show that the collection
of standard fppf coverings of affines satisfies properties
(1), (2) and (3) of
Sites, Definition \ref{sites-definition-site}.
This is clear since for example, given a standard fppf
covering $\{T_i \to T\}_{i\in I}$ and for each
$i$ we have a standard fppf covering $\{T_{ij} \to T_i\}_{j\in J_i}$, then
$\{T_{ij} \to T\}_{i \in I, j\in J_i}$ is a standard fppf covering
because $\bigcup_{i\in I} J_i$ is finite and each $T_{ij}$ is affine.
\end{proof}

\begin{lemma}[Fibre products fppf]
\label{topologies-lemma-fibre-products-fppf}
\source{stacks:021U}
\href{https://stacks.math.columbia.edu/tag/021U}{\texttt{Stacks tag 021U}}
\group{topologies}

Let $S$ be a scheme. Let $\Sch_{fppf}$ be a big fppf
site containing $S$. The underlying categories of the sites
$\Sch_{fppf}$, $(\Sch/S)_{fppf}$,
and $(\textit{Aff}/S)_{fppf}$ have fibre products.
In each case the obvious functor into the category $\Sch$ of
all schemes commutes with taking fibre products. The category
$(\Sch/S)_{fppf}$ has a final object, namely $S/S$.
\end{lemma}

\begin{proof}
For $\Sch_{fppf}$ it is true by construction, see
Sets, Lemma \ref{sets-lemma-what-is-in-it}.
Suppose we have $U \to S$, $V \to U$, $W \to U$ morphisms
of schemes with $U, V, W \in \Ob(\Sch_{fppf})$.
The fibre product $V \times_U W$ in $\Sch_{fppf}$
is a fibre product in $\Sch$ and
is the fibre product of $V/S$ with $W/S$ over $U/S$ in
the category of all schemes over $S$, and hence also a
fibre product in $(\Sch/S)_{fppf}$.
This proves the result for $(\Sch/S)_{fppf}$.
If $U, V, W$ are affine, so is $V \times_U W$ and hence the
result for $(\textit{Aff}/S)_{fppf}$.
\end{proof}

Next, we check that the big affine site defines the same
topos as the big site.

\begin{lemma}[Affine big site fppf]
\label{topologies-lemma-affine-big-site-fppf}
\source{stacks:021V}
\href{https://stacks.math.columbia.edu/tag/021V}{\texttt{Stacks tag 021V}}
\group{topologies}

Let $S$ be a scheme. Let $\Sch_{fppf}$ be a big fppf
site containing $S$.
The functor $(\textit{Aff}/S)_{fppf} \to (\Sch/S)_{fppf}$
is cocontinuous and induces an equivalence of topoi from
$\Sh((\textit{Aff}/S)_{fppf})$ to
$\Sh((\Sch/S)_{fppf})$.
\end{lemma}

\begin{proof}
\uses{topologies-lemma-fppf-affine}

The notion of a special cocontinuous functor is introduced in
Sites, Definition \ref{sites-definition-special-cocontinuous-functor}.
Thus we have to verify assumptions (1) -- (5) of
Sites, Lemma \ref{sites-lemma-equivalence}.
Denote the inclusion functor
$u : (\textit{Aff}/S)_{fppf} \to (\Sch/S)_{fppf}$.
Being cocontinuous just means that any fppf covering of
$T/S$, $T$ affine, can be refined by a standard fppf covering of $T$.
This is the content of
Lemma \ref{topologies-lemma-fppf-affine}.
Hence (1) holds. We see $u$ is continuous simply because a standard
fppf covering is a fppf covering. Hence (2) holds.
Parts (3) and (4) follow immediately from the fact that $u$ is
fully faithful. And finally condition (5) follows from the
fact that every scheme has an affine open covering.
\end{proof}

Next, we establish some relationships between the topoi
associated to these sites.

\begin{lemma}[Morphism big fppf]
\label{topologies-lemma-morphism-big-fppf}
\source{stacks:021W}
\href{https://stacks.math.columbia.edu/tag/021W}{\texttt{Stacks tag 021W}}
\group{topologies}

Let $\Sch_{fppf}$ be a big fppf site.
Let $f : T \to S$ be a morphism in $\Sch_{fppf}$.
The functor
$$
u : (\Sch/T)_{fppf} \longrightarrow (\Sch/S)_{fppf},
\quad
V/T \longmapsto V/S
$$
is cocontinuous, and has a continuous right adjoint
$$
v : (\Sch/S)_{fppf} \longrightarrow (\Sch/T)_{fppf},
\quad
(U \to S) \longmapsto (U \times_S T \to T).
$$
They induce the same morphism of topoi
$$
f_{big} :
\Sh((\Sch/T)_{fppf})
\longrightarrow
\Sh((\Sch/S)_{fppf})
$$
We have $f_{big}^{-1}(\mathcal{G})(U/T) = \mathcal{G}(U/S)$.
We have $f_{big, *}(\mathcal{F})(U/S) = \mathcal{F}(U \times_S T/T)$.
Also, $f_{big}^{-1}$ has a left adjoint $f_{big!}$ which commutes with
fibre products and equalizers.
\end{lemma}

\begin{proof}
The functor $u$ is cocontinuous, continuous, and commutes with fibre products
and equalizers. Hence
Sites, Lemmas \ref{sites-lemma-when-shriek} and
\ref{sites-lemma-preserve-equalizers}
apply and we deduce the formula
for $f_{big}^{-1}$ and the existence of $f_{big!}$. Moreover,
the functor $v$ is a right adjoint because given $U/T$ and $V/S$
we have $\Mor_S(u(U), V) = \Mor_T(U, V \times_S T)$
as desired. Thus we may apply
Sites, Lemmas \ref{sites-lemma-have-functor-other-way} and
\ref{sites-lemma-have-functor-other-way-morphism} to get the
formula for $f_{big, *}$.
\end{proof}

\begin{lemma}[Composition fppf]
\label{topologies-lemma-composition-fppf}
\source{stacks:021X}
\href{https://stacks.math.columbia.edu/tag/021X}{\texttt{Stacks tag 021X}}
\group{topologies}

Given schemes $X$, $Y$, $Z$ in $(\Sch/S)_{fppf}$
and morphisms $f : X \to Y$, $g : Y \to Z$ we have
$g_{big} \circ f_{big} = (g \circ f)_{big}$.
\end{lemma}

\begin{proof}
\uses{topologies-lemma-morphism-big-fppf}

This follows from the simple description of pushforward
and pullback for the functors on the big sites from
Lemma \ref{topologies-lemma-morphism-big-fppf}.
\end{proof}

\section{The ph topology}
\label{topologies-section-ph}
In this section we define the ph topology. This is the topology
generated by Zariski coverings and proper surjective morphisms, see
Lemma \ref{topologies-lemma-characterize-sheaf}.

\medskip\noindent
We borrow our notation/terminology
from the paper \cite{ph} by Goodwillie and Lichtenbaum.
These authors show that if we restrict to the subcategory of
Noetherian schemes, then the ph topology is the same as the
``h topology'' as originally defined by Voevodsky: this is
the topology generated by Zariski open coverings and
finite type morphisms which are universally submersive.
They also show that the two topologies do not agree on
non-Noetherian schemes, see \cite[Example 4.5]{ph}.
We return to (our version of) the h topology in
More on Flatness, Section \ref{flat-section-h}.

\medskip\noindent
Before we can define the coverings in our topology we need to do
a bit of work.

\begin{definition}[Standard ph covering]
\label{topologies-definition-standard-ph-covering}
\source{stacks:0DBD}
\href{https://stacks.math.columbia.edu/tag/0DBD}{\texttt{Stacks tag 0DBD}}
\group{topologies}

Let $T$ be an affine scheme. A {\it standard ph covering} is a family
$\{f_j : U_j \to T\}_{j = 1, \ldots, m}$ constructed from a
proper surjective morphism $f : U \to T$ and an affine open covering
$U = \bigcup_{j = 1, \ldots, m} U_j$ by setting $f_j = f|_{U_j}$.
\end{definition}

It follows immediately from Chow's lemma that we can refine a
standard ph covering by a standard ph covering corresponding to
a surjective projective morphism.

\begin{lemma}[Base change standard ph]
\label{topologies-lemma-base-change-standard-ph}
\source{stacks:0DBE}
\href{https://stacks.math.columbia.edu/tag/0DBE}{\texttt{Stacks tag 0DBE}}
\group{topologies}

Let $\{f_j : U_j \to T\}_{j = 1, \ldots, m}$ be a standard ph covering.
Let $T' \to T$ be a morphism of affine schemes.
Then $\{U_j \times_T T' \to T'\}_{j = 1, \ldots, m}$ is a
standard ph covering.
\end{lemma}

\begin{proof}
\uses{topologies-definition-standard-ph-covering}

Let $f : U \to T$ be proper surjective and let an affine open covering
$U = \bigcup_{j = 1, \ldots, m} U_j$ be given as in
Definition \ref{topologies-definition-standard-ph-covering}.
Then $U \times_T T' \to T'$ is proper surjective
(Morphisms, Lemmas \ref{morphisms-lemma-base-change-surjective} and
\ref{morphisms-lemma-base-change-proper}).
Also, $U \times_T T' = \bigcup_{j = 1, \ldots, m} U_j \times_T T'$
is an affine open covering.
This concludes the proof.
\end{proof}

\begin{lemma}[Refine by standard ph]
\label{topologies-lemma-refine-by-standard-ph}
\source{stacks:0DBF}
\href{https://stacks.math.columbia.edu/tag/0DBF}{\texttt{Stacks tag 0DBF}}
\group{topologies}

Let $T$ be an affine scheme. Each of the following types of families
of maps with target $T$ has a refinement by a standard ph covering:
\begin{enumerate}
\item any Zariski open covering of $T$,
\item $\{W_{ji} \to T\}_{j = 1, \ldots, m, i = 1, \ldots n_j}$
where $\{W_{ji} \to U_j\}_{i = 1, \ldots, n_j}$
and $\{U_j \to T\}_{j = 1, \ldots, m}$ are standard ph coverings.
\end{enumerate}
\end{lemma}

\begin{proof}
\uses{topologies-definition-standard-ph-covering}

Part (1) follows from the fact that any Zariski open covering of $T$
can be refined by a finite affine open covering.

\medskip\noindent
Proof of (3). Choose $U \to T$ proper surjective and
$U = \bigcup_{j = 1, \ldots, m} U_j$ as in
Definition \ref{topologies-definition-standard-ph-covering}.
Choose $W_j \to U_j$ proper surjective and $W_j = \bigcup W_{ji}$
as in Definition \ref{topologies-definition-standard-ph-covering}.
By Chow's lemma (Limits, Lemma \ref{limits-lemma-chow-finite-type})
we can find $W'_j \to W_j$ proper surjective
and closed immersions $W'_j \to \mathbf{P}^{e_j}_{U_j}$.
Thus, after replacing $W_j$ by $W'_j$ and $W_j = \bigcup W_{ji}$
by a suitable affine open covering of $W'_j$, we may assume
there is a closed immersion $W_j \subset \mathbf{P}^{e_j}_{U_j}$
for all $j = 1, \ldots, m$.

\medskip\noindent
Let $\overline{W}_j \subset \mathbf{P}^{e_j}_U$ be the
scheme theoretic closure of $W_j$. Then $W_j \subset \overline{W}_j$
is an open subscheme; in fact $W_j$ is the inverse image of
$U_j \subset U$ under the morphism $\overline{W}_j \to U$.
(To see this use that $W_j \to \mathbf{P}^{e_j}_U$ is quasi-compact
and hence formation of the scheme theoretic image commutes
with restriction to opens, see
Morphisms, Section \ref{morphisms-section-scheme-theoretic-image}.)
Let $Z_j = U \setminus U_j$
with reduced induced closed subscheme structure.
Then
$$
V_j = \overline{W}_j \amalg Z_j \to U
$$
is proper surjective and the open subscheme $W_j \subset V_j$
is the inverse image of $U_j$. Hence for $v \in V_j$, $v \not \in W_j$
we can pick an affine open neighbourhood $v \in V_{j, v} \subset V_j$
which maps into $U_{j'}$ for some $1 \leq j' \leq m$.

\medskip\noindent
To finish the proof we consider the proper surjective morphism
$$
V = V_1 \times_U V_2 \times_U \ldots \times_U V_m
\longrightarrow U \longrightarrow T
$$
and the covering of $V$ by the affine opens
$$
V_{1, v_1} \times_U \ldots \times_U V_{j - 1, v_{j - 1}}
\times_U W_{j i} \times_U
V_{j + 1, v_{j + 1}} \times_U \ldots \times_U V_{m, v_m}
$$
These do indeed form a covering, because each point of $U$ is
in some $U_j$ and the inverse image of $U_j$ in $V$ is equal to
$V_1 \times \ldots \times V_{j - 1} \times W_j \times V_{j + 1} \times
\ldots \times V_m$. Observe that the morphism from
the affine open displayed above to $T$ factors through $W_{ji}$
thus we obtain a refinement. Finally, we only need a finite number
of these affine opens as $V$ is quasi-compact (as a scheme proper
over the affine scheme $T$).
\end{proof}

\begin{definition}[Ph covering]
\label{topologies-definition-ph-covering}
\source{stacks:0DBG}
\href{https://stacks.math.columbia.edu/tag/0DBG}{\texttt{Stacks tag 0DBG}}
\group{topologies}

Let $T$ be a scheme. A {\it ph covering of $T$} is a family
of morphisms $\{f_i : T_i \to T\}_{i \in I}$ of schemes such
that $f_i$ is locally of finite type and such that for every
affine open $U \subset T$ there exists a standard ph covering
$\{U_j \to U\}_{j = 1, \ldots, m}$ refining the family
$\{T_i \times_T U \to U\}_{i \in I}$.
\end{definition}

A standard ph covering is a ph covering by
Lemma \ref{topologies-lemma-base-change-standard-ph}.

\begin{lemma}[Zariski ph]
\label{topologies-lemma-zariski-ph}
\source{stacks:0DBH}
\href{https://stacks.math.columbia.edu/tag/0DBH}{\texttt{Stacks tag 0DBH}}
\group{topologies}

A Zariski covering is a ph covering\footnote{We will see
in More on Morphisms, Lemma \ref{more-morphisms-lemma-fppf-ph} that
fppf coverings (and hence syntomic, smooth, or \'etale coverings)
are ph coverings as well.}.
\end{lemma}

\begin{proof}
\uses{topologies-lemma-refine-by-standard-ph}

This is true because a Zariski covering of an affine scheme
can be refined by a standard ph covering by
Lemma \ref{topologies-lemma-refine-by-standard-ph}.
\end{proof}

\begin{lemma}[Surjective proper ph]
\label{topologies-lemma-surjective-proper-ph}
\source{stacks:0DES}
\href{https://stacks.math.columbia.edu/tag/0DES}{\texttt{Stacks tag 0DES}}
\group{topologies}

Let $f : Y \to X$ be a surjective proper morphism of schemes.
Then $\{Y \to X\}$ is a ph covering.
\end{lemma}

\begin{proof}
Omitted.
\end{proof}

\begin{lemma}[Refine by ph]
\label{topologies-lemma-refine-by-ph}
\source{stacks:0ET9}
\href{https://stacks.math.columbia.edu/tag/0ET9}{\texttt{Stacks tag 0ET9}}
\group{topologies}

Let $T$ be a scheme. Let $\{f_i : T_i \to T\}_{i \in I}$ be a family
of morphisms such that $f_i$ is locally of finite type for all $i$.
The following are equivalent
\begin{enumerate}
\item $\{T_i \to T\}_{i \in I}$ is a ph covering,
\item there is a ph covering which refines $\{T_i \to T\}_{i \in I}$, and
\item $\{\coprod_{i \in I} T_i \to T\}$ is a ph covering.
\end{enumerate}
\end{lemma}

\begin{proof}
\uses{topologies-definition-ph-covering}

The equivalence of (1) and (2) follows immediately from
Definition \ref{topologies-definition-ph-covering}
and the fact that a refinement of a refinement is a refinement.
Because of the equivalence of (1) and (2) and since
$\{T_i \to T\}_{i \in I}$ refines $\{\coprod_{i \in I} T_i \to T\}$
we see that (1) implies (3). Finally, assume (3) holds.
Let $U \subset T$ be an affine open and let
$\{U_j \to U\}_{j = 1, \ldots, m}$ be a standard ph covering
which refines $\{U \times_T \coprod_{i \in I} T_i \to U\}$.
This means that for each $j$ we have a morphism
$$
h_j :
U_j
\longrightarrow
U \times_T \coprod\nolimits_{i \in I} T_i =
\coprod\nolimits_{i \in I} U \times_T T_i
$$
over $U$. Since $U_j$ is quasi-compact, we get
disjoint union decompositions $U_j = \coprod_{i \in I} U_{j, i}$
by open and closed subschemes almost all of which are empty
such that $h_j|_{U_{j, i}}$ maps $U_{j, i}$ into $U \times_T T_i$.
It follows that
$$
\{U_{j, i} \to U\}_{j = 1, \ldots, m,\ i \in I,\ U_{j, i} \not = \emptyset}
$$
is a standard ph covering (small detail omitted) refining
$\{U \times_T T_i \to U\}_{i \in I}$. Thus (1) holds.
\end{proof}

Next, we show that this notion satisfies the conditions of
Sites, Definition \ref{sites-definition-site}.

\begin{lemma}[Ph]
\label{topologies-lemma-ph}
\source{stacks:0DBI}
\href{https://stacks.math.columbia.edu/tag/0DBI}{\texttt{Stacks tag 0DBI}}
\group{topologies}

Let $T$ be a scheme.
\begin{enumerate}
\item If $T' \to T$ is an isomorphism then $\{T' \to T\}$
is a ph covering of $T$.
\item If $\{T_i \to T\}_{i\in I}$ is a ph covering and for each
$i$ we have a ph covering $\{T_{ij} \to T_i\}_{j\in J_i}$, then
$\{T_{ij} \to T\}_{i \in I, j\in J_i}$ is a ph covering.
\item If $\{T_i \to T\}_{i\in I}$ is a ph covering
and $T' \to T$ is a morphism of schemes then
$\{T' \times_T T_i \to T'\}_{i\in I}$ is a ph covering.
\end{enumerate}
\end{lemma}

\begin{proof}
\uses{topologies-lemma-base-change-standard-ph, topologies-lemma-refine-by-standard-ph}

Assertion (1) is clear.

\medskip\noindent
Proof of (3). The base change $T_i \times_T T' \to T'$ is
locally of finite type by
Morphisms, Lemma \ref{morphisms-lemma-base-change-finite-type}.
hence we only need to check the condition on affine opens.
Let $U' \subset T'$ be an affine open subscheme.
Since $U'$ is quasi-compact we can find a finite affine open
covering $U' = U'_1 \cup \ldots \cup U'$ such that
$U'_j \to T$ maps into an affine open $U_j \subset T$.
Choose a standard ph covering $\{U_{jl} \to U_j\}_{l = 1, \ldots, n_j}$
refining $\{T_i \times_T U_j \to U_j\}$.
By Lemma \ref{topologies-lemma-base-change-standard-ph}
the base change $\{U_{jl} \times_{U_j} U'_j \to U'_j\}$
is a standard ph covering. Note that $\{U'_j \to U'\}$
is a standard ph covering as well.
By Lemma \ref{topologies-lemma-refine-by-standard-ph}
the family $\{U_{jl} \times_{U_j} U'_j \to U'\}$
can be refined by a standard ph covering. Since
$\{U_{jl} \times_{U_j} U'_j \to U'\}$ refines $\{T_i \times_T U' \to U'\}$
we conclude.

\medskip\noindent
Proof of (2). Composition preserves being locally of finite type,
see Morphisms, Lemma \ref{morphisms-lemma-composition-finite-type}.
Hence we only need to check the condition on affine opens.
Let $U \subset T$ be affine open. First we pick a standard ph covering
$\{U_k \to U\}_{k = 1, \ldots, m}$ refining
$\{T_i \times_T U \to U\}$.
Say the refinement is given by morphisms $U_k \to T_{i_k}$ over $T$.
Then
$$
\{T_{i_kj} \times_{T_{i_k}} U_k \to U_k\}_{j \in J_{i_k}}
$$
is a ph covering by part (3). As $U_k$ is affine,
we can find a standard ph covering
$\{U_{ka} \to U_k\}_{a = 1, \ldots, b_k}$ refining this family.
Then we apply Lemma \ref{topologies-lemma-refine-by-standard-ph}
to see that $\{U_{ka} \to U\}$ can be refined by a
standard ph covering. Since $\{U_{ka} \to U\}$
refines $\{T_{ij} \times_T U \to U\}$ this finishes the proof.
\end{proof}

\begin{definition}[Big ph site]
\label{topologies-definition-big-ph-site}
\source{stacks:0DBJ}
\href{https://stacks.math.columbia.edu/tag/0DBJ}{\texttt{Stacks tag 0DBJ}}
\group{topologies}

A {\it big ph site} is any site $\Sch_{ph}$ as in
Sites, Definition \ref{sites-definition-site} constructed as follows:
\begin{enumerate}
\item Choose any set of schemes $S_0$, and any set of ph coverings
$\text{Cov}_0$ among these schemes.
\item As underlying category take any category $\Sch_\alpha$
constructed as in Sets, Lemma \ref{sets-lemma-construct-category}
starting with the set $S_0$.
\item Choose any set of coverings as in
Sets, Lemma \ref{sets-lemma-coverings-site} starting with the
category $\Sch_\alpha$ and the class of ph coverings,
and the set $\text{Cov}_0$ chosen above.
\end{enumerate}
\end{definition}

See the remarks following Definition \ref{topologies-definition-big-zariski-site}
for motivation and explanation regarding the definition of big sites.

\medskip\noindent
Before we continue with the introduction of the big ph site of
a scheme $S$, let us point out that the topology on a big ph site
$\Sch_{ph}$ is in some sense induced from the ph topology
on the category of all schemes.

\begin{lemma}[Ph induced]
\label{topologies-lemma-ph-induced}
\source{stacks:0DBK}
\href{https://stacks.math.columbia.edu/tag/0DBK}{\texttt{Stacks tag 0DBK}}
\group{topologies}

\uses{topologies-definition-big-ph-site}
Let $\Sch_{ph}$ be a big ph site as in
Definition \ref{topologies-definition-big-ph-site}.
Let $T \in \Ob(\Sch_{ph})$.
Let $\{T_i \to T\}_{i \in I}$ be an arbitrary ph covering of $T$.
\begin{enumerate}
\item There exists a covering $\{U_j \to T\}_{j \in J}$ of $T$ in the site
$\Sch_{ph}$ which refines $\{T_i \to T\}_{i \in I}$.
\item If $\{T_i \to T\}_{i \in I}$ is a standard ph covering, then
it is tautologically equivalent to a covering of $\Sch_{ph}$.
\item If $\{T_i \to T\}_{i \in I}$ is a Zariski covering, then
it is tautologically equivalent to a covering of $\Sch_{ph}$.
\end{enumerate}
\end{lemma}

\begin{proof}
\uses{topologies-lemma-zariski-ph, topologies-lemma-ph}

For each $i$ choose an affine open covering $T_i = \bigcup_{j \in J_i} T_{ij}$
such that each $T_{ij}$ maps into an affine open subscheme of $T$. By
Lemmas \ref{topologies-lemma-zariski-ph} and \ref{topologies-lemma-ph}
the refinement $\{T_{ij} \to T\}_{i \in I, j \in J_i}$ is a ph covering
of $T$ as well. Hence we may assume each $T_i$ is affine, and maps into
an affine open $W_i$ of $T$. Applying
Sets, Lemma \ref{sets-lemma-what-is-in-it}
we see that $W_i$ is isomorphic to an object of $\Sch_{ph}$.
But then $T_i$ as a finite type scheme over $W_i$
is isomorphic to an object $V_i$ of $\Sch_{ph}$ by a second
application of
Sets, Lemma \ref{sets-lemma-what-is-in-it}.
The covering $\{V_i \to T\}_{i \in I}$ refines $\{T_i \to T\}_{i \in I}$
(because they are isomorphic).
Moreover, $\{V_i \to T\}_{i \in I}$ is combinatorially equivalent to a
covering $\{U_j \to T\}_{j \in J}$ of $T$ in the site
$\Sch_{ph}$ by
Sets, Lemma \ref{sets-lemma-what-is-in-it}.
The covering $\{U_j \to T\}_{j \in J}$ is a refinement as in (1).
In the situation of (2), (3) each of the
schemes $T_i$ is isomorphic to an object of $\Sch_{ph}$ by
Sets, Lemma \ref{sets-lemma-what-is-in-it},
and another application of
Sets, Lemma \ref{sets-lemma-coverings-site}
gives what we want.
\end{proof}

\begin{definition}[Big small ph]
\label{topologies-definition-big-small-ph}
\source{stacks:0DBL}
\href{https://stacks.math.columbia.edu/tag/0DBL}{\texttt{Stacks tag 0DBL}}
\group{topologies}

Let $S$ be a scheme. Let $\Sch_{ph}$ be a big ph site containing $S$.
\begin{enumerate}
\item The {\it big ph site of $S$}, denoted
$(\Sch/S)_{ph}$, is the site $\Sch_{ph}/S$
introduced in Sites, Section \ref{sites-section-localize}.
\item The {\it big affine ph site of $S$}, denoted
$(\textit{Aff}/S)_{ph}$, is the full subcategory of
$(\Sch/S)_{ph}$ whose objects are affine $U/S$.
A covering of $(\textit{Aff}/S)_{ph}$ is any finite covering
$\{U_i \to U\}$ of $(\Sch/S)_{ph}$ with $U_i$ and $U$ affine.
\end{enumerate}
\end{definition}

Observe that the coverings in $(\textit{Aff}/S)_{ph}$ are {\it not}
given by standard ph coverings. The reason is simply that this would
fail the second axiom of Sites, Definition \ref{sites-definition-site}.
Rather, the coverings in $(\textit{Aff}/S)_{ph}$ are those finite families
$\{U_i \to U\}$ of finite type morphisms between affine objects
of $(\Sch/S)_{ph}$ which can be refined by a standard ph covering.
We explicitly state and prove that the big affine ph site is a site.

\begin{lemma}[Verify site ph]
\label{topologies-lemma-verify-site-ph}
\source{stacks:0DBM}
\href{https://stacks.math.columbia.edu/tag/0DBM}{\texttt{Stacks tag 0DBM}}
\group{topologies}

Let $S$ be a scheme. Let $\Sch_{ph}$ be a big ph
site containing $S$. Then $(\textit{Aff}/S)_{ph}$ is a site.
\end{lemma}

\begin{proof}
\uses{topologies-lemma-verify-site-etale, topologies-lemma-ph}

Reasoning as in the proof of Lemma \ref{topologies-lemma-verify-site-etale}
it suffices to show that the collection of finite ph coverings
$\{U_i \to U\}$ with $U$, $U_i$ affine satisfies properties
(1), (2) and (3) of
Sites, Definition \ref{sites-definition-site}.
This is clear since for example, given a finite ph covering
$\{T_i \to T\}_{i\in I}$ with $T_i, T$ affine, and for each
$i$ a finite ph covering $\{T_{ij} \to T_i\}_{j\in J_i}$ with $T_{ij}$
affine , then $\{T_{ij} \to T\}_{i \in I, j\in J_i}$ is a ph covering
(Lemma \ref{topologies-lemma-ph}), $\bigcup_{i\in I} J_i$ is finite and
each $T_{ij}$ is affine.
\end{proof}

\begin{lemma}[Fibre products ph]
\label{topologies-lemma-fibre-products-ph}
\source{stacks:0DBN}
\href{https://stacks.math.columbia.edu/tag/0DBN}{\texttt{Stacks tag 0DBN}}
\group{topologies}

Let $S$ be a scheme. Let $\Sch_{ph}$ be a big ph
site containing $S$. The underlying categories of the sites
$\Sch_{ph}$, $(\Sch/S)_{ph}$, and $(\textit{Aff}/S)_{ph}$ have fibre products.
In each case the obvious functor into the category $\Sch$ of
all schemes commutes with taking fibre products. The category
$(\Sch/S)_{ph}$ has a final object, namely $S/S$.
\end{lemma}

\begin{proof}
For $\Sch_{ph}$ it is true by construction, see
Sets, Lemma \ref{sets-lemma-what-is-in-it}.
Suppose we have $U \to S$, $V \to U$, $W \to U$ morphisms
of schemes with $U, V, W \in \Ob(\Sch_{ph})$.
The fibre product $V \times_U W$ in $\Sch_{ph}$
is a fibre product in $\Sch$ and
is the fibre product of $V/S$ with $W/S$ over $U/S$ in
the category of all schemes over $S$, and hence also a
fibre product in $(\Sch/S)_{ph}$.
This proves the result for $(\Sch/S)_{ph}$.
If $U, V, W$ are affine, so is $V \times_U W$ and hence the
result for $(\textit{Aff}/S)_{ph}$.
\end{proof}

Next, we check that the big affine site defines the same
topos as the big site.

\begin{lemma}[Affine big site ph]
\label{topologies-lemma-affine-big-site-ph}
\source{stacks:0DBP}
\href{https://stacks.math.columbia.edu/tag/0DBP}{\texttt{Stacks tag 0DBP}}
\group{topologies}

Let $S$ be a scheme. Let $\Sch_{ph}$ be a big ph
site containing $S$.
The functor $(\textit{Aff}/S)_{ph} \to (\Sch/S)_{ph}$
is cocontinuous and induces an equivalence of topoi from
$\Sh((\textit{Aff}/S)_{ph})$ to
$\Sh((\Sch/S)_{ph})$.
\end{lemma}

\begin{proof}
The notion of a special cocontinuous functor is introduced in
Sites, Definition \ref{sites-definition-special-cocontinuous-functor}.
Thus we have to verify assumptions (1) -- (5) of
Sites, Lemma \ref{sites-lemma-equivalence}.
Denote the inclusion functor
$u : (\textit{Aff}/S)_{ph} \to (\Sch/S)_{ph}$.
Being cocontinuous follows because any ph covering of
$T/S$, $T$ affine, can be refined by a standard ph covering of $T$
by definition. Hence (1) holds. We see $u$ is continuous simply
because a finite ph covering of an affine by affines is a ph covering.
Hence (2) holds.
Parts (3) and (4) follow immediately from the fact that $u$ is
fully faithful. And finally condition (5) follows from the
fact that every scheme has an affine open covering (which is
a ph covering).
\end{proof}

\begin{lemma}[Characterize sheaf]
\label{topologies-lemma-characterize-sheaf}
\source{stacks:0DBQ}
\href{https://stacks.math.columbia.edu/tag/0DBQ}{\texttt{Stacks tag 0DBQ}}
\group{topologies}

Let $\mathcal{F}$ be a presheaf on $(\Sch/S)_{ph}$.
Then $\mathcal{F}$ is a sheaf if and only if
\begin{enumerate}
\item $\mathcal{F}$ satisfies the sheaf condition for
Zariski coverings, and
\item if $f : V \to U$ is proper surjective, then
$\mathcal{F}(U)$ maps bijectively to the equalizer
of the two maps $\mathcal{F}(V) \to \mathcal{F}(V \times_U V)$.
\end{enumerate}
Moreover, in the presence of (1) property (2) is equivalent to property
\begin{enumerate}
\item[(2')] the sheaf property for $\{V \to U\}$ as in (2) with $U$ affine.
\end{enumerate}
\end{lemma}

\begin{proof}
We will show that if (1) and (2) hold, then $\mathcal{F}$ is sheaf.
Let $\{T_i \to T\}$ be a ph covering, i.e., a covering in $(\Sch/S)_{ph}$.
We will verify the sheaf condition for this covering.
Let $s_i \in \mathcal{F}(T_i)$ be sections which restrict to the same
section over $T_i \times_T T_{i'}$. We will show that there exists a
unique section $s \in \mathcal{F}(T)$ restricting to $s_i$ over $T_i$.
Let $T = \bigcup U_j$ be an affine open covering.
By property (1) it suffices to produce sections $s_j \in \mathcal{F}(U_j)$
which agree on $U_j \cap U_{j'}$ in order to produce $s$.
Consider the ph coverings $\{T_i \times_T U_j \to U_j\}$.
Then $s_{ji} = s_i|_{T_i \times_T U_j}$ are sections agreeing
over $(T_i \times_T U_j) \times_{U_j} (T_{i'} \times_T U_j)$.
Choose a proper surjective morphism $V_j \to U_j$ and a finite affine
open covering $V_j = \bigcup V_{jk}$ such that the standard ph covering
$\{V_{jk} \to U_j\}$ refines $\{T_i \times_T U_j \to U_j\}$.
If $s_{jk} \in \mathcal{F}(V_{jk})$
denotes the pullback of $s_{ji}$ to $V_{jk}$ by the
implied morphisms, then we find that $s_{jk}$ glue to a section
$s'_j \in \mathcal{F}(V_j)$. Using the agreement on overlaps
once more, we find that $s'_j$ is in the equalizer of the two
maps $\mathcal{F}(V_j) \to \mathcal{F}(V_j \times_{U_j} V_j)$.
Hence by (2) we find that $s'_j$ comes from a unique section
$s_j \in \mathcal{F}(U_j)$. We omit the verification that these
sections $s_j$ have all the desired properties.

\medskip\noindent
Proof of the equivalence of (2) and (2') in the presence of (1).
Suppose $V \to U$ is a morphism of $(\Sch/S)_{ph}$ which is
proper and surjective. Choose an
affine open covering $U = \bigcup U_i$ and set $V_i = V \times_U U_i$.
Then we see that $\mathcal{F}(U) \to \mathcal{F}(V)$
is injective because we know $\mathcal{F}(U_i) \to \mathcal{F}(V_i)$
is injective by (2') and we know $\mathcal{F}(U) \to \prod \mathcal{F}(U_i)$
is injective by (1). Finally, suppose that we are given an
$t \in \mathcal{F}(V)$ in the equalizer of the two maps
$\mathcal{F}(V) \to \mathcal{F}(V \times_U V)$.
Then $t|_{V_i}$ is in the equalizer of the two maps
$\mathcal{F}(V_i) \to \mathcal{F}(V_i \times_{U_i} V_i)$
for all $i$. Hence we obtain a unique section $s_i \in \mathcal{F}(U_i)$
mapping to $t|_{V_i}$ for all $i$ by (2').
We omit the verification that $s_i|_{U_i \cap U_j} = s_j|_{U_i \cap U_j}$
for all $i, j$; this uses the uniqueness property just shown.
By the sheaf property for the covering $U = \bigcup U_i$ we obtain
a section $s \in \mathcal{F}(U)$. We omit the proof that $s$
maps to $t$ in $\mathcal{F}(V)$.
\end{proof}

Next, we establish some relationships between the topoi
associated to these sites.

\begin{lemma}[Morphism big ph]
\label{topologies-lemma-morphism-big-ph}
\source{stacks:0DBR}
\href{https://stacks.math.columbia.edu/tag/0DBR}{\texttt{Stacks tag 0DBR}}
\group{topologies}

Let $\Sch_{ph}$ be a big ph site.
Let $f : T \to S$ be a morphism in $\Sch_{ph}$.
The functor
$$
u : (\Sch/T)_{ph} \longrightarrow (\Sch/S)_{ph},
\quad
V/T \longmapsto V/S
$$
is cocontinuous, and has a continuous right adjoint
$$
v : (\Sch/S)_{ph} \longrightarrow (\Sch/T)_{ph},
\quad
(U \to S) \longmapsto (U \times_S T \to T).
$$
They induce the same morphism of topoi
$$
f_{big} :
\Sh((\Sch/T)_{ph})
\longrightarrow
\Sh((\Sch/S)_{ph})
$$
We have $f_{big}^{-1}(\mathcal{G})(U/T) = \mathcal{G}(U/S)$.
We have $f_{big, *}(\mathcal{F})(U/S) = \mathcal{F}(U \times_S T/T)$.
Also, $f_{big}^{-1}$ has a left adjoint $f_{big!}$ which commutes with
fibre products and equalizers.
\end{lemma}

\begin{proof}
The functor $u$ is cocontinuous, continuous, and commutes with fibre products
and equalizers. Hence
Sites, Lemmas \ref{sites-lemma-when-shriek} and
\ref{sites-lemma-preserve-equalizers}
apply and we deduce the formula
for $f_{big}^{-1}$ and the existence of $f_{big!}$. Moreover,
the functor $v$ is a right adjoint because given $U/T$ and $V/S$
we have $\Mor_S(u(U), V) = \Mor_T(U, V \times_S T)$
as desired. Thus we may apply
Sites, Lemmas \ref{sites-lemma-have-functor-other-way} and
\ref{sites-lemma-have-functor-other-way-morphism} to get the
formula for $f_{big, *}$.
\end{proof}

\begin{lemma}[Composition ph]
\label{topologies-lemma-composition-ph}
\source{stacks:0DBS}
\href{https://stacks.math.columbia.edu/tag/0DBS}{\texttt{Stacks tag 0DBS}}
\group{topologies}

Given schemes $X$, $Y$, $Y$ in $(\Sch/S)_{ph}$
and morphisms $f : X \to Y$, $g : Y \to Z$ we have
$g_{big} \circ f_{big} = (g \circ f)_{big}$.
\end{lemma}

\begin{proof}
\uses{topologies-lemma-morphism-big-ph}

This follows from the simple description of pushforward
and pullback for the functors on the big sites from
Lemma \ref{topologies-lemma-morphism-big-ph}.
\end{proof}

\section{The fpqc topology}
\label{topologies-section-fpqc}
\begin{definition}[Fpqc covering]
\label{topologies-definition-fpqc-covering}
\source{stacks:022B}
\href{https://stacks.math.columbia.edu/tag/022B}{\texttt{Stacks tag 022B}}
\group{topologies}

Let $T$ be a scheme. An {\it fpqc covering of $T$} is a family
of morphisms $\{f_i : T_i \to T\}_{i \in I}$ of schemes
such that each $f_i$ is flat and such that for every affine open
$U \subset T$ there exists $n \geq 0$, a map
$a : \{1, \ldots, n\} \to I$ and affine opens
$V_j \subset T_{a(j)}$, $j = 1, \ldots, n$
with $\bigcup_{j = 1}^n f_{a(j)}(V_j) = U$.
\end{definition}

To be sure this condition implies that $T = \bigcup f_i(T_i)$.
It is slightly harder to recognize an fpqc covering, hence we provide
some lemmas to do so.

\begin{lemma}[Recognize fpqc covering]
\label{topologies-lemma-recognize-fpqc-covering}
\source{stacks:03L7}
\href{https://stacks.math.columbia.edu/tag/03L7}{\texttt{Stacks tag 03L7}}
\group{topologies}

Let $T$ be a scheme. Let $\{f_i : T_i \to T\}_{i \in I}$ be a family of
morphisms of schemes with target $T$. The following are equivalent
\begin{enumerate}
\item $\{f_i : T_i \to T\}_{i \in I}$ is an fpqc covering,
\item each $f_i$ is flat and for every affine open $U \subset T$
there exist quasi-compact opens
$U_i \subset T_i$ which are almost all empty,
such that $U = \bigcup f_i(U_i)$,
\item each $f_i$ is flat and there exists an affine open covering
$T = \bigcup_{\alpha \in A} U_\alpha$ and for each $\alpha \in A$
there exist $i_{\alpha, 1}, \ldots, i_{\alpha, n(\alpha)} \in I$
and quasi-compact opens $U_{\alpha, j} \subset T_{i_{\alpha, j}}$ such that
$U_\alpha =
\bigcup_{j = 1, \ldots, n(\alpha)} f_{i_{\alpha, j}}(U_{\alpha, j})$.
\end{enumerate}
If $T$ is quasi-separated, these are also equivalent to
\begin{enumerate}
\item[(4)] each $f_i$ is flat, and for every $t \in T$ there exist
$i_1, \ldots, i_n \in I$ and quasi-compact opens $U_j \subset T_{i_j}$
such that $\bigcup_{j = 1, \ldots, n} f_{i_j}(U_j)$ is a
(not necessarily open) neighbourhood of $t$ in $T$.
\end{enumerate}
\end{lemma}

\begin{proof}
We omit the proof of the equivalence of (1), (2), and (3).
From now on assume $T$ is quasi-separated.
We prove (4) implies (2). Let $U \subset T$ be an affine open.
To prove (2) it suffices to show that for every $t \in U$ there exist
finitely many quasi-compact opens $U_j \subset T_{i_j}$ such that
$f_{i_j}(U_j) \subset U$ and such that $\bigcup f_{i_j}(U_j)$
is a neighbourhood of $t$ in $U$. By assumption there do exist
finitely many quasi-compact opens $U'_j \subset T_{i_j}$ such that
such that $\bigcup f_{i_j}(U'_j)$ is a neighbourhood of $t$ in $T$.
Since $T$ is quasi-separated we see that
$U_j = U'_j \cap f_{i_j}^{-1}(U) = (f_{i_j}|_{U'_j})^{-1}(U)$
is quasi-compact by
Schemes, Lemma \ref{schemes-lemma-quasi-compact-permanence}.
Thus (2) holds. Since it is clear that (2) implies
(4) the proof is finished.
\end{proof}

\begin{example}[Not fpqc]
\label{topologies-example-not-fpqc}
\source{stacks:0H7E}
\href{https://stacks.math.columbia.edu/tag/0H7E}{\texttt{Stacks tag 0H7E}}
\group{topologies}

\uses{topologies-lemma-recognize-fpqc-covering}
Let $X = X_1 \cup X_2$ where $X_1 = X_2 = \Spec(k[t_1,t_2,\dots])$
be the non-quasi-separated scheme from
Schemes, Example \ref{schemes-example-not-quasi-separated}, and
$Y = X_1 \amalg \Spec(\mathcal{O}_{X_2, 0})$.
Then the obvious morphism $f : Y \to X$
is flat and surjective, and (since $Y$ is quasi-compact)
trivially satisfies (4) of Lemma \ref{topologies-lemma-recognize-fpqc-covering}.
But $\{f : Y \to X\}$ isn't an fpqc covering. Namely, we claim there
is no quasi-compact open $W$ of $Y$ whose image by $f$ is $X_2$.
Namely, $f^{-1}(X_2) = X_1 \setminus \{0\} \amalg \Spec(\mathcal{O}_{X_2, 0})$.
Hence if $W$ exists, then $W = U \amalg \Spec(\mathcal{O}_{X_2, 0})$
for some quasi-compact open $U$ of $X_1 \setminus \{0\} = X_2 \setminus \{0\}$.
But then we may write $U = D(f_1) \cup \ldots \cup D(f_n)$ for
some $f_i \in k[t_1, t_2, \ldots]$ vanishing at $0$.
Suppose that $f_1, \ldots, f_n$ use only the variables
$x_1, \ldots, x_m$. Then $p = (0, \ldots, 0, 1, 0, \ldots)$ with $1$
in the $(m + 1)$st spot is a closed point of $X_2$ not in the image
of $W$.
\end{example}

\begin{lemma}[Disjoint union is fpqc covering]
\label{topologies-lemma-disjoint-union-is-fpqc-covering}
\source{stacks:040I}
\href{https://stacks.math.columbia.edu/tag/040I}{\texttt{Stacks tag 040I}}
\group{topologies}

Let $T$ be a scheme. Let $\{f_i : T_i \to T\}_{i \in I}$ be a family of
morphisms of schemes with target $T$. The following are equivalent
\begin{enumerate}
\item $\{f_i : T_i \to T\}_{i \in I}$ is an fpqc covering, and
\item setting $T' = \coprod_{i \in I} T_i$, and $f = \coprod_{i \in I} f_i$
the family $\{f : T' \to T\}$ is an fpqc covering.
\end{enumerate}
\end{lemma}

\begin{proof}
\uses{topologies-lemma-recognize-fpqc-covering}

Suppose that $U \subset T$ is an affine open. If (1) holds, then we find
$i_1, \ldots, i_n \in I$ and affine opens $U_j \subset T_{i_j}$ such that
$U = \bigcup_{j = 1, \ldots, n} f_{i_j}(U_j)$. Then
$U_1 \amalg \ldots \amalg U_n \subset T'$ is a quasi-compact open surjecting
onto $U$. Thus $\{f : T' \to T\}$ is an fpqc covering by
Lemma \ref{topologies-lemma-recognize-fpqc-covering}.
Conversely, if (2) holds then there exists a quasi-compact open
$U' \subset T'$ with $U = f(U')$. Then $U_j = U' \cap T_j$ is quasi-compact
open in $T_j$ and empty for almost all $j$. By
Lemma \ref{topologies-lemma-recognize-fpqc-covering} we see that (1) holds.
\end{proof}

\begin{lemma}[Family flat dominated covering]
\label{topologies-lemma-family-flat-dominated-covering}
\source{stacks:03L8}
\href{https://stacks.math.columbia.edu/tag/03L8}{\texttt{Stacks tag 03L8}}
\group{topologies}

Let $T$ be a scheme. Let $\{f_i : T_i \to T\}_{i \in I}$ be a family of
morphisms of schemes with target $T$. Assume that
\begin{enumerate}
\item each $f_i$ is flat, and
\item the family $\{f_i : T_i \to T\}_{i \in I}$ can be refined by an
fpqc covering of $T$.
\end{enumerate}
Then $\{f_i : T_i \to T\}_{i \in I}$ is an fpqc covering of $T$.
\end{lemma}

\begin{proof}
\uses{topologies-lemma-recognize-fpqc-covering}

Let $\{g_j : X_j \to T\}_{j \in J}$ be an fpqc covering refining
$\{f_i : T_i \to T\}$. Suppose that $U \subset T$ is affine open.
Choose $j_1, \ldots, j_m \in J$ and $V_k \subset X_{j_k}$ affine
open such that $U = \bigcup g_{j_k}(V_k)$. For each $j$ pick $i_j \in I$
and a morphism $h_j : X_j \to T_{i_j}$ such that $g_j = f_{i_j} \circ h_j$.
Since $h_{j_k}(V_k)$ is quasi-compact we can find a quasi-compact
open $h_{j_k}(V_k) \subset U_k \subset f_{i_{j_k}}^{-1}(U)$.
Then $U = \bigcup f_{i_{j_k}}(U_k)$. We conclude that
$\{f_i : T_i \to T\}_{i \in I}$ is an fpqc covering by
Lemma \ref{topologies-lemma-recognize-fpqc-covering}.
\end{proof}

\begin{lemma}[Family flat fpqc local covering]
\label{topologies-lemma-family-flat-fpqc-local-covering}
\source{stacks:03L9}
\href{https://stacks.math.columbia.edu/tag/03L9}{\texttt{Stacks tag 03L9}}
\group{topologies}

Let $T$ be a scheme. Let $\{f_i : T_i \to T\}_{i \in I}$ be a family of
morphisms of schemes with target $T$. Assume that
\begin{enumerate}
\item each $f_i$ is flat, and
\item there exists an fpqc covering
$\{g_j : S_j \to T\}_{j \in J}$ such that each
$\{S_j \times_T T_i \to S_j\}_{i \in I}$ is an fpqc covering.
\end{enumerate}
Then $\{f_i : T_i \to T\}_{i \in I}$ is an fpqc covering of $T$.
\end{lemma}

\begin{proof}
\uses{topologies-lemma-recognize-fpqc-covering, topologies-lemma-family-flat-dominated-covering}

We will use Lemma \ref{topologies-lemma-recognize-fpqc-covering} without further
mention. Let $U \subset T$ be an affine open. By (2) we can find
quasi-compact opens $V_j \subset S_j$ for $j \in J$, almost all empty, such that
$U = \bigcup g_j(V_j)$. Then for each $j$ we can choose quasi-compact
opens $W_{ij} \subset S_j \times_T T_i$ for $i \in I$, almost all empty,
with $V_j = \bigcup_i \text{pr}_1(W_{ij})$. Thus
$\{S_j \times_T T_i \to T\}$ is an fpqc covering.
Since this covering refines $\{f_i : T_i \to T\}$ we conclude by
Lemma \ref{topologies-lemma-family-flat-dominated-covering}.
\end{proof}

\begin{lemma}[Zariski etale smooth syntomic fppf fpqc]
\label{topologies-lemma-zariski-etale-smooth-syntomic-fppf-fpqc}
\source{stacks:022C}
\href{https://stacks.math.columbia.edu/tag/022C}{\texttt{Stacks tag 022C}}
\group{topologies}

Any fppf covering is an fpqc covering, and a fortiori,
any syntomic, smooth, \'etale or Zariski covering is an fpqc covering.
\end{lemma}

\begin{proof}
\uses{topologies-lemma-zariski-etale-smooth-syntomic-fppf, topologies-definition-fpqc-covering}

We will show that an fppf covering is an fpqc covering, and then the
rest follows from
Lemma \ref{topologies-lemma-zariski-etale-smooth-syntomic-fppf}.
Let $\{f_i : U_i \to U\}_{i \in I}$ be an fppf covering.
By definition this means that the $f_i$ are flat which checks the first
condition of Definition \ref{topologies-definition-fpqc-covering}. To check the
second, let $V \subset U$ be an affine open subset.
Write $f_i^{-1}(V) = \bigcup_{j \in J_i} V_{ij}$
for some affine opens $V_{ij} \subset U_i$. Since each $f_i$ is open
(Morphisms, Lemma \ref{morphisms-lemma-fppf-open}), we see that
$V = \bigcup_{i\in I} \bigcup_{j \in J_i} f_i(V_{ij})$
is an open covering of $V$.
Since $V$ is quasi-compact, this covering has a finite
refinement. This finishes the proof.
\end{proof}

The fpqc\footnote{The letters fpqc stand for
``fid\`element plat quasi-compacte''.}
topology cannot be treated in the same way as the fppf
topology\footnote{A more precise statement would be that the analogue of
Lemma \ref{topologies-lemma-fppf-induced} for the fpqc topology does not hold.}.
Namely, suppose that $R$ is a nonzero ring. We will see in
Lemma \ref{topologies-lemma-no-set-of-fpqc-covers-is-initial}
that there does not exist a set $A$ of fpqc-coverings of $\Spec(R)$
such that every fpqc-covering can be refined by an element of $A$.
If $R = k$ is a field, then the reason for this unboundedness is that
there does not exist a field extension of $k$ such that every field extension
of $k$ is contained in it.

\medskip\noindent
If you ignore set theoretic difficulties, then you run into presheaves
which do not have a sheafification, see
\cite[Theorem 5.5]{Waterhouse-fpqc-sheafification}.
A mildly interesting option is to consider only those faithfully flat ring
extensions $R \to R'$ where the cardinality of $R'$ is suitably bounded.
(And if you consider all schemes in a fixed universe as in SGA4 then you
are bounding the cardinality by a strongly inaccessible cardinal.)
However, it is not so clear what happens if you change the cardinal
to a bigger one.

\medskip\noindent
For these reasons we do not introduce fpqc sites and we will not consider
cohomology with respect to the fpqc-topology.

\medskip\noindent
On the other hand, given a contravariant functor
$F : \Sch^{opp} \to \textit{Sets}$
it does make sense to ask whether $F$ satisfies the sheaf property
for the fpqc topology, see below.
Moreover, we can wonder about descent of object
in the fpqc topology, etc. Simply put, for certain results the correct
generality is to work with fpqc coverings.

\begin{lemma}[Fpqc]
\label{topologies-lemma-fpqc}
\source{stacks:022D}
\href{https://stacks.math.columbia.edu/tag/022D}{\texttt{Stacks tag 022D}}
\group{topologies}

Let $T$ be a scheme.
\begin{enumerate}
\item If $T' \to T$ is an isomorphism then $\{T' \to T\}$
is an fpqc covering of $T$.
\item If $\{T_i \to T\}_{i\in I}$ is an fpqc covering and for each
$i$ we have an fpqc covering $\{T_{ij} \to T_i\}_{j\in J_i}$, then
$\{T_{ij} \to T\}_{i \in I, j\in J_i}$ is an fpqc covering.
\item If $\{T_i \to T\}_{i\in I}$ is an fpqc covering
and $T' \to T$ is a morphism of schemes then
$\{T' \times_T T_i \to T'\}_{i\in I}$ is an fpqc covering.
\end{enumerate}
\end{lemma}

\begin{proof}
\uses{topologies-lemma-recognize-fpqc-covering}

Part (1) is immediate. Recall that the composition of flat morphisms
is flat and that the base change of a flat morphism is flat
(Morphisms, Lemmas \ref{morphisms-lemma-base-change-flat} and
\ref{morphisms-lemma-composition-flat}).
Thus we can apply Lemma \ref{topologies-lemma-recognize-fpqc-covering}
in each case to check that our families of morphisms are fpqc coverings.

\medskip\noindent
Proof of (2). Assume $\{T_i \to T\}_{i\in I}$ is an fpqc covering and for each
$i$ we have an fpqc covering $\{f_{ij} : T_{ij} \to T_i\}_{j\in J_i}$.
Let $U \subset T$ be an affine open. We can find
quasi-compact opens $U_i \subset T_i$ for $i \in I$, almost all empty,
such that $U = \bigcup f_i(U_i)$. Then for each $i$ we can choose
quasi-compact opens $U_{ij} \subset T_{ij}$ for $j \in J_i$, almost all empty,
with $U_i = \bigcup_j f_{ij}(U_{ij})$. Thus
$\{T_{ij} \to T\}$ is an fpqc covering.

\medskip\noindent
Proof of (3). Assume $\{T_i \to T\}_{i\in I}$ is an fpqc covering
and $T' \to T$ is a morphism of schemes. Let $U' \subset T'$ be an affine
open which maps into the affine open $U \subset T$. Choose
quasi-compact opens $U_i \subset T_i$, almost all empty,
such that $U = \bigcup f_i(U_i)$. Then $U' \times_U U_i$ is
a quasi-compact open of $T' \times_T T_i$ and
$U' = \bigcup \text{pr}_1(U' \times_U U_i)$. Since $T'$
can be covered by such affine opens $U' \subset T'$ we see
that $\{T' \times_T T_i \to T'\}_{i\in I}$ is an fpqc covering by
Lemma \ref{topologies-lemma-recognize-fpqc-covering}.
\end{proof}

\begin{lemma}[Fpqc affine]
\label{topologies-lemma-fpqc-affine}
\source{stacks:022E}
\href{https://stacks.math.columbia.edu/tag/022E}{\texttt{Stacks tag 022E}}
\group{topologies}

Let $T$ be an affine scheme.
Let $\{T_i \to T\}_{i \in I}$ be an fpqc covering of $T$.
Then there exists an fpqc covering
$\{U_j \to T\}_{j = 1, \ldots, n}$ which is a refinement
of $\{T_i \to T\}_{i \in I}$ such that each $U_j$ is an affine
scheme. Moreover, we may choose each $U_j$ to be open affine
in one of the $T_i$.
\end{lemma}

\begin{proof}
This follows directly from the definition.
\end{proof}

\begin{definition}[Standard fpqc]
\label{topologies-definition-standard-fpqc}
\source{stacks:022F}
\href{https://stacks.math.columbia.edu/tag/022F}{\texttt{Stacks tag 022F}}
\group{topologies}

Let $T$ be an affine scheme. A {\it standard fpqc covering}
of $T$ is a family $\{f_j : U_j \to T\}_{j = 1, \ldots, n}$
with each $U_j$ is affine, flat over $T$ and $T = \bigcup f_j(U_j)$.
\end{definition}

Since we do not introduce the affine site we have to show directly
that the collection of all standard fpqc coverings satisfies the
axioms.

\begin{lemma}[Fpqc affine axioms]
\label{topologies-lemma-fpqc-affine-axioms}
\source{stacks:03LA}
\href{https://stacks.math.columbia.edu/tag/03LA}{\texttt{Stacks tag 03LA}}
\group{topologies}

Let $T$ be an affine scheme.
\begin{enumerate}
\item If $T' \to T$ is an isomorphism then $\{T' \to T\}$
is a standard fpqc covering of $T$.
\item If $\{T_i \to T\}_{i\in I}$ is a standard fpqc covering and for each
$i$ we have a standard fpqc covering $\{T_{ij} \to T_i\}_{j\in J_i}$, then
$\{T_{ij} \to T\}_{i \in I, j\in J_i}$ is a standard fpqc covering.
\item If $\{T_i \to T\}_{i\in I}$ is a standard fpqc covering
and $T' \to T$ is a morphism of affine schemes then
$\{T' \times_T T_i \to T'\}_{i\in I}$ is a standard fpqc covering.
\end{enumerate}
\end{lemma}

\begin{proof}
This follows formally from the fact that compositions and base changes
of flat morphisms are flat
(Morphisms, Lemmas \ref{morphisms-lemma-base-change-flat} and
\ref{morphisms-lemma-composition-flat})
and that fibre products of affine schemes are affine
(Schemes, Lemma \ref{schemes-lemma-fibre-product-affines}).
\end{proof}

\begin{lemma}[Fpqc covering affines mapping in]
\label{topologies-lemma-fpqc-covering-affines-mapping-in}
\source{stacks:03LB}
\href{https://stacks.math.columbia.edu/tag/03LB}{\texttt{Stacks tag 03LB}}
\group{topologies}

Let $T$ be a scheme. Let $\{f_i : T_i \to T\}_{i \in I}$ be a family of
morphisms of schemes with target $T$. Assume that
\begin{enumerate}
\item each $f_i$ is flat, and
\item every affine scheme
$Z$ and morphism $h : Z \to T$ there exists a standard fpqc covering
$\{Z_j \to Z\}_{j = 1, \ldots, n}$ which refines the family
$\{T_i \times_T Z \to Z\}_{i \in I}$.
\end{enumerate}
Then $\{f_i : T_i \to T\}_{i \in I}$ is an fpqc covering of $T$.
\end{lemma}

\begin{proof}
\uses{topologies-lemma-family-flat-dominated-covering, topologies-lemma-family-flat-fpqc-local-covering}

Let $T = \bigcup U_\alpha$ be an affine open covering.
For each $\alpha$ the pullback family $\{T_i \times_T U_\alpha \to U_\alpha\}$
can be refined by a standard fpqc covering, hence is an
fpqc covering by Lemma
\ref{topologies-lemma-family-flat-dominated-covering}.
As $\{U_\alpha \to T\}$ is an fpqc covering we conclude that
$\{T_i \to T\}$ is an fpqc covering by
Lemma \ref{topologies-lemma-family-flat-fpqc-local-covering}.
\end{proof}

\begin{definition}[Sheaf property fpqc]
\label{topologies-definition-sheaf-property-fpqc}
\source{stacks:022G}
\href{https://stacks.math.columbia.edu/tag/022G}{\texttt{Stacks tag 022G}}
\group{topologies}

Let $F$ be a contravariant functor on the category
of schemes with values in sets.
\begin{enumerate}
\item Let $\{U_i \to T\}_{i \in I}$ be a family of morphisms
of schemes with fixed target.
We say that $F$ {\it satisfies the sheaf property for the given family}
if for any collection of elements $\xi_i \in F(U_i)$ such that
$\xi_i|_{U_i \times_T U_j} = \xi_j|_{U_i \times_T U_j}$
there exists a unique element
$\xi \in F(T)$ such that $\xi_i = \xi|_{U_i}$ in $F(U_i)$.
\item We say that $F$ {\it satisfies the sheaf property for the
fpqc topology} if it satisfies the sheaf property for any
fpqc covering.
\end{enumerate}
\end{definition}

We try to avoid using the terminology ``$F$ is a sheaf'' in this
situation since we are not defining a category of fpqc sheaves
as we explained above.

\begin{lemma}[Sheaf property fpqc]
\label{topologies-lemma-sheaf-property-fpqc}
\source{stacks:022H}
\href{https://stacks.math.columbia.edu/tag/022H}{\texttt{Stacks tag 022H}}
\group{topologies}

Let $F$ be a contravariant functor on the category
of schemes with values in sets. Then $F$ satisfies
the sheaf property for the fpqc topology if and only
if it satisfies
\begin{enumerate}
\item the sheaf property for every Zariski covering, and
\item the sheaf property for any standard fpqc covering.
\end{enumerate}
Moreover, in the presence of (1) property (2) is equivalent to
property
\begin{enumerate}
\item[(2')] the sheaf property for $\{V \to U\}$
with $V$, $U$ affine and $V \to U$ faithfully flat.
\end{enumerate}
\end{lemma}

\begin{proof}
\uses{topologies-lemma-fpqc-affine}

Assume (1) and (2) hold.
Let $\{f_i : T_i \to T\}_{i \in I}$ be an fpqc covering. Let $s_i \in F(T_i)$
be a family of elements such that $s_i$ and $s_j$ map to the same element
of $F(T_i \times_T T_j)$. Let $W \subset T$ be the maximal open subset
such that there exists a unique $s \in F(W)$ with
$s|_{f_i^{-1}(W)} = s_i|_{f_i^{-1}(W)}$ for all $i$.
Such a maximal open exists because $F$ satisfies the
sheaf property for Zariski coverings; in fact $W$ is the
union of all opens with this property. Let $t \in T$.
We will show $t \in W$. To do this we pick an affine open
$t \in U \subset T$ and we will show there is a unique
$s \in F(U)$ with
$s|_{f_i^{-1}(U)} = s_i|_{f_i^{-1}(U)}$ for all $i$.

\medskip\noindent
By Lemma \ref{topologies-lemma-fpqc-affine} we can find a standard fpqc covering
$\{U_j \to U\}_{j = 1, \ldots, n}$ refining $\{U \times_T T_i \to U\}$,
say by morphisms $h_j : U_j \to T_{i_j}$. By (2) we obtain a unique element
$s \in F(U)$ such that $s|_{U_j} = F(h_j)(s_{i_j})$. Note that for any
scheme $V \to U$ over $U$ there is a unique section $s_V \in F(V)$
which restricts to $F(h_j \circ \text{pr}_2)(s_{i_j})$ on
$V \times_U U_j$ for $j = 1, \ldots, n$. Namely, this is true if $V$
is affine by (2) as $\{V \times_U U_j \to V\}$ is a standard fpqc covering
and in general this follows from (1) and the affine case by choosing an
affine open covering of $V$. In particular, $s_V = s|_V$.
Now, taking $V = U \times_T T_i$ and using that
$s_{i_j}|_{T_{i_j} \times_T T_i} = s_i|_{T_{i_j} \times_T T_i}$
we conclude that $s|_{U \times_T T_i} = s_V = s_i|_{U \times_T T_i}$
which is what we had to show.

\medskip\noindent
Proof of the equivalence of (2) and (2') in the presence of (1).
Suppose $\{T_i \to T\}$ is a standard fpqc covering, then
$\coprod T_i \to T$ is a faithfully flat morphism of affine schemes.
In the presence of (1) we have $F(\coprod T_i) = \prod F(T_i)$
and similarly
$F((\coprod T_i) \times_T (\coprod T_i)) = \prod F(T_i \times_T T_{i'})$.
Thus the sheaf condition for $\{T_i \to T\}$ and $\{\coprod T_i \to T\}$
is the same.
\end{proof}

The following lemma is here just to point out set theoretical difficulties
do indeed arise and should be ignored by most readers.

\begin{lemma}[No set of fpqc covers is initial]
\label{topologies-lemma-no-set-of-fpqc-covers-is-initial}
\source{stacks:0BBK}
\href{https://stacks.math.columbia.edu/tag/0BBK}{\texttt{Stacks tag 0BBK}}
\group{topologies}

Let $R$ be a nonzero ring. There does not exist a set $A$ of
fpqc-coverings of $\Spec(R)$ such that every fpqc-covering can
be refined by an element of $A$.
\end{lemma}

\begin{proof}
Let us first explain this when $R = k$ is a field. For any set $I$ consider
the purely transcendental field extension
$k_I = k(\{t_i\}_{i \in I})/k$. Since $k \to k_I$ is faithfully flat
we see that $\{\Spec(k_I) \to \Spec(k)\}$ is an fpqc covering.
Let $A$ be a set and for each $\alpha \in A$ let
$\mathcal{U}_\alpha = \{S_{\alpha, j} \to \Spec(k)\}_{j \in J_\alpha}$ be an
fpqc covering. If $\mathcal{U}_\alpha$ refines $\{\Spec(k_I) \to \Spec(k)\}$
then the morphisms $S_{\alpha, j} \to \Spec(k)$ factor through
$\Spec(k_I)$. Since $\mathcal{U}_\alpha$ is a covering,
at least some $S_{\alpha, j}$ is nonempty. Pick a
point $s \in S_{\alpha, j}$. Since we have the factorization
$S_{\alpha, j} \to \Spec(k_I) \to \Spec(k)$
we obtain a homomorphism of fields $k_I \to \kappa(s)$.
In particular, we see that the cardinality of $\kappa(s)$
is at least the cardinality of $I$. Thus if we take $I$ to be a set
of cardinality bigger than the cardinalities of the residue fields
of all the schemes $S_{\alpha, j}$, then such a factorization does
not exist and the lemma holds for $R = k$.

\medskip\noindent
General case. Since $R$ is nonzero it has a maximal prime ideal
$\mathfrak m$ with residue field $\kappa$. Let $I$ be a set and
consider $R_I = S_I^{-1} R[\{t_i\}_{i \in I}]$
where $S_I \subset R[\{t_i\}_{i \in I}]$ is the multiplicative
subset of $f \in R[\{t_i\}_{i \in I}]$ such that $f$ maps to
a nonzero element of $R/\mathfrak p[\{t_i\}_{i \in I}]$ for
all primes $\mathfrak p$ of $R$. Then $R_I$ is a faithfully
flat $R$-algebra and $\{\Spec(R_I) \to \Spec(R)\}$ is an
fpqc covering. We leave it as an exercise to the reader to show that
$R_I \otimes_R \kappa \cong \kappa(\{t_i\}_{i \in I}) = \kappa_I$
with notation as above (hint: use that $R \to \kappa$ is surjective
and that any $f \in R[\{t_i\}_{i \in I}]$ one of whose monomials occurs
with coefficient $1$ is an element of $S_I$). Let $A$ be a set and
for each $\alpha \in A$ let
$\mathcal{U}_\alpha = \{S_{\alpha, j} \to \Spec(R)\}_{j \in J_\alpha}$ be an
fpqc covering. If $\mathcal{U}_\alpha$ refines $\{\Spec(R_I) \to \Spec(R)\}$,
then by base change we conclude that
$\{S_{\alpha, j} \times_{\Spec(R)} \Spec(\kappa) \to \Spec(\kappa)\}$
refines $\{\Spec(\kappa_I) \to \Spec(\kappa)\}$.
Hence by the result of the previous paragraph, there exists an $I$
such that this is not the case and the lemma is proved.
\end{proof}

\section{The V topology}
\label{topologies-section-V}
The V topology is stronger than all other topologies in this chapter.
Roughly speaking it is generated by Zariski coverings and
by quasi-compact morphisms satisfying a lifting property
for specializations (Lemma \ref{topologies-lemma-refine-qcqs-V}).
However, the procedure we will use to define V coverings is a bit different.
We will first define standard V coverings of affines and then
use these to define V coverings in general. Typographical point:
in the literature sometimes ``$v$-covering'' is used instead of ``V covering''.

\begin{definition}[Standard V covering]
\label{topologies-definition-standard-V-covering}
\source{stacks:0ETB}
\href{https://stacks.math.columbia.edu/tag/0ETB}{\texttt{Stacks tag 0ETB}}
\group{topologies}

Let $T$ be an affine scheme. A {\it standard V covering} is a finite family
$\{T_j \to T\}_{j = 1, \ldots, m}$ with $T_j$ affine
such that for every morphism $g : \Spec(V) \to T$ where $V$
is a valuation ring, there is an extension $V \subset W$ of valuation rings
(More on Algebra, Definition
\ref{more-algebra-definition-extension-valuation-rings}),
an index $1 \leq j \leq m$, and a commutative diagram
$$
\begin{array}{cc}\Spec(W) \longrightarrow\downarrow & T_j \downarrow \\ \Spec(V) \overset{g}{\longrightarrow} & T\end{array}
$$
\end{definition}

We first prove a few basic lemmas about this notion.

\begin{lemma}[Standard fpqc standard V]
\label{topologies-lemma-standard-fpqc-standard-V}
\source{stacks:0ETC}
\href{https://stacks.math.columbia.edu/tag/0ETC}{\texttt{Stacks tag 0ETC}}
\group{topologies}

A standard fpqc covering is a standard V covering.
\end{lemma}

\begin{proof}
\uses{topologies-definition-standard-fpqc}

Let $\{X_i \to X\}_{i = 1, \ldots, n}$ be a standard fpqc covering
(Definition \ref{topologies-definition-standard-fpqc}). Let $g : \Spec(V) \to X$
be a morphism where $V$ is a valuation ring. Let $x \in X$ be the
image of the closed point of $\Spec(V)$. Choose an $i$ and a point
$x_i \in X_i$ mapping to $x$. Then $\Spec(V) \times_X X_i$
has a point $x'_i$ mapping to the closed point of $\Spec(V)$.
Since $\Spec(V) \times_X X_i \to \Spec(V)$ is flat
we can find a specialization $x''_i \leadsto x'_i$ of
points of $\Spec(V) \times_X X_i$ with $x''_i$ mapping to
the generic point of $\Spec(V)$, see
Morphisms, Lemma \ref{morphisms-lemma-generalizations-lift-flat}. By
Schemes, Lemma \ref{schemes-lemma-points-specialize}
we can choose a valuation ring $W$ and a morphism
$h : \Spec(W) \to \Spec(V) \times_X X_i$ such that $h$
maps the generic point of $\Spec(W)$ to $x''_i$ and the
closed point of $\Spec(W)$ to $x'_i$. We obtain a
commutative diagram
$$
\begin{array}{cc}\Spec(W) \longrightarrow\downarrow & X_i \downarrow \\ \Spec(V) \longrightarrow & X\end{array}
$$
where $V \to W$ is an extension of valuation rings.
This proves the lemma.
\end{proof}

\begin{lemma}[Standard ph standard V]
\label{topologies-lemma-standard-ph-standard-V}
\source{stacks:0ETD}
\href{https://stacks.math.columbia.edu/tag/0ETD}{\texttt{Stacks tag 0ETD}}
\group{topologies}

A standard ph covering is a standard V covering.
\end{lemma}

\begin{proof}
\uses{topologies-definition-standard-ph-covering}

Let $T$ be an affine scheme. Let $f : U \to T$ be a proper surjective
morphism. Let $U = \bigcup_{j = 1, \ldots, m} U_j$ be a finite
affine open covering. We have to show that $\{U_j \to T\}$
is a standard V covering, see
Definition \ref{topologies-definition-standard-ph-covering}.
Let $g : \Spec(V) \to T$
be a morphism where $V$ is a valuation ring with fraction field $K$.
Since $U \to T$ is surjective, we may choose a field extension
$L/K$ and a commutative diagram
$$
\begin{array}{ccc}\Spec(L) \longrightarrow\downarrow &  & U \downarrow \\ \Spec(K) \longrightarrow & \Spec(V) \overset{g}{\longrightarrow} & T\end{array}
$$
By Algebra, Lemma \ref{algebra-lemma-dominate} we can choose a valuation ring
$W \subset L$ dominating $V$. By the valuative criterion of
properness (Morphisms, Lemma \ref{morphisms-lemma-characterize-proper})
we can then find the morphism $h$ in the commutative
diagram
$$
\begin{array}{ccc}\Spec(L) \longrightarrow\downarrow & \Spec(W) \underset{h}{\longrightarrow}\downarrow & U \downarrow \\ \Spec(K) \longrightarrow & \Spec(V) \overset{g}{\longrightarrow} & X\end{array}
$$
Since $\Spec(W)$ has a unique closed point, we see that $\Im(h)$
is contained in $U_j$ for some $j$. Thus $h : \Spec(W) \to U_j$
is the desired lift and we conclude $\{U_j \to T\}$ is a standard V covering.
\end{proof}

\begin{lemma}[Base change standard V]
\label{topologies-lemma-base-change-standard-V}
\source{stacks:0ETE}
\href{https://stacks.math.columbia.edu/tag/0ETE}{\texttt{Stacks tag 0ETE}}
\group{topologies}

Let $\{T_j \to T\}_{j = 1, \ldots, m}$ be a standard V covering.
Let $T' \to T$ be a morphism of affine schemes.
Then $\{T_j \times_T T' \to T'\}_{j = 1, \ldots, m}$ is a
standard V covering.
\end{lemma}

\begin{proof}
Let $\Spec(V) \to T'$ be a morphism where $V$ is a valuation ring.
By assumption we can find an extension of valuation rings $V \subset W$,
an $i$, and a commutative diagram
$$
\begin{array}{cc}\Spec(W) \longrightarrow\downarrow & T_i \downarrow \\ \Spec(V) \longrightarrow & T\end{array}
$$
By the universal property of fibre products
we obtain a morphism $\Spec(W) \to T' \times_T T_i$
as desired.
\end{proof}

\begin{lemma}[Composition standard V]
\label{topologies-lemma-composition-standard-V}
\source{stacks:0ETF}
\href{https://stacks.math.columbia.edu/tag/0ETF}{\texttt{Stacks tag 0ETF}}
\group{topologies}

Let $T$ be an affine scheme. Let $\{T_j \to T\}_{j = 1, \ldots, m}$
be a standard V covering. Let $\{T_{ji} \to T_j\}_{i = 1, \ldots n_j}$
be a standard V covering. Then $\{T_{ji} \to T\}_{i, j}$
is a standard V covering.
\end{lemma}

\begin{proof}
This follows formally from the observation that if
$V \subset W$ and $W \subset \Omega$ are extensions of
valuation rings, then $V \subset \Omega$ is an extension
of valuation rings.
\end{proof}

\begin{lemma}[Refine standard V]
\label{topologies-lemma-refine-standard-V}
\source{stacks:0ETG}
\href{https://stacks.math.columbia.edu/tag/0ETG}{\texttt{Stacks tag 0ETG}}
\group{topologies}

Let $T$ be an affine scheme. Let $\{T_j \to T\}_{j = 1, \ldots, m}$
be a family of morphisms with $T_j$ affine for all $j$.
The following are equivalent
\begin{enumerate}
\item $\{T_j \to T\}_{j = 1, \ldots, m}$ is a standard V covering,
\item there is a standard V covering which refines
$\{T_j \to T\}_{j = 1, \ldots, m}$, and
\item $\{\coprod_{j = 1, \ldots, m} T_j \to T\}$ is a standard
V covering.
\end{enumerate}
\end{lemma}

\begin{proof}
Omitted. Hints: This follows almost immediately from the definition.
The only slightly interesting point is that
a morphism from the spectrum of a local ring into
$\coprod_{j = 1, \ldots, m} T_j$ must factor through some $T_j$.
\end{proof}

\begin{definition}[V covering]
\label{topologies-definition-V-covering}
\source{stacks:0ETH}
\href{https://stacks.math.columbia.edu/tag/0ETH}{\texttt{Stacks tag 0ETH}}
\group{topologies}

Let $T$ be a scheme. A {\it V covering of $T$} is a family
of morphisms $\{T_i \to T\}_{i \in I}$ of schemes such that for every
affine open $U \subset T$ there exists a standard V covering
$\{U_j \to U\}_{j = 1, \ldots, m}$ refining the family
$\{T_i \times_T U \to U\}_{i \in I}$.
\end{definition}

The V topology has the same set theoretical problems as
the fpqc topology. Thus we refrain from defining V sites
and we will not consider cohomology with respect to the V topology.
On the other hand, given a
$F : \Sch^{opp} \to \textit{Sets}$
it does make sense to ask whether $F$ satisfies the sheaf property
for the V topology, see below.
Moreover, we can wonder about descent of object
in the V topology, etc.

\begin{lemma}[Refine by V]
\label{topologies-lemma-refine-by-V}
\source{stacks:0ETI}
\href{https://stacks.math.columbia.edu/tag/0ETI}{\texttt{Stacks tag 0ETI}}
\group{topologies}

Let $T$ be a scheme. Let $\{f_i : T_i \to T\}_{i \in I}$ be a family
of morphisms. The following are equivalent
\begin{enumerate}
\item $\{T_i \to T\}_{i \in I}$ is a V covering,
\item there is a V covering which refines $\{T_i \to T\}_{i \in I}$, and
\item $\{\coprod_{i \in I} T_i \to T\}$ is a V covering.
\end{enumerate}
\end{lemma}

\begin{proof}
\uses{topologies-lemma-refine-by-ph}

Omitted. Hint: compare with the proof of
Lemma \ref{topologies-lemma-refine-by-ph}.
\end{proof}

\begin{lemma}[V]
\label{topologies-lemma-V}
\source{stacks:0ETJ}
\href{https://stacks.math.columbia.edu/tag/0ETJ}{\texttt{Stacks tag 0ETJ}}
\group{topologies}

Let $T$ be a scheme.
\begin{enumerate}
\item If $T' \to T$ is an isomorphism then $\{T' \to T\}$
is a V covering of $T$.
\item If $\{T_i \to T\}_{i\in I}$ is a V covering and for each
$i$ we have a V covering $\{T_{ij} \to T_i\}_{j\in J_i}$, then
$\{T_{ij} \to T\}_{i \in I, j\in J_i}$ is a V covering.
\item If $\{T_i \to T\}_{i\in I}$ is a V covering
and $T' \to T$ is a morphism of schemes then
$\{T' \times_T T_i \to T'\}_{i\in I}$ is a V covering.
\end{enumerate}
\end{lemma}

\begin{proof}
\uses{topologies-lemma-base-change-standard-V, topologies-lemma-standard-fpqc-standard-V, topologies-lemma-composition-standard-V}

Assertion (1) is clear.

\medskip\noindent
Proof of (3).
Let $U' \subset T'$ be an affine open subscheme.
Since $U'$ is quasi-compact we can find a finite affine open
covering $U' = U'_1 \cup \ldots \cup U'$ such that
$U'_j \to T$ maps into an affine open $U_j \subset T$.
Choose a standard V covering $\{U_{jl} \to U_j\}_{l = 1, \ldots, n_j}$
refining $\{T_i \times_T U_j \to U_j\}$.
By Lemma \ref{topologies-lemma-base-change-standard-V}
the base change $\{U_{jl} \times_{U_j} U'_j \to U'_j\}$
is a standard V covering. Note that $\{U'_j \to U'\}$
is a standard V covering
(for example by Lemma \ref{topologies-lemma-standard-fpqc-standard-V}).
By Lemma \ref{topologies-lemma-composition-standard-V}
the family $\{U_{jl} \times_{U_j} U'_j \to U'\}$
is a standard V covering. Since
$\{U_{jl} \times_{U_j} U'_j \to U'\}$ refines $\{T_i \times_T U' \to U'\}$
we conclude.

\medskip\noindent
Proof of (2).
Let $U \subset T$ be affine open. First we pick a standard V covering
$\{U_k \to U\}_{k = 1, \ldots, m}$ refining
$\{T_i \times_T U \to U\}$.
Say the refinement is given by morphisms $U_k \to T_{i_k}$ over $T$.
Then
$$
\{T_{i_kj} \times_{T_{i_k}} U_k \to U_k\}_{j \in J_{i_k}}
$$
is a V covering by part (3). As $U_k$ is affine,
we can find a standard V covering
$\{U_{ka} \to U_k\}_{a = 1, \ldots, b_k}$ refining this family.
Then we apply Lemma \ref{topologies-lemma-composition-standard-V}
to see that $\{U_{ka} \to U\}$ is a standard V covering which
refines $\{T_{ij} \times_T U \to U\}$. This finishes the proof.
\end{proof}

\begin{lemma}[Zariski etale smooth syntomic fppf fpqc ph V]
\label{topologies-lemma-zariski-etale-smooth-syntomic-fppf-fpqc-ph-V}
\source{stacks:0ETK}
\href{https://stacks.math.columbia.edu/tag/0ETK}{\texttt{Stacks tag 0ETK}}
\group{topologies}

Any fpqc covering is a V covering. A fortiori,
any fppf, syntomic, smooth, \'etale or Zariski covering is a V covering.
Also, a ph covering is a V covering.
\end{lemma}

\begin{proof}
\uses{topologies-lemma-fpqc-affine, topologies-lemma-standard-fpqc-standard-V, topologies-lemma-zariski-etale-smooth-syntomic-fppf-fpqc, topologies-lemma-standard-ph-standard-V}

An fpqc covering can affine locally be refined by a
standard fpqc covering, see Lemmas \ref{topologies-lemma-fpqc-affine}.
A standard fpqc covering is a standard V covering, see
Lemma \ref{topologies-lemma-standard-fpqc-standard-V}. Hence the first statement
follows from our definition of V covers in terms of standard V coverings.
The conclusion for fppf, syntomic, smooth, \'etale or Zariski coverings
follows as these are fpqc coverings, see
Lemma \ref{topologies-lemma-zariski-etale-smooth-syntomic-fppf-fpqc}.

\medskip\noindent
The statement on ph coverings follows from
Lemma \ref{topologies-lemma-standard-ph-standard-V} in the same manner.
\end{proof}

\begin{definition}[Sheaf property V]
\label{topologies-definition-sheaf-property-V}
\source{stacks:0ETL}
\href{https://stacks.math.columbia.edu/tag/0ETL}{\texttt{Stacks tag 0ETL}}
\group{topologies}

\uses{topologies-definition-sheaf-property-fpqc}
Let $F$ be a contravariant functor on the category
of schemes with values in sets. We say that
$F$ {\it satisfies the sheaf property for the V topology}
if it satisfies the sheaf property for any V covering
(see Definition \ref{topologies-definition-sheaf-property-fpqc}).
\end{definition}

We try to avoid using the terminology ``$F$ is a sheaf'' in this
situation since we are not defining a category of V sheaves
as we explained above.

\begin{lemma}[Sheaf property V]
\label{topologies-lemma-sheaf-property-V}
\source{stacks:0ETM}
\href{https://stacks.math.columbia.edu/tag/0ETM}{\texttt{Stacks tag 0ETM}}
\group{topologies}

Let $F$ be a contravariant functor on the category
of schemes with values in sets. Then $F$ satisfies
the sheaf property for the V topology if and only
if it satisfies
\begin{enumerate}
\item the sheaf property for every Zariski covering, and
\item the sheaf property for any standard V covering.
\end{enumerate}
Moreover, in the presence of (1) property (2) is equivalent to
property
\begin{enumerate}
\item[(2')] the sheaf property for a standard V covering
of the form $\{V \to U\}$, i.e., consisting of a single arrow.
\end{enumerate}
\end{lemma}

\begin{proof}
\uses{topologies-lemma-base-change-standard-V}

Assume (1) and (2) hold.
Let $\{f_i : T_i \to T\}_{i \in I}$ be a V covering. Let $s_i \in F(T_i)$
be a family of elements such that $s_i$ and $s_j$ map to the same element
of $F(T_i \times_T T_j)$. Let $W \subset T$ be the maximal open subset
such that there exists a unique $s \in F(W)$ with
$s|_{f_i^{-1}(W)} = s_i|_{f_i^{-1}(W)}$ for all $i$.
Such a maximal open exists because $F$ satisfies the
sheaf property for Zariski coverings; in fact $W$ is the
union of all opens with this property. Let $t \in T$.
We will show $t \in W$. To do this we pick an affine open
$t \in U \subset T$ and we will show there is a unique
$s \in F(U)$ with
$s|_{f_i^{-1}(U)} = s_i|_{f_i^{-1}(U)}$ for all $i$.

\medskip\noindent
We can find a standard V covering
$\{U_j \to U\}_{j = 1, \ldots, n}$ refining $\{U \times_T T_i \to U\}$,
say by morphisms $h_j : U_j \to T_{i_j}$. By (2) we obtain a unique element
$s \in F(U)$ such that $s|_{U_j} = F(h_j)(s_{i_j})$. Note that for any
scheme $V \to U$ over $U$ there is a unique section $s_V \in F(V)$
which restricts to $F(h_j \circ \text{pr}_2)(s_{i_j})$ on
$V \times_U U_j$ for $j = 1, \ldots, n$. Namely, this is true if $V$
is affine by (2) as $\{V \times_U U_j \to V\}$ is a standard V covering
(Lemma \ref{topologies-lemma-base-change-standard-V})
and in general this follows from (1) and the affine case by choosing an
affine open covering of $V$. In particular, $s_V = s|_V$.
Now, taking $V = U \times_T T_i$ and using that
$s_{i_j}|_{T_{i_j} \times_T T_i} = s_i|_{T_{i_j} \times_T T_i}$
we conclude that $s|_{U \times_T T_i} = s_V = s_i|_{U \times_T T_i}$
which is what we had to show.

\medskip\noindent
Proof of the equivalence of (2) and (2') in the presence of (1).
Suppose $\{T_i \to T\}_{i = 1, \ldots, n}$ is a standard V covering, then
$\coprod_{i = 1, \ldots, n} T_i \to T$ is a morphism of affine schemes
which is clearly also a standard V covering.
In the presence of (1) we have $F(\coprod T_i) = \prod F(T_i)$
and similarly
$F((\coprod T_i) \times_T (\coprod T_i)) = \prod F(T_i \times_T T_{i'})$.
Thus the sheaf condition for $\{T_i \to T\}$ and $\{\coprod T_i \to T\}$
is the same.
\end{proof}

The following lemma shows that being a V covering is
related to the possibility of lifting specializations.

\begin{lemma}[Refine qcqs V]
\label{topologies-lemma-refine-qcqs-V}
\source{stacks:0ETN}
\href{https://stacks.math.columbia.edu/tag/0ETN}{\texttt{Stacks tag 0ETN}}
\group{topologies}

Let $X \to Y$ be a quasi-compact morphism of schemes.
The following are equivalent
\begin{enumerate}
\item $\{X \to Y\}$ is a V covering,
\item for any valuation ring $V$ and morphism $g : \Spec(V) \to Y$
there exists an extension of valuation rings $V \subset W$
and a commutative diagram
$$
\begin{array}{cc}\Spec(W) \longrightarrow\downarrow & X \downarrow \\ \Spec(V) \longrightarrow & Y\end{array}
$$
\item for any morphism $Z \to Y$ and specialization $z' \leadsto z$
of points in $Z$, there is a specialization $w' \leadsto w$
of points in $Z \times_Y X$ mapping to $z' \leadsto z$.
\end{enumerate}
\end{lemma}

\begin{proof}
\uses{topologies-definition-V-covering, topologies-definition-standard-V-covering}

Assume (1) and let $g : \Spec(V) \to Y$ be as in (2).
Since $V$ is a local ring there is an affine open $U \subset Y$
such that $g$ factors through $U$. By Definition \ref{topologies-definition-V-covering}
we can find a standard V covering $\{U_j \to U\}$ refining
$\{X \times_Y U \to U\}$. By Definition \ref{topologies-definition-standard-V-covering}
we can find a $j$, an extension of valuation rings $V \subset W$
and a commutative diagram
$$
\begin{array}{ccc}\Spec(W) \longrightarrow\downarrow & U_j \downarrow\dashrightarrow & X \swarrow \\ \Spec(V) \longrightarrow & Y & \end{array}
$$
We have the dotted arrow making the diagram commute by the refinement
property of the covering and we see that (2) holds.

\medskip\noindent
Assume (2) and let $Z \to Y$ and $z' \leadsto z$ be as in (3).
By Schemes, Lemma \ref{schemes-lemma-points-specialize}
we can find a valuation ring $V$ and a morphism $\Spec(V) \to Z$
such that the closed point of $\Spec(V)$ maps to $z$ and the
generic point of $\Spec(V)$ maps to $z'$. By (2) we can find an
extension of valuation rings $V \subset W$ and a
commutative diagram
$$
\begin{array}{ccc}\Spec(W) \longrightarrow\downarrow &  & X \downarrow \\ \Spec(V) \longrightarrow & Z \longrightarrow & Y\end{array}
$$
The generic and closed points of $\Spec(W)$ map to points
$w' \leadsto w$ in $Z \times_Y X$ via the
induced morphism $\Spec(W) \to Z \times_Y X$.
This shows that (3) holds.

\medskip\noindent
Assume (3) holds and let $U \subset Y$ be an affine open.
Choose a finite affine open covering
$U \times_Y X = \bigcup_{j = 1, \ldots, m} U_j$.
This is possible as $X \to Y$ is quasi-compact.
We claim that $\{U_j \to U\}$ is a standard V covering.
The claim implies (1) is true and finishes the proof of the lemma.
In order to prove the claim, let $V$ be a valuation ring
and let $g : \Spec(V) \to U$ be a morphism.
By (3) we find a specialization $w' \leadsto w$ of
points of
$$
T = \Spec(V) \times_X Y = \Spec(V) \times_U (U \times_X Y)
$$
such that $w'$ maps to the generic point of $\Spec(V)$
and $w$ maps to the closed point of $\Spec(V)$.
By Schemes, Lemma \ref{schemes-lemma-points-specialize}
we can find a valuation ring $W$ and a morphism
$\Spec(W) \to T$ such that the generic point of $\Spec(W)$
maps to $w'$ and the closed point of $\Spec(W)$ maps to $w$.
The composition $\Spec(W) \to T \to \Spec(V)$ corresponds
to an inclusion $V \subset W$ which presents $W$ as an
extension of the valuation ring $V$.
Since $T = \bigcup \Spec(V) \times_U U_j$ is an open
covering, we see that $\Spec(W) \to T$ factors through
$\Spec(V) \times_U U_j$ for some $j$. Thus we obtain
a commutative diagram
$$
\begin{array}{cc}\Spec(W) \downarrow\longrightarrow & U_j \downarrow \\ \Spec(V) \longrightarrow & U\end{array}
$$
and the proof of the claim is complete.
\end{proof}

A V covering gives a universally submersive family of maps.
The converse of this lemma is false, see
Examples, Section \ref{examples-section-universally-submersive-not-V}.

\begin{lemma}[V covering universally submersive]
\label{topologies-lemma-V-covering-universally-submersive}
\source{stacks:0ETP}
\href{https://stacks.math.columbia.edu/tag/0ETP}{\texttt{Stacks tag 0ETP}}
\group{topologies}

Let $\{f_i : X_i \to X\}_{i \in I}$ be a V covering.
Then
$$
\coprod\nolimits_{i \in I} f_i :
\coprod\nolimits_{i \in I} X_i
\longrightarrow
X
$$
is a universally submersive morphism of schemes (Morphisms, Definition
\ref{morphisms-definition-submersive}).
\end{lemma}

\begin{proof}
\uses{topologies-lemma-V, topologies-lemma-refine-qcqs-V}

We will use without further mention that the base change of a V covering
is a V covering (Lemma \ref{topologies-lemma-V}). In particular it suffices
to show that the morphism is submersive.
Being submersive is clearly Zariski local on the base.
Thus we may assume $X$ is affine. Then $\{X_i \to X\}$
can be refined by a standard V covering $\{Y_j \to X\}$.
If we can show that $\coprod Y_j \to X$ is submersive,
then since there is a factorization $\coprod Y_j \to \coprod X_i \to X$
we conclude that $\coprod X_i \to X$ is submersive.
Set $Y = \coprod Y_j$ and consider the morphism of affines
$f : Y \to X$.
By Lemma \ref{topologies-lemma-refine-qcqs-V} we know that we can lift any
specialization $x' \leadsto x$ in $X$ to some specialization
$y' \leadsto y$ in $Y$.
Thus if $T \subset X$ is a subset such that $f^{-1}(T)$
is closed in $Y$, then $T \subset X$ is closed under specialization.
Since $f^{-1}(T) \subset Y$
with the reduced induced closed subscheme structure
is an affine scheme, we conclude that $T \subset X$ is closed by
Algebra, Lemma \ref{algebra-lemma-image-stable-specialization-closed}.
Hence $f$ is submersive.
\end{proof}

\section{Change of topologies}
\label{topologies-section-change-topologies}
Let $f : X \to Y$ be a morphism of schemes over a base scheme $S$.
In this case we have the following morphisms of sites\footnote{We have
not included the comparison between the ph topology and the others;
for this see
More on Morphisms, Remark \ref{more-morphisms-remark-change-topologies}.}
(with suitable choices of sites as in Remark \ref{topologies-remark-choice-sites}
below):
\begin{enumerate}
\item $(\Sch/X)_{fppf} \longrightarrow (\Sch/Y)_{fppf}$,
\item $(\Sch/X)_{fppf} \longrightarrow (\Sch/Y)_{syntomic}$,
\item $(\Sch/X)_{fppf} \longrightarrow (\Sch/Y)_{smooth}$,
\item $(\Sch/X)_{fppf} \longrightarrow
(\Sch/Y)_\etale$,
\item $(\Sch/X)_{fppf} \longrightarrow (\Sch/Y)_{Zar}$,
\item $(\Sch/X)_{syntomic} \longrightarrow (\Sch/Y)_{syntomic}$,
\item $(\Sch/X)_{syntomic} \longrightarrow (\Sch/Y)_{smooth}$,
\item $(\Sch/X)_{syntomic} \longrightarrow
(\Sch/Y)_\etale$,
\item $(\Sch/X)_{syntomic} \longrightarrow (\Sch/Y)_{Zar}$,
\item $(\Sch/X)_{smooth} \longrightarrow (\Sch/Y)_{smooth}$,
\item $(\Sch/X)_{smooth} \longrightarrow
(\Sch/Y)_\etale$,
\item $(\Sch/X)_{smooth} \longrightarrow (\Sch/Y)_{Zar}$,
\item $(\Sch/X)_\etale \longrightarrow
(\Sch/Y)_\etale$,
\item $(\Sch/X)_\etale \longrightarrow (\Sch/Y)_{Zar}$,
\item $(\Sch/X)_{Zar} \longrightarrow (\Sch/Y)_{Zar}$,
\item $(\Sch/X)_{fppf} \longrightarrow Y_\etale$,
\item $(\Sch/X)_{syntomic} \longrightarrow Y_\etale$,
\item $(\Sch/X)_{smooth} \longrightarrow Y_\etale$,
\item $(\Sch/X)_\etale \longrightarrow Y_\etale$,
\item $(\Sch/X)_{fppf} \longrightarrow Y_{Zar}$,
\item $(\Sch/X)_{syntomic} \longrightarrow Y_{Zar}$,
\item $(\Sch/X)_{smooth} \longrightarrow Y_{Zar}$,
\item $(\Sch/X)_\etale \longrightarrow Y_{Zar}$,
\item $(\Sch/X)_{Zar} \longrightarrow Y_{Zar}$,
\item $X_\etale \longrightarrow Y_\etale$,
\item $X_\etale \longrightarrow Y_{Zar}$,
\item $X_{Zar} \longrightarrow Y_{Zar}$,
\end{enumerate}
In each case the underlying continuous functor
$\Sch/Y \to \Sch/X$, or
$Y_\tau \to \Sch/X$
is the functor $Y'/Y \mapsto X \times_Y Y'/X$. Namely, in the sections
above we have seen the morphisms
$f_{big} : (\Sch/X)_\tau \to (\Sch/Y)_\tau$
and
$f_{small} : X_\tau \to Y_\tau$
for $\tau$ as above.
We also have seen the morphisms of sites
$\pi_Y : (\Sch/Y)_\tau \to Y_\tau$ for
$\tau \in \{\etale, Zariski\}$.
On the other hand, it is clear that the identity functor
$(\Sch/X)_\tau \to (\Sch/X)_{\tau'}$ defines
a morphism of sites when $\tau$ is a stronger topology than
$\tau'$. Hence composing these gives the list of possible morphisms
above.

\medskip\noindent
Because of the simple description of the underlying functor it
is clear that given morphisms of schemes $X \to Y \to Z$ the
composition of two of the morphisms of sites above, e.g.,
$$
(\Sch/X)_{\tau_0} \longrightarrow
(\Sch/Y)_{\tau_1} \longrightarrow
(\Sch/Z)_{\tau_2}
$$
is the corresponding morphism of sites associated to the morphism
of schemes $X \to Z$.

\begin{remark}[Choice sites]
\label{topologies-remark-choice-sites}
\source{stacks:03FF}
\href{https://stacks.math.columbia.edu/tag/03FF}{\texttt{Stacks tag 03FF}}
\group{topologies}

Take any category $\Sch_\alpha$ constructed as in
Sets, Lemma \ref{sets-lemma-construct-category}
starting with the set of schemes $\{X, Y, S\}$. Choose any set of
coverings $\text{Cov}_{fppf}$ on $\Sch_\alpha$ as in
Sets, Lemma \ref{sets-lemma-coverings-site}
starting with the category $\Sch_\alpha$ and the class of fppf
coverings. Let $\Sch_{fppf}$ denote the big fppf site so
obtained. Next, for $\tau \in \{Zariski, \etale, smooth, syntomic\}$
let $\Sch_\tau$ have the same underlying category as
$\Sch_{fppf}$ with coverings
$\text{Cov}_\tau \subset \text{Cov}_{fppf}$ simply the subset of
$\tau$-coverings. It is straightforward to check that this gives rise
to a big site $\Sch_\tau$.
\end{remark}

\section{Change of big sites}
\label{topologies-section-change-alpha}
In this section we explain what happens on changing the big
Zariski/fppf/\'etale sites.

\medskip\noindent
Let $\tau, \tau' \in \{Zariski, \etale, smooth, syntomic, fppf\}$.
Given two big sites $\Sch_\tau$ and $\Sch'_{\tau'}$
we say that
{\it $\Sch_\tau$ is contained in $\Sch'_{\tau'}$} if
$\Ob(\Sch_\tau) \subset \Ob(\Sch'_{\tau'})$
and
$\text{Cov}(\Sch_\tau) \subset \text{Cov}(\Sch'_{\tau'})$.
In this case $\tau$ is stronger than $\tau'$, for example, no fppf
site can be contained in an \'etale site.

\begin{lemma}[Contained in]
\label{topologies-lemma-contained-in}
\source{stacks:022J}
\href{https://stacks.math.columbia.edu/tag/022J}{\texttt{Stacks tag 022J}}
\group{topologies}

Any set of big Zariski sites is contained in a common big Zariski site.
The same is true, mutatis mutandis, for big fppf and big \'etale sites.
\end{lemma}

\begin{proof}
This is true because the union of a set of sets is a set, and the
constructions in Sets, Lemmas \ref{sets-lemma-construct-category} and
\ref{sets-lemma-coverings-site}
allow one to start with any initially given set of schemes
and coverings.
\end{proof}

\begin{lemma}[Change alpha]
\label{topologies-lemma-change-alpha}
\source{stacks:022K}
\href{https://stacks.math.columbia.edu/tag/022K}{\texttt{Stacks tag 022K}}
\group{topologies}

Let $\tau \in \{Zariski, \etale, smooth, syntomic, fppf\}$.
Suppose given big sites $\Sch_\tau$ and $\Sch'_\tau$.
Assume that $\Sch_\tau$ is contained in $\Sch'_\tau$.
The inclusion functor $\Sch_\tau \to \Sch'_\tau$ satisfies
the assumptions of Sites, Lemma \ref{sites-lemma-bigger-site}.
There are morphisms of topoi
\begin{eqnarray*}
g : \Sh(\Sch_\tau) &
\longrightarrow &
\Sh(\Sch'_\tau) \\
f : \Sh(\Sch'_\tau) &
\longrightarrow &
\Sh(\Sch_\tau)
\end{eqnarray*}
such that $f \circ g \cong \text{id}$. For any object $S$
of $\Sch_\tau$ the inclusion functor
$(\Sch/S)_\tau \to (\Sch'/S)_\tau$ satisfies
the assumptions of Sites, Lemma \ref{sites-lemma-bigger-site}
also. Hence similarly we obtain morphisms
\begin{eqnarray*}
g : \Sh((\Sch/S)_\tau) &
\longrightarrow &
\Sh((\Sch'/S)_\tau) \\
f : \Sh((\Sch'/S)_\tau) &
\longrightarrow &
\Sh((\Sch/S)_\tau)
\end{eqnarray*}
with $f \circ g \cong \text{id}$.
\end{lemma}

\begin{proof}
\uses{topologies-lemma-zariski-induced, topologies-lemma-etale-induced, topologies-lemma-smooth-induced, topologies-lemma-syntomic-induced, topologies-lemma-fppf-induced}

Assumptions (b), (c), and (e) of
Sites, Lemma \ref{sites-lemma-bigger-site}
are immediate for the functors
$\Sch_\tau \to \Sch'_\tau$ and
$(\Sch/S)_\tau \to (\Sch'/S)_\tau$. Property (a) holds by
Lemma \ref{topologies-lemma-zariski-induced},
\ref{topologies-lemma-etale-induced},
\ref{topologies-lemma-smooth-induced},
\ref{topologies-lemma-syntomic-induced}, or
\ref{topologies-lemma-fppf-induced}.
Property (d) holds because
fibre products in the categories $\Sch_\tau$, $\Sch'_\tau$
exist and are compatible with fibre products in the category of schemes.
\end{proof}

Discussion:
The functor $g^{-1} = f_*$ is simply the restriction functor which associates
to a sheaf $\mathcal{G}$ on $\Sch'_\tau$ the restriction
$\mathcal{G}|_{\Sch_\tau}$. Hence this lemma simply says that given
any sheaf of sets $\mathcal{F}$ on $\Sch_\tau$ there exists a
canonical sheaf $\mathcal{F}'$ on $\Sch'_\tau$ such that
$\mathcal{F}|_{\Sch'_\tau} = \mathcal{F}'$. In fact the sheaf
$\mathcal{F}'$ has the following description: it is the sheafification
of the presheaf
$$
\Sch'_\tau \longrightarrow \textit{Sets}, \quad
V \longmapsto \colim_{V \to U} \mathcal{F}(U)
$$
where $U$ is an object of $\Sch_\tau$. This is true because
$\mathcal{F}' = f^{-1}\mathcal{F} = (u_p\mathcal{F})^\#$ according to
Sites, Lemmas \ref{sites-lemma-when-shriek} and \ref{sites-lemma-bigger-site}.

\section{Extending functors}
\label{topologies-section-extending-functors}
Let us start with a simple example which explains what we are doing.
Let $R$ be a ring. Suppose $F$ is a functor defined on the category
$\mathcal{C}$ of $R$-algebras of the form
$$
A = R[x_1, \ldots, x_n]/(f_1, \ldots, f_m)
$$
for $n, m \geq 0$ integers and $f_1, \ldots, f_m \in R[x_1, \ldots, x_n]$
elements. Then for any $R$-algebra $B$ we can define
$$
F'(B) = \colim_{A \to B,\ A \in \mathcal{C}} F(A)
$$
It turns out $F'$ is the unique functor on the category
of all $R$-algebras which extends $F$ and commutes with filtered
colimits. The same procedure works in the category of schemes
if we impose that our functor is a Zariski sheaf.

\begin{lemma}[Extend]
\label{topologies-lemma-extend}
\source{stacks:0EUW}
\href{https://stacks.math.columbia.edu/tag/0EUW}{\texttt{Stacks tag 0EUW}}
\group{topologies}

Let $S$ be a scheme. Let $\mathcal{C}$ be a full subcategory
of the category $\Sch/S$ of all schemes over $S$. Assume
\begin{enumerate}
\item if $X \to S$ is an object of $\mathcal{C}$ and
$U \subset X$ is an affine open, then $U \to S$ is isomorphic
to an object of $\mathcal{C}$,
\item if $V$ is an affine scheme lying over an affine open $U \subset S$
such that $V \to U$ is of finite presentation, then $V \to S$ is isomorphic
to an object of $\mathcal{C}$.
\end{enumerate}
Let $F : \mathcal{C}^{opp} \to \textit{Sets}$ be a functor.
Assume
\begin{enumerate}
\item[(a)] for any Zariski covering $\{f_i : X_i \to X\}_{i \in I}$
with $X, X_i$ objects of $\mathcal{C}$ we have
the sheaf condition for $F$ and this family\footnote{As we
do not know that $X_i \times_X X_j$ is in $\mathcal{C}$
this has to be interpreted as follows: by property (1)
there exist Zariski coverings $\{U_{ijk} \to X_i \times_X X_j\}_{k \in K_{ij}}$
with $U_{ijk}$ an object of $\mathcal{C}$. Then the sheaf condition
says that $F(X)$ is the equalizer of the two maps
from $\prod F(X_i)$ to $\prod F(U_{ijk})$.},
\item[(b)] if $X = \lim X_i$ is a directed limit of affine schemes
over $S$ with $X, X_i$ objects of $\mathcal{C}$, then
$F(X) = \colim F(X_i)$.
\end{enumerate}
Then there is a unique way to extend $F$ to a functor
$F' : (\Sch/S)^{opp} \to \textit{Sets}$ satisfying the analogues of
(a) and (b), i.e., $F'$ satisfies the sheaf condition for any
Zariski covering and $F'(X) = \colim F'(X_i)$ whenever $X = \lim X_i$
is a directed limit of affine schemes over $S$.
\end{lemma}

\begin{proof}
The idea will be to first extend $F$ to a sufficiently large collection of
affine schemes over $S$ and then use the Zariski sheaf property to extend
to all schemes.

\medskip\noindent
Suppose that $V$ is an affine scheme over $S$ whose structure morphism
$V \to S$ factors through some affine open $U \subset S$. In this case
we can write
$$
V = \lim V_i
$$
as a cofiltered limit with $V_i \to U$ of finite presentation
and $V_i$ affine. See Algebra, Lemma \ref{algebra-lemma-ring-colimit-fp}.
By conditions (1) and (2)
we may replace our $V_i$ by objects of $\mathcal{C}$.
Observe that $V_i \to S$ is locally of finite presentation
(if $S$ is quasi-separated, then these morphisms are actually
of finite presentation). Then we set
$$
F'(V) = \colim F(V_i)
$$
Actually, we can give a more canonical expression, namely
$$
F'(V) = \colim_{V \to V'} F(V')
$$
where the colimit is over the category of morphisms $V \to V'$ over $S$
where $V'$ is an object of $\mathcal{C}$ whose structure
morphism $V' \to S$ is locally of finite presentation.
The reason this is the same as the first formula is that by
Limits, Proposition
\ref{limits-proposition-characterize-locally-finite-presentation}
our inverse system $V_i$ is cofinal in this category!
Finally, note that if $V$ were an object of $\mathcal{C}$,
then $F'(V) = F(V)$ by assumption (b).

\medskip\noindent
The second formula turns $F'$ into a contravariant functor
on the category formed by affine schemes $V$ over $S$ whose
structure morphism factors through an affine open of $S$.
Let $V$ be such an affine scheme over $S$ and
suppose that $V = \bigcup_{k = 1, \ldots, n} V_k$ is a finite open covering
by affines. Then it makes sense to ask if the sheaf condition
holds for $F'$ and this open covering.
This is true and easy to show: write $V = \lim V_i$ as in
the previous paragraph. By Limits, Lemma \ref{limits-lemma-descend-opens}
for all sufficiently large $i$ we can find affine opens
$V_{i, k} \subset V_i$ compatible with transition maps
pulling back to $V_k$ in $V$. Thus
$$
F'(V_k) = \colim F(V_{i, k})
\quad\text{and}\quad
F'(V_k \cap V_l) = \colim F(V_{i, k} \cap V_{i, l})
$$
Strictly speaking in these formulas we need to replace $V_{i, k}$ and
$V_{i, k} \cap V_{i, l}$ by isomorphic affine objects of $\mathcal{C}$
before applying the functor $F$.
Since $I$ is directed the colimits pass through equalizers.
Hence the sheaf condition (b) for $F$ and the Zariski coverings
$\{V_{i, k} \to V_i\}$ implies the sheaf
condition for $F'$ and this covering.

\medskip\noindent
Let $X$ be a general scheme over $S$. Let $\mathcal{B}_X$ denote
the collection of affine opens of $X$ whose structure morphism
to $S$ maps into an affine open of $S$. It is clear that
$\mathcal{B}_X$ is a basis for the topology of $X$.
By the result of the previous paragraph and Sheaves, Lemma
\ref{sheaves-lemma-cofinal-systems-coverings-standard-case}
we see that $F'$ is a sheaf on $\mathcal{B}_X$.
Hence $F'$ restricted to $\mathcal{B}_X$
extends uniquely to a sheaf $F'_X$ on $X$, see
Sheaves, Lemma \ref{sheaves-lemma-extend-off-basis}.
If $X$ is an object of $\mathcal{C}$ then we have a canonical
identification $F'_X(X) = F(X)$ by the agreement of $F'$ and $F$
on the objects for which they are both defined and
the fact that $F$ satisfies the sheaf condition for
Zariski coverings.

\medskip\noindent
Let $f : X \to Y$ be a morphism of schemes over $S$.
We get a unique $f$-map from $F'_Y$ to $F'_X$ compatible
with the maps $F'(V) \to F'(U)$ for all
$U \in \mathcal{B}_X$ and $V \in \mathcal{B}_Y$
with $f(U) \subset V$, see
Sheaves, Lemma \ref{sheaves-lemma-f-map-basis-above-and-below-structures}.
We omit the verification that these maps compose
correctly given morphisms $X \to Y \to Z$ of schemes over $S$.
We also omit the verification that if $f$ is a morphism
of $\mathcal{C}$, then the induced map $F'_Y(Y) \to F'_X(X)$
is the same as the map $F(Y) \to F(X)$ via the identifications
$F'_X(X) = F(X)$ and $F'_Y(Y) = F(Y)$ above.
In this way we see that the desired extension of
$F$ is the functor which sends $X/S$ to $F'_X(X)$.

\medskip\noindent
Property (a) for the functor $X \mapsto F'_X(X)$ is almost immediate
from the construction; we omit the details.
Suppose that $X = \lim_{i \in I} X_i$
is a directed limit of affine schemes over $S$. We have to show that
$$
F'_X(X) = \colim_{i \in I} F'_{X_i}(X_i)
$$
First assume that there is some $i \in I$ such that
$X_i \to S$ factors through an affine open $U \subset S$.
Then $F'$ is defined on $X$ and on $X_{i'}$ for $i' \geq i$
and we see that $F'_{X_{i'}}(X_{i'}) = F'(X_{i'})$ for
$i' \geq i$ and $F'_X(X) = F'(X)$. In this case every arrow
$X \to V$ with $V$ locally of finite presentation
over $S$ factors as $X \to X_{i'} \to V$ for some
$i' \geq i$, see Limits, Proposition
\ref{limits-proposition-characterize-locally-finite-presentation}.
Thus we have
\begin{align*}
F'_X(X)
& =
F'(X)  \\
& = \colim_{X \to V} F(V) \\
& =
\colim_{i' \geq i} \colim_{X_{i'} \to V} F(V) \\
& =
\colim_{i' \geq i} F'(X_{i'}) \\
& =
\colim_{i' \geq i} F'_{X_{i'}}(X_{i'}) \\
& =
\colim_{i' \in I} F'_{X_{i'}}(X_{i'})
\end{align*}
as desired. Finally, in general we pick any $i \in I$ and we choose
a finite affine open covering $V_i = V_{i, 1} \cup \ldots \cup V_{i, n}$
such that $V_{i, k} \to S$ factors through an affine open of $S$.
Let $V_k \subset V$ and $V_{i', k}$ for $i' \geq i$
be the inverse images of $V_{i, k}$.
By the previous case we see that
$$
F'_{V_k}(V_k) = \colim_{i' \geq i} F'_{V_{i', k}}(V_{i', k})
$$
and
$$
F'_{V_k \cap V_l}(V_k \cap V_l) =
\colim_{i' \geq i}
F'_{V_{i', k} \cap V_{i', l}}(V_{i', k} \cap V_{i', l})
$$
By the sheaf property and exactness of filtered colimits
we find that $F'_X(X) = \colim_{i \in I} F'_{X_i}(X_i)$
also in this case. This finishes the proof of property (b)
and hence finishes the proof of the lemma.
\end{proof}

\begin{lemma}[Limit fppf topology]
\label{topologies-lemma-limit-fppf-topology}
\source{stacks:049N}
\href{https://stacks.math.columbia.edu/tag/049N}{\texttt{Stacks tag 049N}}
\group{topologies}

\uses{topologies-definition-standard-Zariski, topologies-definition-standard-etale, topologies-definition-standard-smooth, topologies-definition-standard-syntomic, topologies-definition-standard-fppf}
Let $\tau \in \{Zariski, \etale, smooth, syntomic, fppf\}$.
Let $T$ be an affine scheme which is written as a limit
$T = \lim_{i \in I} T_i$ of a directed inverse system of affine schemes.
\begin{enumerate}
\item Let $\mathcal{V} = \{V_j \to T\}_{j = 1, \ldots, m}$ be a
standard $\tau$-covering of $T$, see Definitions
\ref{topologies-definition-standard-Zariski},
\ref{topologies-definition-standard-etale},
\ref{topologies-definition-standard-smooth},
\ref{topologies-definition-standard-syntomic}, and
\ref{topologies-definition-standard-fppf}.
Then there exists an index $i$ and a standard $\tau$-covering
$\mathcal{V}_i = \{V_{i, j} \to T_i\}_{j = 1, \ldots, m}$
whose base change $T \times_{T_i} \mathcal{V}_i$ to $T$
is isomorphic to $\mathcal{V}$.
\item Let $\mathcal{V}_i$, $\mathcal{V}'_i$ be a pair of standard
$\tau$-coverings of $T_i$. If
$f : T \times_{T_i} \mathcal{V}_i \to T \times_{T_i} \mathcal{V}'_i$ is
a morphism of coverings of $T$, then there exists an index
$i' \geq i$ and a morphism
$f_{i'} : T_{i'} \times_{T_i} \mathcal{V} \to
T_{i'} \times_{T_i} \mathcal{V}'_i$
whose base change to $T$ is $f$.
\item If
$f, g : \mathcal{V} \to \mathcal{V}'_i$
are morphisms of standard $\tau$-coverings of $T_i$ whose
base changes $f_T, g_T$ to $T$ are equal then there exists an
index $i' \geq i$ such that $f_{T_{i'}} = g_{T_{i'}}$.
\end{enumerate}
In other words, the category of standard $\tau$-coverings of $T$ is
the colimit over $I$ of the categories of standard $\tau$-coverings of $T_i$.
\end{lemma}

\begin{proof}
Let us prove this for $\tau = fppf$.
By Limits, Lemma \ref{limits-lemma-descend-finite-presentation}
the category of schemes of finite presentation over $T$ is the
colimit over $I$ of the categories of finite presentation over $T_i$. By
Limits, Lemmas \ref{limits-lemma-descend-affine-finite-presentation}
and \ref{limits-lemma-descend-flat-finite-presentation}
the same is true for category of schemes which are affine, flat and
of finite presentation over $T$.
To finish the proof of the lemma it suffices to show that if
$\{V_{j, i} \to T_i\}_{j = 1, \ldots, m}$ is a finite family of
flat finitely presented morphisms with $V_{j, i}$ affine, and the
base change $\coprod_j T \times_{T_i} V_{j, i} \to T$ is surjective,
then for some $i' \geq i$ the morphism
$\coprod T_{i'} \times_{T_i} V_{j, i} \to T_{i'}$ is surjective.
Denote $W_{i'} \subset T_{i'}$, resp.\ $W \subset T$ the image.
Of course $W = T$ by assumption.
Since the morphisms are flat and of finite presentation we see that
$W_i$ is a quasi-compact open of $T_i$, see
Morphisms, Lemma \ref{morphisms-lemma-fppf-open}.
Moreover, $W = T \times_{T_i} W_i$ (formation of image commutes
with base change). Hence by
Limits, Lemma \ref{limits-lemma-descend-opens}
we conclude that $W_{i'} = T_{i'}$ for some large enough $i'$
and we win.

\medskip\noindent
For $\tau \in \{Zariski, \etale, smooth, syntomic\}$ a standard $\tau$-covering
is a standard fppf covering. Hence the fully faithfulness of the functor
holds. The only issue is to show that given a standard fppf covering
$\mathcal{V}_i$ for some $i$ such that $\mathcal{V}_i \times_{T_i} T$
is a standard $\tau$-covering, then $\mathcal{V}_i \times_{T_i} T_{i'}$
is a standard $\tau$-covering for all $i' \gg i$. This follows immediately
from Limits, Lemmas
\ref{limits-lemma-descend-open-immersion},
\ref{limits-lemma-descend-etale},
\ref{limits-lemma-descend-smooth}, and
\ref{limits-lemma-descend-syntomic}.
\end{proof}

\begin{lemma}[Extend sheaf general]
\label{topologies-lemma-extend-sheaf-general}
\source{stacks:0GDW}
\href{https://stacks.math.columbia.edu/tag/0GDW}{\texttt{Stacks tag 0GDW}}
\group{topologies}

\uses{topologies-lemma-extend}
Let $S$, $\mathcal{C}$, $F$ satisfy conditions (1), (2), (a), and (b) of
Lemma \ref{topologies-lemma-extend} and denote $F' : (\Sch/S)^{opp} \to \textit{Sets}$
the unique extension constructed in the lemma. Let
$\tau \in \{Zariski, \etale, smooth, syntomic, fppf\}$. Assume
\begin{enumerate}
\item[(c)] for any standard $\tau$-covering $\{V_i \to V\}_{i = 1, \ldots, n}$
of affines in $\Sch/S$ such that $V \to S$ factors through an affine open
$U \subset S$ and $V \to U$ is of finite presentation, the sheaf condition
hold for $F$ and $\{V_i \to V\}_{i = 1, \ldots, n}$\footnote{This makes
sense as $V$, $V_i$, and $V_i \times_V V_j$ are isomorphic to objects
of $\mathcal{C}$ by (2).}.
\end{enumerate}
Then $F'$ satisfies the sheaf condition for all $\tau$-coverings.
\end{lemma}

\begin{proof}
\uses{topologies-lemma-limit-fppf-topology}

Let $X$ be a scheme over $S$ and let $\{X_i \to X\}_{i \in I}$
be a $\tau$-covering. Let $s_i \in F'(X_i)$ be elements such that
$s_i$ and $s_j$ map to the same element of $F'(X_i \times_X X_j)$
for all $i, j \in I$. We have to show that there is a
unique element $s \in F'(X)$ restricting to $s_i \in F'(X_i)$ for all $i \in I$.

\medskip\noindent
Special case: $X$ is an affine such that the structure morphism
maps into an affine open $U$ of $S$ and the covering
$\{X_i \to X\}_{i \in I}$ is a standard $\tau$-covering.
In this case we can write
$$
X = \lim V_k
$$
as a cofiltered limit with $V_k \to U$ of finite presentation
and $V_k$ affine. See Algebra, Lemma \ref{algebra-lemma-ring-colimit-fp}.
By Lemma \ref{topologies-lemma-limit-fppf-topology} there exists a $k$ and a standard
$\tau$-covering $\{V_{k, i} \to V_k\}_{i \in I}$ whose base change
to $X$ is the given covering. For $k' \geq k$ denote
$\{V_{k', i} \to V_{k'}\}_{i \in I}$ the base change to $V_{k'}$
of our covering. Then we see that
\begin{align*}
F'(X)
& =
\colim_{k' \geq k} F(V_k) \\
& =
\colim_{k' \geq k}
\text{Equalizer}(
\begin{array}{cc}\prod F(V_{k', i})
\longrightarrow\longrightarrow & \prod F(V_{k', i} \times_{V_{k'}} V_{k', j})\end{array}
\\
& =
\text{Equalizer}(
\begin{array}{cc}\colim_{k' \geq k}
\prod F(V_{k', i})
\longrightarrow\longrightarrow & \colim_{k' \geq k}
\prod F(V_{k', i} \times_{V_{k'}} V_{k', j})\end{array}
\\
& =
\text{Equalizer}(
\begin{array}{cc}\prod F'(X_i)
\longrightarrow\longrightarrow & \prod F'(X_i \times_X X_j)\end{array}
\end{align*}
The first equality holds by construction of $F'$. The second holds by
assumption (c). The third holds because filtered colimits are exact.
The fourth again holds by construction of $F'$.
In this way we find that the sheaf property holds for $F'$
with respect to $\{X_i \to X\}_{i \in I}$.

\medskip\noindent
General case. Choose an affine open covering $X = \bigcup U_k$ such that
each $U_k$ maps into an affine open of $S$. For every $k$ we may choose
a standard $\tau$-covering $\{V_{k, j} \to U_k\}_{j = 1, \ldots, m_k}$
which refines $\{X_i \times_X U_k \to U_k\}_{i \in I}$. For each
$j \in \{1, \ldots, m_k\}$ choose an index $i_{k, j} \in I$
and a morphism $g_{k, j} : V_{k, j} \to X_{i_{k, j}}$ over $X$.
Let $s_{k, j}$ be the element of $F'(V_{k, j})$ we get by restricting
$s_{i_{k, j}}$ via $g_{k, j}$. Observe that $s_{k, j}$ and $s_{k', j'}$
restrict to the same element of $F'(V_{k, j} \times_X V_{k', j'})$
for all $k$ and $k'$ and all $j \in \{1, \ldots, m_k\}$
and $j' \in \{1, \ldots, m_{k'}\}$; verification omitted.
In particular, by the result of the previous paragraph there is a unique
element $s_k \in F'(U_k)$ restricting to $s_{k, j}$ for all $j$.
With this notation we are ready to finish the proof.

\medskip\noindent
Proof of uniqueness of $s$: this is true because $F'$ satisfies
the sheaf property for Zariski coverings and $s|_{U_k}$ must be equal
to $s_k$ because both restrict to $s_{k, j}$ for all $j$.
This uniqueness then shows that $s_k$ and $s_{k'}$ must restrict
to the same section of $F'$ over (the non-affine scheme) $U_k \cap U_{k'}$
because these sections restrict to the same section over the
$\tau$-covering $\{V_{k, j} \times_X V_{k', j'} \to U_k \cap U_{k'}\}$.
Thus by the sheaf property for Zariski coverings, there is a unique
section $s$ of $F'$ over $X$ whose restriction to $U_k$ is $s_k$.
We omit the verification (similar to the above) that $s$ restricts
to $s_i$ over $X_i$.
\end{proof}

\begin{lemma}[Extend sheaf]
\label{topologies-lemma-extend-sheaf}
\source{stacks:0EUX}
\href{https://stacks.math.columbia.edu/tag/0EUX}{\texttt{Stacks tag 0EUX}}
\group{topologies}

\uses{topologies-lemma-extend}
Let $\tau \in \{Zariski, \etale, smooth, syntomic, fppf\}$.
Let $S$ be a scheme contained in a big site $\Sch_\tau$.
Let $F : (\Sch/S)_\tau^{opp} \to \textit{Sets}$ be a $\tau$-sheaf
satisfying property (b) of Lemma \ref{topologies-lemma-extend} with
$\mathcal{C} = (\Sch/S)_\tau$. Then the extension
$F'$ of $F$ to the category of all schemes over $S$
satisfies the sheaf condition for all $\tau$-coverings.
\end{lemma}

\begin{proof}
\uses{topologies-lemma-extend-sheaf-general, topologies-lemma-extend, topologies-lemma-zariski-induced, topologies-lemma-etale-induced, topologies-lemma-smooth-induced, topologies-lemma-syntomic-induced, topologies-lemma-fppf-induced}

This follows from Lemma \ref{topologies-lemma-extend-sheaf-general}
applied with $\mathcal{C} = (\Sch/S)_\tau$.
Conditions (1), (2), (a), and (b) of Lemma \ref{topologies-lemma-extend}
hold; we omit the details. Thus we get our unique extension $F'$
to the category of all schemes over $S$. Finally, observe that
any standard $\tau$-covering is tautologically equivalent to a
covering in $(\Sch/S)_\tau$, see Sets, Lemma \ref{sets-lemma-what-is-in-it}
as well as Lemmas \ref{topologies-lemma-zariski-induced},
\ref{topologies-lemma-etale-induced},
\ref{topologies-lemma-smooth-induced},
\ref{topologies-lemma-syntomic-induced}, and
\ref{topologies-lemma-fppf-induced}. By
Sites, Lemma \ref{sites-lemma-tautological-same-sheaf}
the sheaf property passes through tautological equivalence
of coverings. Hence the fact that $F$ is a $\tau$-sheaf implies
that property (c) of Lemma \ref{topologies-lemma-extend-sheaf-general} holds
and we conclude.
\end{proof}


\input{chapters}
